% Generated mechanically from frozen input inputs/p03/ja/tex/modules.tex.
% Do not edit this generated file; edit the bound locale source instead.

% OK, start here.
%

\setcounter{chapter}{16}
\chapter{加群の層}



\phantomsection
\label{modules-section-phantom}


\section{序論}
\label{modules-section-introduction}

\noindent
本章では，加群の層に関する基本概念を詳述する。
とりわけアーベル群の層もこれに含まれる。実際，そのような層は
$\underline{\mathbf{Z}}$-加群の層とみなすことができる。
基本的な参考文献は \cite{FAC}, \cite{EGA}, \cite{SGA4} である。

\medskip\noindent
環付きトポス上の加群の層については別の章で扱う（《サイト上の加群》第
\ref{sites-modules-section-introduction} 節を参照）。ただし，そこでは
本章の議論の大部分を繰り返すことになる。





\section{病的な場合}
\label{modules-section-pathology}

\noindent
環付き空間とは，位相空間 $X$ と環の層 $\mathcal{O}$ からなる組である。
この定義では $\mathcal{O} = 0$ も許す。この場合，加群の圏はただ一つの
対象（すなわち $0$）しかもたない。それでもアーベル圏ではあるが，
やや退化している。同様に，層 $\mathcal{O}$ が $X$ のある開部分集合上で
零になることもありうる。

\medskip\noindent
局所環付き空間を考える場合（後にそうする）には，この現象は起こらない。








\section{加群の層のアーベル圏}
\label{modules-section-kernels}

\noindent
$(X, \mathcal{O}_X)$ を環付き空間とする（《層》定義
\ref{sheaves-definition-ringed-space} を参照）。$\mathcal{F}$ と
$\mathcal{G}$ を $\mathcal{O}_X$-加群の層とする（《層》定義
\ref{sheaves-definition-sheaf-modules} を参照）。
$\varphi, \psi : \mathcal{F} \to \mathcal{G}$ を
$\mathcal{O}_X$-加群の層の射とする。各開集合 $U \subset X$ 上で
$\varphi$ と $\psi$ が誘導する写像の和となる射を
$\varphi + \psi : \mathcal{F} \to \mathcal{G}$ と定義する。
これは明らかに，やはり $\mathcal{O}_X$-加群の層の射である。
また，$\mathcal{O}_X$-加群の層の射の合成がこの加法に関して
双線形であることも明らかである。したがって
$\textit{Mod}(\mathcal{O}_X)$ は前加法圏である（《ホモロジー》定義
\ref{homology-definition-preadditive} を参照）。

\medskip\noindent
すべての開集合 $U \subset X$ に対して定値 $\{0\}$ をとる
$\mathcal{O}_X$-加群の層を $0$ と書く。これは明らかに
$\textit{Mod}(\mathcal{O}_X)$ の終対象であると同時に始対象でもある。
$\mathcal{O}_X$-加群の層の射
$\varphi : \mathcal{F} \to \mathcal{G}$ に対して，次の条件は同値である：
(a) $\varphi$ は零射である，(b) $\varphi$ は $0$ を経由して分解する，
(c) $\varphi$ は各開集合 $U$ 上の切断において零である，
(d) すべての $x \in X$ に対して $\varphi_x = 0$ である。
《層》補題 \ref{sheaves-lemma-points-exactness} を参照せよ。

\medskip\noindent
さらに，二つの $\mathcal{O}_X$-加群の層
$\mathcal{F}$, $\mathcal{G}$ が与えられたとき，直和を
$$
\mathcal{F} \oplus \mathcal{G} = \mathcal{F} \times \mathcal{G}
$$
と定義し，《ホモロジー》定義 \ref{homology-definition-direct-sum}
における明らかな射 $(i, j, p, q)$ を備えることができる。したがって
$\textit{Mod}(\mathcal{O}_X)$ は加法圏である（《ホモロジー》定義
\ref{homology-definition-additive-category} を参照）。

\medskip\noindent
$\varphi : \mathcal{F} \to \mathcal{G}$ を
$\mathcal{O}_X$-加群の層の射とする。$\Ker(\varphi)$ を，切断が
$$
\Ker(\varphi)(U) =
\{ s \in \mathcal{F}(U) \mid \varphi(s) = 0 \text{（} \mathcal{G}(U)\text{ において）}\}
$$
となる $\mathcal{F}$ の部分層として定義する。ここで $U$ は
$X$ のすべての開集合を走る。これが実際に
$\mathcal{O}_X$-加群の層の圏における核であることは容易に分かる。
言い換えると，射 $\alpha : \mathcal{H} \to \mathcal{F}$ が
$\Ker(\varphi)$ を経由して分解することと
$\varphi \circ \alpha = 0$ であることは同値である。さらに茎のレベルでは
$\Ker(\varphi)_x = \Ker(\varphi_x)$ が成り立つ。

\medskip\noindent
一方，$\Coker(\varphi)$ を，次の規則で定まる
$\mathcal{O}_X$-加群の前層に付随する $\mathcal{O}_X$-加群の層として定義する：
$$
U
\longmapsto
\Coker(\mathcal{F}(U)\to \mathcal{G}(U)) =
\mathcal{G}(U)/\varphi(\mathcal{F}(U)).
$$
茎をとることと層化は可換なので（《層》補題
\ref{sheaves-lemma-stalk-sheafification} を参照），
$\Coker(\varphi)_x = \Coker(\varphi_x)$ である。したがって写像
$\mathcal{G} \to \Coker(\varphi)$ は（集合の層の射として）全射である。
《層》第 \ref{sheaves-section-exactness-points} 節を参照せよ。
これが余核であることを示すため，$\beta \circ \varphi$ が零となる
$\mathcal{O}_X$-加群の層の射
$\beta : \mathcal{G} \to \mathcal{H}$ を考える。このとき各開集合
$U \subset X$ に対し，$\beta$ は
$\mathcal{G}(U)/\varphi(\mathcal{F}(U))$ から $\mathcal{H}(U)$ への
写像を誘導する。層化の普遍性（《層》補題
\ref{sheaves-lemma-sheafification-presheaf-modules} を参照）により，
もとの $\beta$ が合成
$\mathcal{G} \to \Coker(\varphi) \to \mathcal{H}$ に等しくなるような
標準的な射 $\Coker(\varphi) \to \mathcal{H}$ を得る。
この射は，上で述べた全射性により一意である。

\begin{lemma}
\label{modules-lemma-abelian}
$(X, \mathcal{O}_X)$ を環付き空間とする。圏
$\textit{Mod}(\mathcal{O}_X)$ はアーベル圏である。さらに，複体
$$
\mathcal{F} \to \mathcal{G} \to \mathcal{H}
$$
が $\mathcal{G}$ において完全であるための必要十分条件は，すべての
$x \in X$ に対して複体
$$
\mathcal{F}_x \to \mathcal{G}_x \to \mathcal{H}_x
$$
が $\mathcal{G}_x$ において完全となることである。
\end{lemma}

\begin{proof}
《ホモロジー》定義 \ref{homology-definition-abelian-category} により，
像と余像が一致することを示せばよい。《層》補題
\ref{sheaves-lemma-points-exactness} により，像と余像が各
$x \in X$ において同じ茎をもつことを示せば十分である。
上の核と余核の構成から，これらの茎は $\mathcal{O}_{X, x}$-加群の圏に
おける余像と像である。したがって，環上の加群の圏がアーベル圏であることから
結論が従う。
\end{proof}

\noindent
実際，圏 $\textit{Mod}(\mathcal{O}_X)$ はさらに多くの性質をもつ。
ここでは二つの構成を挙げる。
\begin{enumerate}
\item 任意の集合 $I$ と，各 $i \in I$ に対する
$\mathcal{O}_X$-加群の層 $\mathcal{F}_i$ が与えられたとき，積
$$
\prod\nolimits_{i \in I} \mathcal{F}_i
$$
をつくることができる。これは，各開集合 $U$ に加群
$\mathcal{F}_i(U)$ の積を対応させる層である。またこれは，
《圏論》定義 \ref{categories-definition-product} における圏論的な積でもある。
\item 任意の集合 $I$ と，各 $i \in I$ に対する
$\mathcal{O}_X$-加群の層 $\mathcal{F}_i$ が与えられたとき，直和
$$
\bigoplus\nolimits_{i \in I} \mathcal{F}_i
$$
をつくることができる。これは，各開集合 $U$ に加群
$\mathcal{F}_i(U)$ の直和を対応させる前層の{\it 層化}である。
またこれは，《圏論》定義 \ref{categories-definition-coproduct}
における圏論的な余積でもある。このことは層化の普遍性から分かる。
\end{enumerate}
以上から，$\textit{Mod}(\mathcal{O}_X)$ にはすべての極限と余極限が
存在することが分かる。

\begin{lemma}
\label{modules-lemma-limits-colimits}
$(X, \mathcal{O}_X)$ を環付き空間とする。
\begin{enumerate}
\item $\textit{Mod}(\mathcal{O}_X)$ にはすべての極限が存在する。
極限は $\mathcal{O}_X$-加群の前層の対応する極限と同じである
（すなわち，開集合上の切断をとることと可換である）。
\item $\textit{Mod}(\mathcal{O}_X)$ にはすべての余極限が存在する。
余極限は，前層の圏における対応する余極限の層化である。
余極限をとることと茎をとることは可換である。
\item フィルター余極限は完全である。
\item 有限直和は，$\mathcal{O}_X$-加群の前層の対応する有限直和と同じである。
\end{enumerate}
\end{lemma}

\begin{proof}
$\textit{Mod}(\mathcal{O}_X)$ はアーベル圏なので（補題
\ref{modules-lemma-abelian}），すべての有限極限と有限余極限をもつ
（《ホモロジー》補題 \ref{homology-lemma-colimit-abelian-category}）。
したがって，極限と余極限の存在およびその記述は，上で述べた積と余積の存在と
その記述，ならびに《圏論》補題
\ref{categories-lemma-limits-products-equalizers} および
\ref{categories-lemma-colimits-coproducts-coequalizers} から従う。
層化と茎をとることは可換なので，余極限をとることと茎をとることも可換である。
(3) の意味は次のとおりである。有向集合 $I$ 上で
$\mathcal{O}_X$-加群の層の完全列
$0 \to \mathcal{F}_i \to \mathcal{G}_i \to \mathcal{H}_i \to 0$
からなる系が与えられたとき，列
$0 \to \colim \mathcal{F}_i \to \colim \mathcal{G}_i \to
\colim \mathcal{H}_i \to 0$ も完全である。完全性は茎で判定できるので
（補題 \ref{modules-lemma-abelian}），これは加群の場合，すなわち《代数》補題
\ref{algebra-lemma-directed-colimit-exact} から従う。(4) の証明は省略する。
\end{proof}

\noindent
極限と余極限が存在するので，《圏論》第
\ref{categories-section-exact-functor} 節のように，
$\mathcal{O}$-加群の層の圏上で定義された関手の完全性を，
極限と余極限によって考えることができる。短完全列による完全性の記述については，
《ホモロジー》補題 \ref{homology-lemma-exact-functor} を参照せよ。

\begin{lemma}
\label{modules-lemma-exactness-pushforward-pullback}
$f : (X, \mathcal{O}_X) \to (Y, \mathcal{O}_Y)$ を
環付き空間の射とする。
\begin{enumerate}
\item 関手
$f_* : \textit{Mod}(\mathcal{O}_X) \to \textit{Mod}(\mathcal{O}_Y)$
は左完全である。実際，すべての極限と可換である。
\item 関手
$f^* : \textit{Mod}(\mathcal{O}_Y) \to \textit{Mod}(\mathcal{O}_X)$
は右完全である。実際，すべての余極限と可換である。
\item アーベル群の層に対する逆像
$f^{-1} : \textit{Ab}(Y) \to \textit{Ab}(X)$ は完全である。
\end{enumerate}
\end{lemma}

\begin{proof}
(1) と (2) は，$(f^*, f_*)$ が随伴関手の組だから成り立つ。
《層》補題 \ref{sheaves-lemma-adjoint-pullback-pushforward-modules}
および《圏論》第 \ref{categories-section-adjoint} 節を参照せよ。
(3) は，完全性を茎で判定できること（補題 \ref{modules-lemma-abelian}）と
逆像の茎の記述から従う。《層》補題
\ref{sheaves-lemma-pullback-abelian-stalk} を参照せよ。
\end{proof}

\begin{lemma}
\label{modules-lemma-j-shriek-exact}
$j : U \to X$ を位相空間の開埋め込みとする。関手
$j_! : \textit{Ab}(U) \to \textit{Ab}(X)$ は完全である。
\end{lemma}

\begin{proof}
《層》補題 \ref{sheaves-lemma-j-shriek-abelian} における茎の記述から従う。
\end{proof}

\begin{lemma}
\label{modules-lemma-section-direct-sum-quasi-compact}
$(X, \mathcal{O}_X)$ を環付き空間とし，$I$ を集合とする。
各 $i \in I$ に対し，$\mathcal{F}_i$ を $\mathcal{O}_X$-加群の層とする。
$U \subset X$ が準コンパクトな開集合ならば，写像
$$
\bigoplus\nolimits_{i \in I} \mathcal{F}_i(U)
\longrightarrow
\left(\bigoplus\nolimits_{i \in I} \mathcal{F}_i\right)(U)
$$
は全単射である。
\end{lemma}

\begin{proof}
$s$ を右辺の元とする。このとき開被覆
$U = \bigcup_{j \in J} U_j$ が存在し，各 $s|_{U_j}$ は
$s_{ji} \in \mathcal{F}_i(U_j)$ を用いた有限和
$\sum_{i \in I_j} s_{ji}$ となる。$U$ は準コンパクトなので，
この被覆が有限，すなわち $J$ が有限であると仮定してよい。
すると $I' = \bigcup_{j \in J} I_j$ は $I$ の有限部分集合である。
明らかに，$s$ は部分層 $\bigoplus_{i \in I'} \mathcal{F}_i$ の切断である。
有限直和については層化が不要であることから結論が従う。
上の補題 \ref{modules-lemma-limits-colimits} を参照せよ。
\end{proof}



\section{加群の層の切断}
\label{modules-section-sections}

\noindent
$(X, \mathcal{O}_X)$ を環付き空間とし，$\mathcal{F}$ を
$\mathcal{O}_X$-加群の層とする。また
$s \in \Gamma(X, \mathcal{F}) = \mathcal{F}(X)$ を大域切断とする。
$s$ に付随する一意な{\it $\mathcal{O}_X$-加群の層の射}
$$
\mathcal{O}_X \longrightarrow \mathcal{F}, \ f \longmapsto fs
$$
が存在する。上の記法が意味するのは，$\mathcal{O}_X$ の局所切断 $f$，
すなわちある開集合 $U$ 上の切断 $f$ が，$f$ と $s$ の $U$ への制限との
積に写されるということである。逆に，任意の射
$\varphi : \mathcal{O}_X \to \mathcal{F}$ は切断 $s = \varphi(1)$ を定め，
$\varphi$ は $s$ に付随する射となる。

\begin{definition}
\label{modules-definition-globally-generated}
$(X, \mathcal{O}_X)$ を環付き空間とし，$\mathcal{F}$ を
$\mathcal{O}_X$-加群の層とする。集合 $I$ と大域切断
$s_i \in \Gamma(X, \mathcal{F})$（$i \in I$）が存在し，写像
$$
\bigoplus\nolimits_{i \in I}
\mathcal{O}_X \longrightarrow \mathcal{F}
$$
が全射となるとき，$\mathcal{F}$ は{\it 大域切断で生成される}という。
ここでこの写像は，$i$ に対応する直和因子上では $s_i$ に付随する射である。
この場合，切断 $s_i$ は $\mathcal{F}$ を{\it 生成する}という。
\end{definition}

\noindent
《層》第 \ref{sheaves-section-stalks} 節で導入した記法の濫用を
しばしば用いる。すなわち，点 $x \in X$ のある開近傍上で定義された
$\mathcal{F}$ の局所切断 $s$ に対し，茎 $\mathcal{F}_x$ における
$s$ の像を $s_x$，あるいは単に $s$ と書く。

\begin{lemma}
\label{modules-lemma-globally-generated}
$(X, \mathcal{O}_X)$ を環付き空間とし，$\mathcal{F}$ を
$\mathcal{O}_X$-加群の層，$I$ を集合とする。また
$s_i \in \Gamma(X, \mathcal{F})$（$i \in I$）を大域切断とする。
切断 $s_i$ が $\mathcal{F}$ を生成するための必要十分条件は，
すべての $x\in X$ に対して元 $s_{i, x} \in \mathcal{F}_x$ が
$\mathcal{O}_{X, x}$-加群 $\mathcal{F}_x$ を生成することである。
\end{lemma}

\begin{proof}
省略する。
\end{proof}

\begin{lemma}
\label{modules-lemma-tensor-product-globally-generated}
\begin{slogan}
大域生成される加群の層のテンソル積は大域生成される。
\end{slogan}
$(X, \mathcal{O}_X)$ を環付き空間とし，$\mathcal{F}$ と
$\mathcal{G}$ を $\mathcal{O}_X$-加群の層とする。
$\mathcal{F}$ と $\mathcal{G}$ が大域切断で生成されるならば，
$\mathcal{F} \otimes_{\mathcal{O}_X} \mathcal{G}$ も同様である。
\end{lemma}

\begin{proof}
省略する。
\end{proof}

\begin{lemma}
\label{modules-lemma-generated-by-local-sections}
$(X, \mathcal{O}_X)$ を環付き空間とし，$\mathcal{F}$ を
$\mathcal{O}_X$-加群の層，$I$ を集合とする。$s_i$（$i \in I$）を
$\mathcal{F}$ の局所切断の族とする。すなわち，ある開集合
$U_i \subset X$ に対して $s_i \in \mathcal{F}(U_i)$ とする。
各 $s_i$ が $\mathcal{G}$ の局所切断に対応するような，最小の
$\mathcal{O}_X$-加群の部分層 $\mathcal{G}$ が一意に存在する。
\end{lemma}

\begin{proof}
次の規則で定まる $\mathcal{O}_X$-加群の部分前層を考える：
$$
U
\longmapsto
\{
\text{有限和 } \sum\nolimits_{i \in J} f_i (s_i|_U)
\text{ であって，}J\text{ は有限，}
U \subset U_i \text{ が } i\in J \text{ に対して成り立ち，かつ }
f_i \in \mathcal{O}_X(U)
\}
$$
$\mathcal{G}$ をこの部分前層の層化とする。《層》補題
\ref{sheaves-lemma-sheafify-epi-mono} により，これは
$\mathcal{F}$ の部分層である。上のすべての有限和は明らかに
$\mathcal{G}$ に含まれなければならないので，これは求める最小の部分層である。
\end{proof}

\begin{definition}
\label{modules-definition-generated-by-local-sections}
$(X, \mathcal{O}_X)$ を環付き空間とし，$\mathcal{F}$ を
$\mathcal{O}_X$-加群の層とする。集合 $I$ と $\mathcal{F}$ の局所切断
$s_i$（$i \in I$）が与えられたとき，上の補題
\ref{modules-lemma-generated-by-local-sections} の部分層 $\mathcal{G}$ を，
{\it $s_i$ によって生成される部分層}という。
\end{definition}

\begin{lemma}
\label{modules-lemma-generated-by-local-sections-stalk}
$(X, \mathcal{O}_X)$ を環付き空間とし，$\mathcal{F}$ を
$\mathcal{O}_X$-加群の層とする。集合 $I$ と $\mathcal{F}$ の局所切断
$s_i$（$i \in I$）が与えられているとする。$\mathcal{G}$ を
$s_i$ によって生成される部分層とし，$x\in X$ とする。このとき
$\mathcal{G}_x$ は，$s_i$ が $x$ で定義されているような $i$ に対する元
$s_{i, x}$ によって生成される，$\mathcal{F}_x$ の
$\mathcal{O}_{X, x}$-部分加群である。
\end{lemma}

\begin{proof}
これは補題 \ref{modules-lemma-generated-by-local-sections} の証明における
$\mathcal{G}$ の構成から明らかである。
\end{proof}










\section{加群の層と切断の台}
\label{modules-section-support}

\begin{definition}
\label{modules-definition-support}
$(X, \mathcal{O}_X)$ を環付き空間とし，$\mathcal{F}$ を
$\mathcal{O}_X$-加群の層とする。
\begin{enumerate}
\item {\it $\mathcal{F}$ の台}とは，$\mathcal{F}_x \not = 0$ となる
点 $x \in X$ の集合である。
\item $\mathcal{F}$ の台を $\text{Supp}(\mathcal{F})$ と書く。
\item $s \in \Gamma(X, \mathcal{F})$ を大域切断とする。
{\it $s$ の台}とは，$s$ の像 $s_x \in \mathcal{F}_x$ が零でないような
点 $x \in X$ の集合である。
\end{enumerate}
\end{definition}

\noindent
局所切断は $\mathcal{F}$ の制限の大域切断なので，これにより
局所切断の台も当然定義される。

\begin{lemma}
\label{modules-lemma-support-section-closed}
$(X, \mathcal{O}_X)$ を環付き空間とし，$\mathcal{F}$ を
$\mathcal{O}_X$-加群の層とする。$U \subset X$ を開集合とする。
\begin{enumerate}
\item $s \in \mathcal{F}(U)$ の台は $U$ の閉集合である。
\item $fs$ の台は，$f \in \mathcal{O}_X(U)$ の台と
$s \in \mathcal{F}(U)$ の台との共通部分に含まれる。
\item $s + s'$ の台は，$s, s' \in \mathcal{F}(U)$ の台の和集合に含まれる。
\item $\mathcal{F}$ の台は，$\mathcal{F}$ のすべての局所切断の台の和集合である。
\item $\varphi : \mathcal{F} \to \mathcal{G}$ が
$\mathcal{O}_X$-加群の層の射ならば，$\varphi(s)$ の台は
$s \in \mathcal{F}(U)$ の台に含まれる。
\end{enumerate}
\end{lemma}

\begin{proof}
実際，$s_x = 0$ ならば，茎の定義により $s$ は $x$ のある開近傍上で
零である。$f$ についても同様である。詳細は省略する。
\end{proof}

\noindent
一般に，加群の層の台は閉集合とは限らない。例えば，
$\mathbf{R}$（通常のアルキメデス的位相をもつ）上のアーベル群の層で，
それぞれ $\mathbf{R}$ の一点 $p_i$ に台をもつ零でない摩天楼層を
無限個直和したものを考えればよい。この層の台は点 $p_i$ の集合となり，
閉集合でないことがある。

\medskip\noindent
別の例として，開埋め込み
$j : U = (0 , \infty) \to \mathbf{R} = X$ とアーベル群の層
$j_!\underline{\mathbf{Z}}_U$ を考える。《層》第
\ref{sheaves-section-open-immersions} 節により，この層の台はちょうど $U$ である。

\begin{lemma}
\label{modules-lemma-support-sheaf-rings-closed}
$X$ を位相空間とする。環の層の台は閉集合である。
\end{lemma}

\begin{proof}
実際，本書の規約では，環が $0$ であることと $1 = 0$ であることは
同値である。したがって，環の層の台は単位切断の台である。
\end{proof}






\section{閉埋め込みとアーベル群の層}
\label{modules-section-closed-immersions}

\noindent
位相空間 $X$ 上のアーベル群の層を
$\underline{\mathbf{Z}}_X$-加群の層とみなしていることを思い出そう。
したがって，加群の層に関する任意の結果や定義をアーベル群の層に適用できる。

\begin{lemma}
\label{modules-lemma-i-star-exact}
$X$ を位相空間，$Z \subset X$ を閉部分集合とし，包含写像を
$i : Z \to X$ と書く。関手
$$
i_* : \textit{Ab}(Z) \longrightarrow \textit{Ab}(X)
$$
は完全かつ充満忠実であり，その本質像は台が $Z$ に含まれる
アーベル群の層全体にちょうど一致する。関手 $i^{-1}$ は $i_*$ の左逆である。
\end{lemma}

\begin{proof}
完全性は，《層》補題 \ref{sheaves-lemma-stalks-closed-pushforward}
における茎の記述と補題 \ref{modules-lemma-abelian} から従う。残りは
《層》補題 \ref{sheaves-lemma-equivalence-categories-closed-abelian}
で示した。
\end{proof}

\noindent
$\mathcal{F}$ を位相空間 $X$ 上のアーベル群の層とする。閉部分集合 $Z$ が
与えられたとき，台が $Z$ に含まれる切断だけからなる
$\mathcal{F}$ の標準的なアーベル部分層が存在する。正確には次のとおりである。

\begin{remark}
\label{modules-remark-sections-support-in-closed}
$X$ を位相空間，$Z \subset X$ を閉部分集合とし，$\mathcal{F}$ を
$X$ 上のアーベル群の層とする。開集合 $U \subset X$ に対し，
$$
\mathcal{H}_Z(\mathcal{F})(U) =
\{s \in \mathcal{F}(U) \mid
\text{ }s\text{ の台が }Z \cap U\text{ に含まれる}\}\text{。}
$$
と定義する。このとき $\mathcal{H}_Z(\mathcal{F})$ は $\mathcal{F}$ の
アーベル部分層であり，台が $Z$ に含まれるもののうち最大である。
補題 \ref{modules-lemma-i-star-exact} により，$\mathcal{H}_Z(\mathcal{F})$ を
$Z$ 上のアーベル群の層とみなすことができ，実際そうみなす。
このようにして左完全関手
$$
\textit{Ab}(X) \longrightarrow \textit{Ab}(Z),\quad
\mathcal{F} \longmapsto \mathcal{H}_Z(\mathcal{F})
\text{（}Z\text{ 上のアーベル群の層とみなす）}
$$
を得る。上の主張はすべて，補題 \ref{modules-lemma-support-section-closed}
から直接従う。
\end{remark}

\noindent
ここで，関手 $i_*$ がアーベル群の層上で右随伴をもつことを示しておこう。

\begin{lemma}
\label{modules-lemma-i-star-right-adjoint}
$i : Z \to X$ を閉部分集合 $Z$ の位相空間 $X$ への包含とする。
注意 \ref{modules-remark-sections-support-in-closed} の関手
$\textit{Ab}(X) \to \textit{Ab}(Z)$，
$\mathcal{F} \mapsto \mathcal{H}_Z(\mathcal{F})$ は
$i_* : \textit{Ab}(Z) \to \textit{Ab}(X)$ の右随伴である。
特に $i_*$ は任意の余極限と可換である。
\end{lemma}

\begin{proof}
$X$ 上の任意のアーベル群の層 $\mathcal{F}$ と $Z$ 上の任意の
アーベル群の層 $\mathcal{G}$ に対して
$$
\Hom_{\textit{Ab}(X)}(i_*\mathcal{G}, \mathcal{F}) =
\Hom_{\textit{Ab}(Z)}(\mathcal{G}, \mathcal{H}_Z(\mathcal{F}))
$$
を示せばよい。これは，$i_*\mathcal{G}$ の任意の切断が $Z$ に台をもつので
明らかである。詳細は省略する。
\end{proof}

\begin{remark}
\label{modules-remark-i-star-right-adjoint}
《層》注意 \ref{sheaves-remark-i-star-not-exact} では，集合の層の圏上の
関手としての $i_*$ は完全でないため，右随伴をもたないことを示した。
しかし，これに非常に近い主張は成り立つ。実際，関手 $i_*$ は
基点付き集合の層上では完全であり，そのような層について $Z$ に台をもつ
切断を定義できる。したがって $\mathcal{H}_Z$ は意味をもち，$i_*$ の
右随伴となる。
\end{remark}









\section{標準完全列}
\label{modules-section-canonical-exact-sequence}

\noindent
この完全列には独立した一節を設ける。

\begin{lemma}
\label{modules-lemma-canonical-exact-sequence}
$X$ を位相空間とし，$U \subset X$ を開部分集合，
$Z \subset X$ をその補集合とする。開埋め込みを $j : U \to X$，
閉埋め込みを $i : Z \to X$ と書く。$X$ 上の任意のアーベル群の層
$\mathcal{F}$ に対して，随伴写像
$j_{!}j^{-1}\mathcal{F} \to \mathcal{F}$ および
$\mathcal{F} \to i_*i^{-1}\mathcal{F}$ は短完全列
$$
0 \to j_{!}j^{-1}\mathcal{F} \to \mathcal{F} \to i_*i^{-1}\mathcal{F} \to 0
$$
を与える。$X$ 上のアーベル群の層の任意の射
$\varphi : \mathcal{F} \to \mathcal{G}$ に対して，短完全列の射
$$
\xymatrix{
0 \ar[r] &
j_{!}j^{-1}\mathcal{F} \ar[r] \ar[d] &
\mathcal{F} \ar[r] \ar[d] &
i_*i^{-1}\mathcal{F} \ar[r] \ar[d] &
0 \\
0 \ar[r] &
j_{!}j^{-1}\mathcal{G} \ar[r] &
\mathcal{G} \ar[r] &
i_*i^{-1}\mathcal{G} \ar[r] &
0
}
$$
を得る。
\end{lemma}

\begin{proof}
短完全列の関手性は，随伴写像の自然性から直ちに従う。
完全性は茎で判定できる（補題 \ref{modules-lemma-abelian}）。
ここに現れる茎の記述については，《層》補題
\ref{sheaves-lemma-j-shriek-abelian} および
\ref{sheaves-lemma-stalks-closed-pushforward} を参照せよ。
\end{proof}








\section{局所的に切断で生成される加群の層}
\label{modules-section-locally-generated}

\noindent
$(X, \mathcal{O}_X)$ を環付き空間とする。本節と次節では，層を
開部分空間 $U \subset X$ にしばしば制限する。《層》第
\ref{sheaves-section-open-immersions} 節を参照せよ。特に，より正確な
記法 $(U, \mathcal{O}_X|_U)$ の代わりに，この開部分空間を
$(U, \mathcal{O}_U)$ と書くことが多い。《層》定義
\ref{sheaves-definition-restriction} を参照せよ。

\medskip\noindent
開埋め込み $j : U = (0 , \infty) \to \mathbf{R} = X$ と
アーベル群の層 $j_!\underline{\mathbf{Z}}_U$ を考える。《層》第
\ref{sheaves-section-open-immersions} 節により，
$j_!\underline{\mathbf{Z}}_U$ の $x = 0$ における茎は $0$ である。
実際，$0$ を含む任意の開区間上で，この層の切断は $0$ である。
したがって，この層が切断で生成されるような点 $0$ の開近傍は存在しない。

\begin{definition}
\label{modules-definition-locally-generated}
$(X, \mathcal{O}_X)$ を環付き空間とし，$\mathcal{F}$ を
$\mathcal{O}_X$-加群の層とする。すべての $x \in X$ に対して
$x$ の開近傍 $U$ が存在し，$\mathcal{F}|_U$ が
$\mathcal{O}_U$-加群の層として大域生成されるとき，
$\mathcal{F}$ は{\it 局所的に切断で生成される}という。
\end{definition}

\noindent
言い換えると，集合 $I$ と，各 $i$ に対する切断
$s_i \in \mathcal{F}(U)$ が存在し，これらに付随する写像
$$
\bigoplus\nolimits_{i \in I} \mathcal{O}_U
\longrightarrow
\mathcal{F}|_U
$$
が全射となる。

\begin{lemma}
\label{modules-lemma-pullback-locally-generated}
$f : (X, \mathcal{O}_X) \to (Y, \mathcal{O}_Y)$ を環付き空間の射とする。
$\mathcal{G}$ が局所的に切断で生成されるならば，引き戻し
$f^*\mathcal{G}$ も局所的に切断で生成される。
\end{lemma}

\begin{proof}
$Y$ の開部分空間 $V$ が与えられたとき，環付き空間の可換図式
$$
\xymatrix{
(f^{-1}V, \mathcal{O}_{f^{-1}V}) \ar[r]_{j'} \ar[d]_{f'} &
(X, \mathcal{O}_X) \ar[d]^f \\
(V, \mathcal{O}_V) \ar[r]^j &
(Y, \mathcal{O}_Y)
}
$$
を考えることができる。《層》補題
\ref{sheaves-lemma-push-pull-composition-modules} により，
$f^*\mathcal{G}|_{f^{-1}V} \cong (f')^*(\mathcal{G}|_V)$ である。
したがって，$\mathcal{G}$ が大域生成されると仮定してよい。

\medskip\noindent
$f^*$ はすべての余極限と可換であり，右完全であることをすでに見た
（補題 \ref{modules-lemma-exactness-pushforward-pullback} を参照）。
したがって，全射
$$
\bigoplus\nolimits_{i \in I}
\mathcal{O}_Y
\to
\mathcal{G}
\to
0
$$
があれば，$f^*$ を適用して全射
$$
\bigoplus\nolimits_{i \in I}
\mathcal{O}_X
\to
f^*\mathcal{G}
\to
0.
$$
を得る。これで補題が従う。
\end{proof}












\section{有限型加群の層}
\label{modules-section-finite-type}

\begin{definition}
\label{modules-definition-finite-type}
$(X, \mathcal{O}_X)$ を環付き空間とし，$\mathcal{F}$ を
$\mathcal{O}_X$-加群の層とする。すべての $x \in X$ に対して
開近傍 $U$ が存在し，$\mathcal{F}|_U$ が有限個の切断で生成されるとき，
$\mathcal{F}$ は{\it 有限型}であるという。
\end{definition}

\begin{lemma}
\label{modules-lemma-pullback-finite-type}
$f : (X, \mathcal{O}_X) \to (Y, \mathcal{O}_Y)$ を環付き空間の射とする。
有限型 $\mathcal{O}_Y$-加群の層 $\mathcal{G}$ の引き戻し
$f^*\mathcal{G}$ は有限型 $\mathcal{O}_X$-加群の層である。
\end{lemma}

\begin{proof}
補題 \ref{modules-lemma-pullback-locally-generated} の証明と同様に議論すれば，
$\mathcal{G}$ は有限個の切断で大域生成されると仮定してよい。
$f^*$ はすべての余極限と可換であり，右完全であることをすでに見た
（補題 \ref{modules-lemma-exactness-pushforward-pullback} を参照）。
したがって，全射
$$
\bigoplus\nolimits_{i = 1, \ldots, n}
\mathcal{O}_Y
\to
\mathcal{G}
\to
0
$$
があれば，$f^*$ を適用して全射
$$
\bigoplus\nolimits_{i = 1, \ldots, n}
\mathcal{O}_X
\to
f^*\mathcal{G}
\to
0.
$$
を得る。これで補題が従う。
\end{proof}

\begin{lemma}
\label{modules-lemma-extension-finite-type}
$X$ を環付き空間とする。有限型 $\mathcal{O}_X$-加群の層の射の像は
有限型である。また，
$0 \to \mathcal{F}_1 \to \mathcal{F}_2 \to \mathcal{F}_3 \to 0$
を $\mathcal{O}_X$-加群の層の短完全列とする。
$\mathcal{F}_1$ と $\mathcal{F}_3$ が有限型ならば，
$\mathcal{F}_2$ も有限型である。
\end{lemma}

\begin{proof}
像についての主張は自明である。短完全列についての主張は，
$\mathcal{F}_3$ の切断が局所的には $\mathcal{F}_2$ の切断へ持ち上がることと，
環上の加群の圏における対応する結果（例えば茎に適用する）から従う。
\end{proof}

\begin{lemma}
\label{modules-lemma-finite-type-surjective-on-stalk}
$X$ を環付き空間とし，
$\varphi : \mathcal{G} \to \mathcal{F}$ を $\mathcal{O}_X$-加群の層の
準同型とする。$x \in X$ とする。$\mathcal{F}$ が有限型であり，茎上の写像
$\varphi_x : \mathcal{G}_x \to \mathcal{F}_x$ が全射であると仮定する。
このとき，$\varphi|_U$ が全射となる開近傍 $x \in U \subset X$ が存在する。
\end{lemma}

\begin{proof}
$U$ 上で $\mathcal{F}$ が $s_1, \ldots, s_n \in \mathcal{F}(U)$ により
生成されるような，$x$ の開近傍 $U \subset X$ を選ぶ。
$\varphi_x$ が全射であるという仮定により，$U$ を縮小すれば，
ある $t_i \in \mathcal{G}(U)$ に対して $s_i = \varphi(t_i)$ と仮定してよい。
この $U$ が条件を満たす。
\end{proof}

\begin{lemma}
\label{modules-lemma-finite-type-stalk-zero}
$X$ を環付き空間とし，$\mathcal{F}$ を $\mathcal{O}_X$-加群の層，
$x \in X$ とする。$\mathcal{F}$ が有限型で $\mathcal{F}_x = 0$ と仮定する。
このとき，$\mathcal{F}|_U$ が零となる開近傍 $x \in U \subset X$ が存在する。
\end{lemma}

\begin{proof}
これは補題 \ref{modules-lemma-finite-type-surjective-on-stalk} を射
$0 \to \mathcal{F}$ に適用した特別な場合である。
\end{proof}

\begin{lemma}
\label{modules-lemma-support-finite-type-closed}
\begin{slogan}
任意の環付き空間上で，有限型加群の層の台は閉集合である。
\end{slogan}
$(X, \mathcal{O}_X)$ を環付き空間とし，$\mathcal{F}$ を
$\mathcal{O}_X$-加群の層とする。$\mathcal{F}$ が有限型ならば，
$\mathcal{F}$ の台は閉集合である。
\end{lemma}

\begin{proof}
これは補題 \ref{modules-lemma-finite-type-stalk-zero} の言い換えである。
\end{proof}

\begin{lemma}
\label{modules-lemma-finite-type-quasi-compact-colimit}
$X$ を環付き空間，$I$ を前順序集合とし，
$(\mathcal{F}_i, f_{ii'})$ を $\mathcal{O}_X$-加群の層からなる $I$ 上の系とする
（《圏論》第 \ref{categories-section-posets-limits} 節を参照）。
$\mathcal{F} = \colim \mathcal{F}_i$ をその余極限とする。次を仮定する：
(a) $I$ は有向である，(b) $\mathcal{F}$ は有限型 $\mathcal{O}_X$-加群の層である，
(c) $X$ は準コンパクトである。このとき，
$\mathcal{F}_i \to \mathcal{F}$ が全射となる $i$ が存在する。
遷移写像 $f_{ii'}$ が単射ならば，ある $i \in I$ に対して
$\mathcal{F} = \mathcal{F}_i$ である。
\end{lemma}

\begin{proof}
$x \in X$ とする。$x$ の開近傍 $U \subset X$ と有限個の切断
$s_j \in \mathcal{F}(U)$（$j = 1, \ldots, m$）が存在し，
$s_1, \ldots, s_m$ は $\mathcal{O}_U$-加群の層として
$\mathcal{F}$ を生成する。必要なら $U$ を $x$ のより小さい開近傍へ縮小し，
各 $s_j$ がある $i \in I$ に対する $\mathcal{F}_i$ の切断から来ると仮定してよい。
$X$ は準コンパクトなので，有限開被覆
$X = \bigcup_{j = 1, \ldots, m} U_j$ をとることができ，各 $j$ に対して
添字 $i_j$ と有限個の切断 $s_{jl} \in \mathcal{F}_{i_j}(U_j)$ を選んで，
それらの像が $\mathcal{F}$ の $U_j$ への制限を生成するようにできる。
明らかに，すべての $i_j$ 以上である任意の添字 $i \in I$ に対して
補題の主張が成り立つ。
\end{proof}

\begin{lemma}
\label{modules-lemma-set-isomorphism-classes-finite-type-modules}
$X$ を環付き空間とする。有限型 $\mathcal{O}_X$-加群の層からなる集合
$\{\mathcal{F}_i\}_{i \in I}$ であって，$X$ 上の各有限型
$\mathcal{O}_X$-加群の層が $\mathcal{F}_i$ のちょうど一つと同型になるものが存在する。
\end{lemma}

\begin{proof}
各開被覆 $\mathcal{U} : X = \bigcup U_j$ に対し，
$\mathcal{O}_X$-加群の層 $\mathcal{F}$ であって，各制限
$\mathcal{F}|_{U_j}$ がある $r_j \geq 0$ に対する
$\mathcal{O}_{U_j}^{\oplus r_j}$ の商となるものを考える。
これらは，部分層 $\mathcal{K}_j \subset \mathcal{O}_{U_j}^{\oplus r_j}$ と
貼り合わせデータ
$$
\varphi_{jj'} :
\mathcal{O}_{U_j \cap U_{j'}}^{\oplus r_j}/
(\mathcal{K}_j|_{U_j \cap U_{j'}})
\longrightarrow
\mathcal{O}_{U_j \cap U_{j'}}^{\oplus r_{j'}}/
(\mathcal{K}_{j'}|_{U_j \cap U_{j'}})
$$
によってパラメータ付けられる。《層》第
\ref{sheaves-section-glueing-sheaves} 節を参照せよ。
すべての貼り合わせデータの集まりは集合をなすことに注意する。
また，写像 $J \to \mathcal{P}(X)$，$j \mapsto U_j$ が単射であるような
すべての被覆 $\mathcal{U} : X = \bigcup_{j \in J} U_i$ の集まりも集合をなす。
したがって，上のような商を貼り合わせて得られるすべての
$\mathcal{O}_X$-加群の層の集まりは集合 $\mathcal{I}$ をなす。
定義により，すべての有限型 $\mathcal{O}_X$-加群の層は
$\mathcal{I}$ のある元と同型である。$\mathcal{I}$ の各同型類から一つずつ
元を選べば，求める層の集合を得る（選択公理を用いる）。
\end{proof}













\section{準連接加群の層}
\label{modules-section-quasi-coherent}

\noindent
本節では，準連接 $\mathcal{O}_X$-加群の層という抽象的概念を導入する。
スキーム上の準連接加群の層は任意のアフィン開集合上でよく記述できるため，
この概念は代数幾何学において非常に有用である。しかし，（局所）環付き空間という
一般の状況では，この概念がまったく良い振る舞いをしないことに注意されたい。
一般に準連接層の圏はアーベル圏ではなく，準連接層の無限直和も
準連接とは限らない，などである。

\begin{definition}
\label{modules-definition-quasi-coherent}
$(X, \mathcal{O}_X)$ を環付き空間とし，$\mathcal{F}$ を
$\mathcal{O}_X$-加群の層とする。すべての点 $x \in X$ に対して
開近傍 $x\in U \subset X$ が存在し，$\mathcal{F}|_U$ が写像
$$
\bigoplus\nolimits_{j \in J}
\mathcal{O}_U
\longrightarrow
\bigoplus\nolimits_{i \in I}
\mathcal{O}_U
$$
の余核と同型であるとき，$\mathcal{F}$ を{\it 準連接
$\mathcal{O}_X$-加群の層}という。準連接 $\mathcal{O}_X$-加群の層の圏を
$\QCoh(\mathcal{O}_X)$ と書く。
\end{definition}

\noindent
この定義は，$X$ が，その上で $\mathcal{F}|_U$ が次の形の
{\it 表示}をもつような開集合 $U$ で被覆されることを意味する：
$$
\bigoplus\nolimits_{j \in J}
\mathcal{O}_U
\longrightarrow
\bigoplus\nolimits_{i \in I}
\mathcal{O}_U
\longrightarrow
\mathcal{F}|_U
\longrightarrow
0.
$$
ここで「表示」とは，表示された列が完全であることを意味する。言い換えると，
\begin{enumerate}
\item $X$ の各点 $x$ に対して開近傍 $U$ が存在し，
$\mathcal{F}|_U$ は大域切断で生成され，かつ
\item これらの切断を適切に選べば，対応する全射の核も大域切断で生成される。
\end{enumerate}

\begin{lemma}
\label{modules-lemma-direct-sum-quasi-coherent}
$(X, \mathcal{O}_X)$ を環付き空間とする。二つの準連接
$\mathcal{O}_X$-加群の層の直和は，準連接 $\mathcal{O}_X$-加群の層である。
\end{lemma}

\begin{proof}
省略する。
\end{proof}

\begin{remark}
\label{modules-remark-infinite-direct-sum-quasi-coherent-not}
注意：一般に，準連接 $\mathcal{O}_X$-加群の層の無限直和が
準連接であるとは限らない。準連接加群の層のさらに奇妙な振る舞いについては，
例 \ref{modules-example-quasi-coherent} を参照せよ。
\end{remark}

\begin{lemma}
\label{modules-lemma-pullback-quasi-coherent}
$f : (X, \mathcal{O}_X) \to (Y, \mathcal{O}_Y)$ を環付き空間の射とする。
準連接 $\mathcal{O}_Y$-加群の層 $\mathcal{G}$ の引き戻し
$f^*\mathcal{G}$ は準連接である。
\end{lemma}

\begin{proof}
補題 \ref{modules-lemma-pullback-locally-generated} の証明と同様に議論すれば，
$\mathcal{G}$ は $\mathcal{O}_Y$ のコピーの直和による大域的表示をもつと
仮定してよい。$f^*$ はすべての余極限と可換であり，右完全であることを
すでに見た（補題 \ref{modules-lemma-exactness-pushforward-pullback} を参照）。
したがって，完全列
$$
\bigoplus\nolimits_{j \in J}
\mathcal{O}_Y
\longrightarrow
\bigoplus\nolimits_{i \in I}
\mathcal{O}_Y
\longrightarrow
\mathcal{G}
\longrightarrow
0
$$
があれば，$f^*$ を適用して完全列
$$
\bigoplus\nolimits_{j \in J}
\mathcal{O}_X
\longrightarrow
\bigoplus\nolimits_{i \in I}
\mathcal{O}_X
\longrightarrow
f^*\mathcal{G}
\longrightarrow
0.
$$
を得る。これで補題が従う。
\end{proof}

\noindent
これにより，準連接層の例が数多く得られる。

\begin{lemma}
\label{modules-lemma-construct-quasi-coherent-sheaves}
$(X, \mathcal{O}_X)$ を環付き空間とする。
$\alpha : R \to \Gamma(X, \mathcal{O}_X)$ を，環 $R$ から
$X$ 上の大域切断の環への環準同型とする。$M$ を $R$-加群とする。
次の三つの構成は，標準的に同型な $\mathcal{O}_X$-加群の層を与える：
\begin{enumerate}
\item $\pi : X \to \{*\}$ が一意な写像であり，$\pi$-写像
$\pi^\sharp$ が与えられた写像
$\alpha : R \to \Gamma(X, \mathcal{O}_X)$ であるような環付き空間の射
$\pi : (X, \mathcal{O}_X) \longrightarrow (\{*\}, R)$ をとる。
$\mathcal{F}_1 = \pi^*M$ とおく。
\item 表示
$\bigoplus_{j \in J} R \to \bigoplus_{i \in I} R \to M \to 0$.
を選び，
$$
\mathcal{F}_2 = \Coker\left(
\bigoplus\nolimits_{j \in J} \mathcal{O}_X
\to
\bigoplus\nolimits_{i \in I} \mathcal{O}_X
\right).
$$
とおく。ここで，$j \in J$ に対応する成分 $\mathcal{O}_X$ 上の写像は，
切断 $\sum_i \alpha(r_{ij})$ によって与えられる。ただし $r_{ij}$ は，
$M$ の表示に現れる写像の行列係数である。
\item $\mathcal{F}_3$ を，前層
$U \mapsto \mathcal{O}_X(U) \otimes_R M$ に付随する層とする。ここで写像
$R \to \mathcal{O}_X(U)$ は，$\alpha$ と制限写像
$\mathcal{O}_X(X) \to \mathcal{O}_X(U)$ の合成である。
\end{enumerate}
この構成は次の性質をもつ：
\begin{enumerate}
\item 得られる $\mathcal{O}_X$-加群の層
$\mathcal{F}_M = \mathcal{F}_1 = \mathcal{F}_2 = \mathcal{F}_3$
は準連接である。
\item この構成は，$R$-加群の圏から $X$ 上の準連接層の圏への関手を与え，
この関手は任意の余極限と可換である。
\item 任意の $x \in X$ に対して，$M$ に関手的な等式
$\mathcal{F}_{M, x} = \mathcal{O}_{X, x} \otimes_R M$
が成り立つ。
\item 任意の $\mathcal{O}_X$-加群の層 $\mathcal{G}$ に対して
$$
\Mor_{\mathcal{O}_X}(\mathcal{F}_M, \mathcal{G})
=
\Hom_R(M, \Gamma(X, \mathcal{G}))
$$
が成り立つ。ここで $\Gamma(X, \mathcal{G})$ 上の $R$-加群構造は，
$\alpha$ を介して $\Gamma(X, \mathcal{O}_X)$-加群構造から得られる。
\end{enumerate}
\end{lemma}

\begin{proof}
$\mathcal{F}_1$ と $\mathcal{F}_3$ の間の同型は，$\pi^*$ が (3) の前層の
層化として定義されることから従う。《層》第
\ref{sheaves-section-ringed-spaces-functoriality-modules} 節を参照せよ。
(2) と (1) の構成の間の同型は，関手 $\pi^*$ が右完全であるため，
$\pi^*(\bigoplus_{j \in J} R) \to \pi^*(\bigoplus_{i \in I} R) \to
\pi^*M \to 0$ が完全となること，$\pi^*$ が任意の直和と可換であること
（補題 \ref{modules-lemma-exactness-pushforward-pullback} を参照），そして
$\pi^*(R) = \mathcal{O}_X$ であることから従う。

\medskip\noindent
(1) は構成 (2) から明らかである。(2) は $\pi^*$ がこれらの性質をもつため
明らかである。(3) は引き戻し層の茎の記述から従う。《層》補題
\ref{sheaves-lemma-stalk-pullback-modules} を参照せよ。
(4) は $\pi_*$ と $\pi^*$ の随伴性から従う。
\end{proof}

\begin{definition}
\label{modules-definition-sheaf-associated}
補題 \ref{modules-lemma-construct-quasi-coherent-sheaves} の状況で，
$\mathcal{F}_M$ を{\it 加群 $M$ と環準同型 $\alpha$ に付随する層}という。
$R = \Gamma(X, \mathcal{O}_X)$ かつ $\alpha = \text{id}_R$ の場合には，
$\mathcal{F}_M$ を単に{\it 加群 $M$ に付随する層}という。
\end{definition}


\begin{lemma}
\label{modules-lemma-restrict-quasi-coherent}
$(X, \mathcal{O}_X)$ を環付き空間とし，
$R = \Gamma(X, \mathcal{O}_X)$ とおく。$M$ を $R$-加群，
$\mathcal{F}_M$ を $M$ に付随する準連接 $\mathcal{O}_X$-加群の層とする。
$g : (Y, \mathcal{O}_Y) \to (X, \mathcal{O}_X)$ が環付き空間の射ならば，
$g^*\mathcal{F}_M$ は $\Gamma(Y, \mathcal{O}_Y)$-加群
$\Gamma(Y, \mathcal{O}_Y) \otimes_R M$ に付随する層である。
\end{lemma}

\begin{proof}
この主張は，補題 \ref{modules-lemma-construct-quasi-coherent-sheaves} における
$\mathcal{F}_M$ の最初の記述 $\pi^*M$ と，次の環付き空間の可換図式から従う：
$$
\xymatrix{
(Y, \mathcal{O}_Y) \ar[r]_-\pi \ar[d]_g &
(\{*\}, \Gamma(Y, \mathcal{O}_Y)) \ar[d]^{g^\sharp\text{ により誘導}} \\
(X, \mathcal{O}_X) \ar[r]^-\pi &
(\{*\}, \Gamma(X, \mathcal{O}_X))
}
$$
（《層》補題 \ref{sheaves-lemma-push-pull-composition-modules} も用いる。）
\end{proof}

\begin{lemma}
\label{modules-lemma-quasi-coherent-module}
$(X, \mathcal{O}_X)$ を環付き空間とし，$x \in X$ を点とする。
$x$ が準コンパクトな近傍からなる基本近傍系をもつと仮定する。
$\mathcal{F}$ を任意の準連接 $\mathcal{O}_X$-加群の層とする。
このとき $x$ の開近傍 $U$ が存在し，$\mathcal{F}|_U$ は，ある
$\Gamma(U, \mathcal{O}_U)$-加群 $M$ に付随する
$(U, \mathcal{O}_U)$ 上の加群の層 $\mathcal{F}_M$ と同型である。
\end{lemma}

\begin{proof}
まず $X$ を $x$ の開近傍で置き換え，$\mathcal{F}$ は写像
$$
\Psi :
\bigoplus\nolimits_{j \in J}
\mathcal{O}_X
\longrightarrow
\bigoplus\nolimits_{i \in I}
\mathcal{O}_X.
$$
の余核と同型であると仮定してよい。問題は，直和の大域切断の加群が
一般には大域切断の加群の直和と異なるため，この写像が「行列」で
与えられるとは限らないことである。

\medskip\noindent
$x \in E \subset X$ を $x$ の準コンパクトな近傍とする
（注意：$E$ は開集合とは限らない）。さらに，$E$ に含まれる $x$ の開近傍を
$x \in U \subset E$ とする。次に，補題
\ref{modules-lemma-section-direct-sum-quasi-compact} の証明と同様に進める。
各 $j \in J$ に対し，$j$ に対応する直和因子 $\mathcal{O}_X$ の切断 $1$ の像を
$s_j \in \Gamma(X, \bigoplus\nolimits_{i \in I} \mathcal{O}_X)$ と書く。
$E \subset \bigcup_{k \in K_j} U_{jk}$ を満たす有限個の開集合
$U_{jk}$（$k \in K_j$）が存在し，各制限 $s_j|_{U_{jk}}$ は有限和
$\sum_{i \in I_{jk}} f_{jki}$ となる。ここで $I_{jk} \subset I$ であり，
$f_{jki}$ は $i \in I$ に対応する直和因子 $\mathcal{O}_X$ に属する。
$I_j = \bigcup_{k \in K_j} I_{jk}$ とおく。これは有限集合である。
$U \subset E \subset \bigcup_{k \in K_j} U_{jk}$ なので，切断 $s_j|_U$ は
有限直和 $\bigoplus_{i \in I_j} \mathcal{O}_X$ の切断である。
補題 \ref{modules-lemma-limits-colimits} により，実際
$s_j|_U = \sum_{i \in I_j} f_{ij}$ と書け，
$f_{ij} \in \mathcal{O}_X(U) = \Gamma(U, \mathcal{O}_U)$ である。

\medskip\noindent
ここで，加群 $M$ を写像
$$
\bigoplus\nolimits_{j \in J}
\Gamma(U, \mathcal{O}_U)
\longrightarrow
\bigoplus\nolimits_{i \in I}
\Gamma(U, \mathcal{O}_U)
$$
の余核として定義できる。この写像の行列は $(f_{ij})$ で与えられる。
補題 \ref{modules-lemma-construct-quasi-coherent-sheaves} の構成 (2) により，
$\mathcal{F}_M$ は $\mathcal{F}|_U$ と同じ表示をもつ。したがって
$\mathcal{F}_M \cong \mathcal{F}|_U$ である。
\end{proof}

\begin{example}
\label{modules-example-quasi-coherent}
$X$ を，実数直線の可算個のコピー $L_1, L_2, L_3, \ldots$ をすべて
$0$ で貼り合わせた空間とする。$0$ の基本近傍系は
$\{U_n\}_{n \in \mathbf{N}}$ であり，$U_n \cap L_i = (-1/n, 1/n)$ とする。
$\mathcal{O}_X$ を実数値連続関数の層とする。
$f : \mathbf{R} \to \mathbf{R}$ を，$(-1, 1)$ 上で恒等的に零，
$(-\infty, -2) \cup (2, \infty)$ 上で恒等的に $1$ となる連続関数とする。
$X$ 上の連続関数で，各 $L_j = \mathbf{R}$ 上では
$x \mapsto f(nx)$ に等しいものを $f_n$ と書く。また $1_{L_j}$ を
$L_j$ の特性関数とする。写像
$$
\bigoplus\nolimits_{j \in \mathbf{N}}
\mathcal{O}_X
\longrightarrow
\bigoplus\nolimits_{j, i \in \mathbf{N}}
\mathcal{O}_X, \quad
e_j \longmapsto \sum\nolimits_{i \in \mathbf{N}} f_j 1_{L_i} e_{ij}
$$
を考える。記法は明らかであろう。$f_j$ は $0$ のある近傍で零なので，
この和は局所有限であり，したがって意味をもつ。$U_n$ 上では，
$j > 2n$ に対する $e_j$ の像は，$g_{ij}$ が連続であるような有限線形結合
$\sum g_{ij} e_{ij}$ ではない。したがって，上の補題
\ref{modules-lemma-quasi-coherent-module} の証明における意味で，表示した写像が
「行列」によって与えられるような $0 \in X$ の近傍は存在しない。

\medskip\noindent
$\bigoplus\nolimits_{j \in \mathbf{N}} \mathcal{O}_X$ は，基底 $e_j$ をもつ
自由加群に付随する層であり，もう一方の直和についても同様である。
したがって，加群に付随する層の射は，一般には $X$ 上で局所的にさえ
加群の射から来るとは限らない。同様に，環付き空間 $X$ と準連接
$\mathcal{O}_X$-加群の層 $\mathcal{F}$ であって，$\mathcal{F}$ が局所的にも
$\mathcal{F}_M$ の形にならない例が存在するはずである。
（そのような例を見つけた場合はメールで知らせてほしい。）
さらに，局所コンパクト空間 $X$ と写像
$\mathcal{F}_M \to \mathcal{F}_N$ であって，局所的にも加群の写像から来ない
例も存在するはずである（補題 \ref{modules-lemma-quasi-coherent-module} の証明から，
$N$ が自由ならばこれは起こらないことが分かる）。
\end{example}












\section{有限表示加群の層}
\label{modules-section-finite-presentation}

\noindent
定義は次のとおりである。

\begin{definition}
\label{modules-definition-finite-presentation}
$(X, \mathcal{O}_X)$ を環付き空間とし，$\mathcal{F}$ を
$\mathcal{O}_X$-加群の層とする。すべての点 $x \in X$ に対して
開近傍 $x\in U \subset X$ と $n, m \in \mathbf{N}$ が存在し，
$\mathcal{F}|_U$ が写像
$$
\bigoplus\nolimits_{j = 1, \ldots, m}
\mathcal{O}_U
\longrightarrow
\bigoplus\nolimits_{i = 1, \ldots, n}
\mathcal{O}_U
$$
の余核と同型であるとき，$\mathcal{F}$ は{\it 有限表示}であるという。
\end{definition}

\noindent
これは，$X$ が，その上で $\mathcal{F}|_U$ が次の形の
{\it 表示}をもつような開集合 $U$ で被覆されることを意味する：
$$
\bigoplus\nolimits_{j = 1, \ldots, m}
\mathcal{O}_U
\longrightarrow
\bigoplus\nolimits_{i = 1, \ldots, n}
\mathcal{O}_U
\to
\mathcal{F}|_U
\to
0.
$$
ここで「表示」とは，表示された列が完全であることを意味する。言い換えると，
\begin{enumerate}
\item $X$ の各点 $x$ に対して開近傍 $U$ が存在し，
$\mathcal{F}|_U$ は有限個の大域切断で生成され，かつ
\item これらの切断を適切に選べば，対応する全射の核も有限個の
大域切断で生成される。
\end{enumerate}

\begin{lemma}
\label{modules-lemma-finite-presentation-quasi-coherent}
$(X, \mathcal{O}_X)$ を環付き空間とする。任意の有限表示
$\mathcal{O}_X$-加群の層は準連接である。
\end{lemma}

\begin{proof}
定義から直ちに従う。
\end{proof}

\begin{lemma}
\label{modules-lemma-cokernel-finite-finite-presentation}
$(X,\mathcal{O}_X)$ を環付き空間とし，$\mathcal{F}$ を有限表示
$\mathcal{O}_X$-加群の層とする。また
$\varphi : \mathcal{G} \to \mathcal{F}$ を $\mathcal{O}_X$-加群の層の射とする。
$\mathcal{G}$ が有限型ならば，$\Coker(\varphi)$ は有限表示である。
\end{lemma}

\begin{proof}
$X$ 上で局所的に
$\mathcal{F} = \mathcal{O}_X^{\oplus n} / \mathcal{M}$ と書ける。
ここで $\mathcal{M} \subset \mathcal{O}_X^{\oplus n}$ は有限型
$\mathcal{O}_X$-部分加群の層である。したがって
$\Im(\varphi) = \mathcal{N} / \mathcal{M}$ と書ける。ここで
$\mathcal{N} \subset \mathcal{O}_X^{\oplus n}$ は $\mathcal{M}$ を含む
$\mathcal{O}_X$-部分加群の層である。$\mathcal{G} \to \Im(\varphi)$ は
全射で $\mathcal{G}$ は有限型だから，$\mathcal{O}_X$-加群の層
$\Im(\varphi)$ は有限型である。補題 \ref{modules-lemma-extension-finite-type}
により $\mathcal{N}$ は有限型である。ゆえに
$\Coker(\varphi) = \mathcal{O}_X^{\oplus n} / \mathcal{N}$ は有限表示である。
\end{proof}

\begin{lemma}
\label{modules-lemma-kernel-surjection-finite-free-onto-finite-presentation}
$(X, \mathcal{O}_X)$ を環付き空間とし，$\mathcal{F}$ を有限表示
$\mathcal{O}_X$-加群の層とする。
\begin{enumerate}
\item $\psi : \mathcal{O}_X^{\oplus r} \to \mathcal{F}$ が全射ならば，
$\Ker(\psi)$ は有限型である。
\item $\theta : \mathcal{G} \to \mathcal{F}$ が全射で，$\mathcal{G}$ が
有限型ならば，$\Ker(\theta)$ は有限型である。
\end{enumerate}
\end{lemma}

\begin{proof}
(1) を証明する。$x \in X$ とする。表示
$$
\mathcal{O}_U^{\oplus m}
\xrightarrow{\chi}
\mathcal{O}_U^{\oplus n}
\xrightarrow{\varphi}
\mathcal{F}|_U
\to
0.
$$
が存在するような $x$ の開近傍 $U \subset X$ を選ぶ。
$e_k$ を $\mathcal{O}_X^{\oplus r}$ の第 $k$ 因子を生成する切断とする。
各 $k = 1, \ldots, r$ に対し，$U$ を $x$ の十分小さい近傍へ縮小すれば，
$\psi(e_k)$ を $U$ 上の $\mathcal{O}_U^{\oplus n}$ の切断
$\tilde e_k$ に持ち上げることができる。これにより，
$\varphi \circ \alpha = \psi$ を満たす層の射
$\alpha : \mathcal{O}_U^{\oplus r} \to \mathcal{O}_U^{\oplus n}$ を得る。
同様に，$U$ を縮小すれば，$\psi \circ \beta = \varphi$ を満たす射
$\beta : \mathcal{O}_U^{\oplus n} \to \mathcal{O}_U^{\oplus r}$ をとれる。
すると写像
$$
\mathcal{O}_U^{\oplus m} \oplus
\mathcal{O}_U^{\oplus r}
\xrightarrow{\beta \circ \chi, 1 - \beta \circ \alpha}
\mathcal{O}_U^{\oplus r}
$$
は $\psi$ の核への全射である。

\medskip\noindent
(2) を証明するには，局所的に全射
$\eta : \mathcal{O}_X^{\oplus r} \to \mathcal{G}$ を選べばよい。
(1) により $\Ker(\theta \circ \eta)$ は有限型である。
$\Ker(\theta) = \eta(\Ker(\theta \circ \eta))$ なので結論を得る。
\end{proof}

\begin{lemma}
\label{modules-lemma-pullback-finite-presentation}
$f : (X, \mathcal{O}_X) \to (Y, \mathcal{O}_Y)$ を環付き空間の射とする。
有限表示加群の層 $\mathcal{G}$ の引き戻し $f^*\mathcal{G}$ は有限表示である。
\end{lemma}

\begin{proof}
補題 \ref{modules-lemma-pullback-quasi-coherent} の証明とまったく同じであるが，
添字集合を有限とする。
\end{proof}

\begin{lemma}
\label{modules-lemma-quasi-coherent-limit-finite-presentation}
$(X, \mathcal{O}_X)$ を環付き空間とし，
$R = \Gamma(X, \mathcal{O}_X)$ とおく。$M$ を $R$-加群とする。
$M$ に付随する $\mathcal{O}_X$-加群の層 $\mathcal{F}_M$ は，
有限表示 $\mathcal{O}_X$-加群の層の有向余極限である。
\end{lemma}

\begin{proof}
これは補題 \ref{modules-lemma-construct-quasi-coherent-sheaves} と，任意の加群が
有限表示加群の有向余極限であるという事実から直ちに従う。
《代数》補題 \ref{algebra-lemma-module-colimit-fp} を参照せよ。
\end{proof}

\begin{lemma}
\label{modules-lemma-finite-presentation-stalk-free}
$(X, \mathcal{O}_X)$ を環付き空間とし，$\mathcal{F}$ を有限表示
$\mathcal{O}_X$-加群の層とする。
$\mathcal{F}_x \cong \mathcal{O}_{X, x}^{\oplus r}$ となる $x \in X$
をとる。このとき，$\mathcal{F}|_U \cong \mathcal{O}_U^{\oplus r}$
となる $x$ の開近傍 $U$ が存在する。
\end{lemma}

\begin{proof}
与えられた同型のもとで $\mathcal{O}_{X, x}^{\oplus r}$ の基底に写る
$s_1, \ldots, s_r \in \mathcal{F}_x$ を選ぶ。各 $s_i$ が
$s_i \in \mathcal{F}(U)$ に持ち上がるような $x$ の開近傍 $U$ を選ぶ。
$U$ を縮小すれば，誘導される写像
$\psi : \mathcal{O}_U^{\oplus r} \to \mathcal{F}|_U$ は全射となる
（補題 \ref{modules-lemma-finite-type-surjective-on-stalk}）。補題
\ref{modules-lemma-kernel-surjection-finite-free-onto-finite-presentation} により，
$\Ker(\psi)$ は有限型である。すると $\Ker(\psi)_x = 0$ だから，
$U$ をもう一度縮小すれば $\Ker(\psi)$ は零になる
（補題 \ref{modules-lemma-finite-type-stalk-zero}）。
\end{proof}





\section{連接加群の層}
\label{modules-section-coherent}

\noindent
本節の参考文献は \cite{FAC} である。

\medskip\noindent
環付き空間 $X$ 上の連接層の圏は，準連接層の圏より扱いやすい対象である。
というのも，$X$ がどのような空間であっても，少なくとも
$\textit{Mod}(\mathcal{O}_X)$ のアーベル部分圏だからである。一方，
環付き空間という一般の状況では，連接加群の層の引き戻しは
「ほとんど決して」連接にならない。

\begin{definition}
\label{modules-definition-coherent}
$(X, \mathcal{O}_X)$ を環付き空間とし，$\mathcal{F}$ を
$\mathcal{O}_X$-加群の層とする。次の二条件が成り立つとき，
$\mathcal{F}$ を{\it 連接 $\mathcal{O}_X$-加群の層}という：
\begin{enumerate}
\item $\mathcal{F}$ は有限型であり，かつ
\item 任意の開集合 $U \subset X$ と任意の有限族
$s_i \in \mathcal{F}(U)$（$i = 1, \ldots, n$）に対して，付随する写像
$\bigoplus_{i = 1, \ldots, n} \mathcal{O}_U \to \mathcal{F}|_U$
の核は有限型である。
\end{enumerate}
連接 $\mathcal{O}_X$-加群の層の圏を $\textit{Coh}(\mathcal{O}_X)$ と書く。
\end{definition}

\begin{lemma}
\label{modules-lemma-coherent-finite-presentation}
$(X, \mathcal{O}_X)$ を環付き空間とする。任意の連接
$\mathcal{O}_X$-加群の層は有限表示であり，したがって準連接である。
\end{lemma}

\begin{proof}
$\mathcal{F}$ を $X$ 上の連接層とし，点 $x \in X$ をとる。
連接性の定義の (1) により，開近傍 $U$ と $U$ 上の $\mathcal{F}$ の切断
$s_i$（$i = 1, \ldots, n$）であって，
$\Psi : \bigoplus_{i = 1, \ldots, n} \mathcal{O}_U \to \mathcal{F}$
が全射となるものを選べる。連接性の定義の (2) により，開近傍
$x \in V \subset U$ と $\bigoplus_{i = 1, \ldots, n} \mathcal{O}_V$ の切断
$t_1, \ldots, t_m$ であって，$\Psi|_V$ の核を生成するものを選べる。
すると $V$ 上で表示
$$
\bigoplus\nolimits_{j = 1, \ldots, m}
\mathcal{O}_V
\longrightarrow
\bigoplus\nolimits_{i = 1, \ldots, n}
\mathcal{O}_V
\to
\mathcal{F}|_V
\to
0
$$
を得る。これが求めるものである。
\end{proof}

\begin{example}
\label{modules-example-coherent-not-Noetherian}
$X$ が一点であると仮定する。この場合，上の定義は環上の加群に対する概念を
与える。連接性の定義は何を意味するだろうか。これはネーター性の概念と
密接に関係するが，同じではない。実際，環
$R = \mathbf{C}[x_1, x_2, x_3, \ldots]$ は自身上の加群として連接であるが，
自身上の加群としてネーターではない。詳しくは《代数》第
\ref{algebra-section-coherent} 節を参照せよ。
\end{example}

\begin{lemma}
\label{modules-lemma-coherent-abelian}
$(X, \mathcal{O}_X)$ を環付き空間とする。
\begin{enumerate}
\item 連接層の任意の有限型部分層は連接である。
\item $\varphi : \mathcal{F} \to \mathcal{G}$ を，有限型層
$\mathcal{F}$ から連接層 $\mathcal{G}$ への射とする。このとき
$\Ker(\varphi)$ は有限型である。
\item $\varphi : \mathcal{F} \to \mathcal{G}$ を連接
$\mathcal{O}_X$-加群の層の射とする。このとき
$\Ker(\varphi)$ と $\Coker(\varphi)$ は連接である。
\item $\mathcal{O}_X$-加群の層の短完全列
$0 \to \mathcal{F}_1 \to \mathcal{F}_2 \to \mathcal{F}_3 \to 0$
が与えられたとき，三つのうち二つが連接ならば残りも連接である。
\item 圏 $\textit{Coh}(\mathcal{O}_X)$ は
$\textit{Mod}(\mathcal{O}_X)$ の弱 Serre 部分圏である。特に，
連接加群の層の圏はアーベル圏であり，包含関手
$\textit{Coh}(\mathcal{O}_X) \to \textit{Mod}(\mathcal{O}_X)$
は完全である。
\end{enumerate}
\end{lemma}

\begin{proof}
定義 \ref{modules-definition-coherent} の条件 (2) は，連接層の任意の部分層について
成り立つ。したがって (1) を得る。

\medskip\noindent
(2) の仮定をおく。$\Ker(\varphi)$ が有限型であることを示そう。
$x \in X$ をとり，$\mathcal{F}|_U$ が $s_1, \ldots, s_n$ で
生成されるような $X$ における $x$ の開近傍 $U$ を選ぶ。
定義 \ref{modules-definition-coherent} により，誘導される写像
$\bigoplus_{i = 1}^n \mathcal{O}_U \to \mathcal{G}$,
$e_i \mapsto \varphi(s_i)$ の核 $\mathcal{K}$ は有限型である。
したがって，合成
$\mathcal{K} \to \bigoplus_{i = 1}^n \mathcal{O}_U \to \mathcal{F}$
の像である $\Ker(\varphi)$ も有限型である。

\medskip\noindent
(3) の仮定をおく。(2) により $\varphi$ の核は有限型であり，
したがって (1) により連接である。

\medskip\noindent
同じ仮定のもとで，$\Coker(\varphi)$ が連接であることを示そう。
$\mathcal{G}$ は有限型なので $\Coker(\varphi)$ も有限型である。
$U \subset X$ を開集合とし，
$\overline{s}_i \in \Coker(\varphi)(U)$（$i = 1, \ldots, n$）を切断とする。
付随する射
$\overline{\Psi} : \bigoplus_{i = 1}^n \mathcal{O}_U \to \Coker(\varphi)$
の核が有限型であることを示さなければならない。$U$ の開被覆で，
各開集合上ですべての切断 $\overline{s}_i$ が $\mathcal{G}$ の切断
$s_i$ に持ち上がるものが存在する。したがって，$U$ 上でそうなっていると
仮定してよい。さらに，$U$ 上の $\Im(\varphi)$ の切断
$t_j$（$j = 1, \ldots, m$）であって，$U$ 上で $\Im(\varphi)$ を
生成するものが存在すると仮定してよい。$t_j$ を用いて
$\Phi : \bigoplus_{j = 1}^m \mathcal{O}_U \to \Im(\varphi)$ を定義し，
$t_j$ と $s_i$ を用いて
$\Psi :
\bigoplus_{j = 1}^m \mathcal{O}_U \oplus
\bigoplus_{i = 1}^n \mathcal{O}_U \to \mathcal{G}$
を定義する。次の可換図式を考える：
$$
\xymatrix{
0 \ar[r] &
\bigoplus_{j = 1}^m \mathcal{O}_U \ar[d]_\Phi \ar[r] &
\bigoplus_{j = 1}^m \mathcal{O}_U \oplus
\bigoplus_{i = 1}^n \mathcal{O}_U \ar[d]_\Psi \ar[r] &
\bigoplus_{i = 1}^n \mathcal{O}_U \ar[d]_{\overline{\Psi}} \ar[r] &
0 \\
0 \ar[r] &
\Im(\varphi) \ar[r] &
\mathcal{G} \ar[r] &
\Coker(\varphi) \ar[r] &
0
}
$$
蛇の補題により完全列
$\Ker(\Psi) \to \Ker(\overline{\Psi}) \to 0$ を得る。
$\Ker(\Psi)$ は有限型加群の層なので，$\Ker(\overline{\Psi})$ も有限型である。

\medskip\noindent
(4) を証明する。
$0 \to \mathcal{F}_1 \to \mathcal{F}_2 \to \mathcal{F}_3 \to 0$
を $\mathcal{O}_X$-加群の層の短完全列とする。(3) により，
$\mathcal{F}_1$ と $\mathcal{F}_3$ が連接ならば $\mathcal{F}_2$ も
連接であることを示せば十分である。補題 \ref{modules-lemma-extension-finite-type}
により $\mathcal{F}_2$ は有限型である。$s_1, \ldots, s_n$ を，
$X$ の共通の開集合 $U$ 上で定義された $\mathcal{F}_2$ の有限個の
局所切断とする。それらの間の関係の加群の層 $\mathcal{K}$ が有限型であることを
示さなければならない。次の可換図式を考える：
$$
\xymatrix{
0 \ar[r] &
0 \ar[r] \ar[d] &
\bigoplus_{i = 1}^{n} \mathcal{O}_U \ar[r] \ar[d] &
\bigoplus_{i = 1}^{n} \mathcal{O}_U \ar[r] \ar[d] &
0 \\
0 \ar[r] &
\mathcal{F}_1 \ar[r] &
\mathcal{F}_2 \ar[r] &
\mathcal{F}_3 \ar[r] &
0
}
$$
記法は明らかであろう。蛇の補題により短完全列
$0 \to \mathcal{K} \to \mathcal{K}_3 \to \mathcal{F}_1$
を得る。ここで $\mathcal{K}_3$ は，切断 $s_i$ の $\mathcal{F}_3$ における
像の間の関係の加群の層である。$\mathcal{F}_1$ は連接なので，
$\mathcal{K}$ は有限型加群の層から連接加群の層への写像の核であり，
したがって (2) により有限型である。

\medskip\noindent
(5) を証明する。(3) と (4) により《ホモロジー》補題
\ref{homology-lemma-characterize-weak-serre-subcategory} を適用できるので，
主張が従う。
\end{proof}

\begin{lemma}
\label{modules-lemma-coherent-structure-sheaf}
$(X, \mathcal{O}_X)$ を環付き空間とし，$\mathcal{F}$ を
$\mathcal{O}_X$-加群の層とする。$\mathcal{O}_X$ が連接
$\mathcal{O}_X$-加群の層であると仮定する。このとき，
$\mathcal{F}$ が連接であることと有限表示であることは同値である。
\end{lemma}

\begin{proof}
省略する。
\end{proof}

\begin{lemma}
\label{modules-lemma-finite-type-to-coherent-injective-on-stalk}
$X$ を環付き空間とし，
$\varphi : \mathcal{G} \to \mathcal{F}$ を $\mathcal{O}_X$-加群の層の
準同型とする。$x \in X$ とする。$\mathcal{G}$ は有限型，
$\mathcal{F}$ は連接であり，茎上の写像
$\varphi_x : \mathcal{G}_x \to \mathcal{F}_x$ は単射であると仮定する。
このとき，$\varphi|_U$ が単射となる開近傍 $x \in U \subset X$ が存在する。
\end{lemma}

\begin{proof}
$\varphi$ の核を $\mathcal{K} \subset \mathcal{G}$ と書く。
補題 \ref{modules-lemma-coherent-abelian} により $\mathcal{K}$ は有限型
$\mathcal{O}_X$-加群の層である。仮定は $\mathcal{K}_x = 0$ ということである。
補題 \ref{modules-lemma-finite-type-stalk-zero} により，$\mathcal{K}|_U = 0$
となる $x$ の開近傍 $U$ が存在する。この $U$ が条件を満たす。
\end{proof}













\section{環付き空間の閉埋め込み}
\label{modules-section-closed-immersion}

\noindent
環付き空間の射
$i : (Z, \mathcal{O}_Z) \to (X, \mathcal{O}_X)$ を，いつ閉埋め込みと
呼ぶべきだろうか。

\medskip\noindent
正規位相空間（連続関手の層を備える）の閉埋め込みや，微分可能多様体
（微分可能関数の層を備える）の閉埋め込みの例に動機づけられ，少なくとも
次を仮定するのが自然に思われる：
\begin{enumerate}
\item 写像 $i$ は位相空間の閉埋め込みである。
\item 付随する写像 $\mathcal{O}_X \to i_*\mathcal{O}_Z$ は全射である。
その核を $\mathcal{I}$ と書く。
\end{enumerate}
これらの条件だけでも，いくつもの好ましい結果が従う。例えば，第
\ref{modules-section-closed-immersions} 節のアーベル群の層に関する結果を一般化して，
$\mathcal{O}_Z$-加群の層の圏が，$\mathcal{I}$ で消去される
$\mathcal{O}_X$-加群の層の圏と同値であることを証明する。

\medskip\noindent
しかし Stacks Project では，$i$ が閉埋め込みで
$(X, \mathcal{O}_X)$ がスキームならば $(Z, \mathcal{O}_Z)$ もスキームに
なることを保証する定義を選ぶ。さらにこの状況では，$i_*$ と $i^*$ が，
準連接 $\mathcal{O}_Z$-加群の層の圏と，$\mathcal{I}$ で消去される準連接
$\mathcal{O}_X$-加群の層の圏との間の同値を与えることを望む。
最低限必要な条件は，$i_*\mathcal{O}_Z$ が準連接
$\mathcal{O}_X$-加群の層であることである。これを保証する良い方法は，
$\mathcal{I}$ が局所的に切断で生成されると仮定することである。
この条件は，「$(Z, \mathcal{O}_Z)$ は $(X, \mathcal{O}_X)$ 上で局所的に，
いくつかの正則関数 $f_i$，すなわち $\mathcal{O}_X$ の局所切断を零とおくことで
定義される」と解釈できる。そこで次の定義に至る。

\begin{definition}
\label{modules-definition-closed-immersion}
{\it 環付き空間の閉埋め込み}\footnote{これは標準的でない用語である。
上の議論を参照せよ。}とは，次の性質をもつ射
$i : (Z, \mathcal{O}_Z) \to (X, \mathcal{O}_X)$
のことである：
\begin{enumerate}
\item 写像 $i$ は位相空間の閉埋め込みである。
\item 付随する写像 $\mathcal{O}_X \to i_*\mathcal{O}_Z$ は全射である。
その核を $\mathcal{I}$ と書く。
\item $\mathcal{O}_X$-加群の層 $\mathcal{I}$ は局所的に切断で生成される。
\end{enumerate}
\end{definition}

\noindent
実際，この定義だけでは，準連接 $\mathcal{O}_Z$-加群の層の $i_*$ が
準連接 $\mathcal{O}_X$-加群の層になることはまだ保証されない。
問題は，準連接 $\mathcal{O}_Z$-加群の層の局所表示を，その順像の局所表示へ
どのように変換すればよいか明らかでないことである。しかし次は自明である。

\begin{lemma}
\label{modules-lemma-i-star-quasi-coherent}
$i : (Z, \mathcal{O}_Z) \to (X, \mathcal{O}_X)$ を環付き空間の
閉埋め込みとし，$\mathcal{F}$ を準連接 $\mathcal{O}_Z$-加群の層とする。
このとき $i_*\mathcal{F}$ は $X$ 上で局所的に，準連接
$\mathcal{O}_X$-加群の層の射の余核である。
\end{lemma}

\begin{proof}
定義により $i_*\mathcal{O}_Z$ は準連接だからである。また $Z$ 上で局所的に，
層 $\mathcal{F}$ は $\mathcal{O}_Z$ のコピーの直和の間の射の余核である。
さらに，{\it 同一の}準連接層のコピーの任意の直和は準連接である。
最後に $i_*$ は任意の余極限と可換である。補題
\ref{modules-lemma-i-star-right-adjoint} を参照せよ。いくつかの詳細は省略する。
\end{proof}

\begin{lemma}
\label{modules-lemma-i-star-reflects-finite-type}
$i : (Z, \mathcal{O}_Z) \to (X, \mathcal{O}_X)$ を環付き空間の射とする。
$i$ は $X$ の閉部分集合への同相写像であり，
$\mathcal{O}_X \to i_*\mathcal{O}_Z$ は全射であると仮定する。
$\mathcal{F}$ を $\mathcal{O}_Z$-加群の層とする。このとき，
$i_*\mathcal{F}$ が有限型であることと $\mathcal{F}$ が有限型であることは同値である。
\end{lemma}

\begin{proof}
$\mathcal{F}$ が有限型であると仮定し，$x \in X$ をとる。
$x \not \in Z$ ならば，$i_*\mathcal{F}$ は $x$ のある近傍上で零であり，
したがってその近傍上で有限生成である。$x = i(z)$ ならば，開近傍
$z \in V \subset Z$ と $V$ 上で $\mathcal{F}$ を生成する切断
$s_1, \ldots, s_n \in \mathcal{F}(V)$ を選ぶ。ある開集合 $U \subset X$
に対して $V = Z \cap U$ と書く。$U$ は $x$ の近傍であることに注意する。
切断 $s_i$ は明らかに $U$ 上の $i_*\mathcal{F}$ の切断 $s_i$ を与える。
得られる写像
$$
\bigoplus\nolimits_{i = 1, \ldots, n} \mathcal{O}_U
\longrightarrow
i_*\mathcal{F}|_U
$$
は，茎上の作用を調べれば全射である（ここで
$\mathcal{O}_X \to i_*\mathcal{O}_Z$ が全射であることを用いる）。
したがって $i_*\mathcal{F}$ は有限型である。

\medskip\noindent
逆に $i_*\mathcal{F}$ が有限型であると仮定する。$z \in Z$ を選び，
$x = i(z)$ とおく。仮定により，$x$ の開近傍 $U \subset X$ と，
$U$ 上で $i_*\mathcal{F}$ を生成する切断
$s_1, \ldots, s_n \in (i_*\mathcal{F})(U)$ が存在する。
$V = Z \cap U$ とおく。$i_*$ の定義により，切断 $s_i$ は $V$ 上の
$\mathcal{F}$ の切断 $s_i$ に対応する。得られる写像
$$
\bigoplus\nolimits_{i = 1, \ldots, n} \mathcal{O}_V
\longrightarrow
\mathcal{F}|_V
$$
は，茎上の作用を調べれば全射である。したがって $\mathcal{F}$ は有限型である。
\end{proof}

\begin{lemma}
\label{modules-lemma-i-star-equivalence}
$i : (Z, \mathcal{O}_Z) \to (X, \mathcal{O}_X)$ を環付き空間の射とする。
$i$ は $X$ の閉部分集合への同相写像であり，
$i^\sharp : \mathcal{O}_X \to i_*\mathcal{O}_Z$ は全射であると仮定する。
$i^\sharp$ の核を $\mathcal{I} \subset \mathcal{O}_X$ と書く。関手
$$
i_* :
\textit{Mod}(\mathcal{O}_Z)
\longrightarrow
\textit{Mod}(\mathcal{O}_X)
$$
は完全かつ充満忠実であり，その本質像は
$\mathcal{I}\mathcal{G} = 0$ を満たす $\mathcal{O}_X$-加群の層
$\mathcal{G}$ 全体である。
\end{lemma}

\begin{proof}
$\mathcal{O}_Z$-加群の層 $\mathcal{F}$ に対する標準写像
$$
i^*i_*\mathcal{F} \longrightarrow \mathcal{F}
$$
が同型であることを主張する。これを茎で確かめる。$z \in Z$ とし，
$x = i(z)$ とおく。このとき
$$
(i^*i_*\mathcal{F})_z =
(i_*\mathcal{F})_x \otimes_{\mathcal{O}_{X, x}} \mathcal{O}_{Z, z} =
\mathcal{F}_z \otimes_{\mathcal{O}_{X, x}} \mathcal{O}_{Z, z} =
\mathcal{F}_z
$$
である。これは《層》補題 \ref{sheaves-lemma-stalk-pullback-modules}，
$\mathcal{O}_{Z, z}$ が $\mathcal{O}_{X, x}$ の商であること，および
《層》補題 \ref{sheaves-lemma-stalks-closed-pushforward} から従う。
したがって $i_*$ は充満忠実である。

\medskip\noindent
$\mathcal{I}\mathcal{G} = 0$ を満たす $\mathcal{O}_X$-加群の層
$\mathcal{G}$ をとる。標準写像
$$
\mathcal{G} \longrightarrow i_*i^*\mathcal{G}
$$
が同型であることを示す。これにより
$\mathcal{F} = i^*\mathcal{G}$ として $\mathcal{G} = i_*\mathcal{F}$ が分かり，
証明が完了する。表示した写像が茎上で同型を誘導することを確かめる。
$x \in X$ かつ $x \not \in i(Z)$ ならば，この場合
$\mathcal{I}_x = \mathcal{O}_{X, x}$ なので $\mathcal{G}_x = 0$ である。
上と同様に，《層》補題 \ref{sheaves-lemma-stalks-closed-pushforward}
により $(i_*i^*\mathcal{G})_x = 0$ である。一方，$x \in Z$ ならば写像
$$
\mathcal{G}_x
\longrightarrow
\mathcal{G}_x \otimes_{\mathcal{O}_{X, x}} \mathcal{O}_{Z, x}
$$
を得る。これは《層》補題 \ref{sheaves-lemma-stalk-pullback-modules}
および \ref{sheaves-lemma-stalks-closed-pushforward} による。
$\mathcal{O}_{Z, x} = \mathcal{O}_{X, x}/\mathcal{I}_x$ であり，仮定により
$\mathcal{G}_x$ は $\mathcal{I}_x$ で消去されるので，この写像は同型である。
\end{proof}

\begin{remark}
\label{modules-remark-sections-support-in-closed-modules}
$(X, \mathcal{O}_X)$ を環付き空間とし，$Z \subset X$ を閉部分集合とする。
$\mathcal{O}_X$-加群の層 $\mathcal{F}$ に対し，{\it $Z$ に台をもつ切断の
部分加群の層} $\mathcal{H}_Z(\mathcal{F})$ を，次の規則で定義できる：
$$
\mathcal{H}_Z(\mathcal{F})(U) =
\{s \in \mathcal{F}(U) \mid \text{Supp}(s) \subset U \cap Z\}
$$
$\mathcal{H}_Z(\mathcal{F})(U)$ は $\mathcal{O}_X(U)$ 上の加群である。
すなわち $\mathcal{H}_Z(\mathcal{F})$ は $\mathcal{O}_X$-加群の層である。
構成により，$\mathcal{H}_Z(\mathcal{F})$ は，台が $Z$ に含まれる
$\mathcal{F}$ の $\mathcal{O}_X$-部分加群の層のうち最大である。
補題 \ref{modules-lemma-i-star-equivalence} を環付き空間の射
$(Z, \mathcal{O}_X|_Z) \to (X, \mathcal{O}_X)$ に適用すると，
$\mathcal{H}_Z(\mathcal{F})$ を $Z$ 上の $\mathcal{O}_X|_Z$-加群の層と
みなすことができ，実際そうみなす。したがって関手
$$
\textit{Mod}(\mathcal{O}_X) \longrightarrow \textit{Mod}(\mathcal{O}_X|_Z),
\quad
\mathcal{F} \longmapsto \mathcal{H}_Z(\mathcal{F})
\text{（}\mathcal{O}_X|_Z\text{-加群の層として }Z\text{ 上でみなす）}
$$
を得る。この関手は左完全であるが，一般には完全でない。
上の主張はすべて補題 \ref{modules-lemma-support-section-closed} から直接従う。
この構成が注意 \ref{modules-remark-sections-support-in-closed} の構成と
両立することは明らかである。
\end{remark}

\begin{lemma}
\label{modules-lemma-adjoint-section-with-support}
$(X, \mathcal{O}_X)$ を環付き空間とし，$i : Z \to X$ を閉部分集合の
包含とする。注意 \ref{modules-remark-sections-support-in-closed-modules} の関手
$\mathcal{H}_Z : \textit{Mod}(\mathcal{O}_X) \to
\textit{Mod}(\mathcal{O}_X|_Z)$ は，
$i_* : \textit{Mod}(\mathcal{O}_X|_Z) \to \textit{Mod}(\mathcal{O}_X)$
の右随伴である。
\end{lemma}

\begin{proof}
任意の $\mathcal{O}_X$-加群の層 $\mathcal{F}$ と任意の
$\mathcal{O}_X|_Z$-加群の層 $\mathcal{G}$ に対して
$$
\Hom_{\mathcal{O}_X|_Z}(\mathcal{G}, \mathcal{H}_Z(\mathcal{F})) =
\Hom_{\mathcal{O}_X}(i_*\mathcal{G}, \mathcal{F})
$$
を示せばよい。$i_*\mathcal{G}$ の任意の切断は $Z$ に台をもつので，
これは明らかである。詳細は省略する。
\end{proof}










\section{局所自由層}
\label{modules-section-locally-free}

\noindent
$(X, \mathcal{O}_X)$ を環付き空間とする。本書の規約では，
（いくつかの）茎 $\mathcal{O}_{X, x}$ が零環であることを許す。
したがって，局所自由層の階数を定義するときには少し注意が必要である。

\begin{definition}
\label{modules-definition-locally-free}
$(X, \mathcal{O}_X)$ を環付き空間とし，$\mathcal{F}$ を
$\mathcal{O}_X$-加群の層とする。
\begin{enumerate}
\item すべての点 $x \in X$ に対して集合 $I$ と開近傍
$x \in U \subset X$ が存在し，$\mathcal{F}|_U$ が
$\mathcal{O}_X|_U$-加群の層として
$\bigoplus_{i \in I} \mathcal{O}_X|_U$ と同型であるとき，
$\mathcal{F}$ は{\it 局所自由}であるという。
\item 添字集合 $I$ を有限に選べるとき，$\mathcal{F}$ は
{\it 有限局所自由}であるという。
\item 添字集合 $I$ の濃度を $r$ に選べるとき，$\mathcal{F}$ は
{\it 階数 $r$ の有限局所自由}であるという。
\end{enumerate}
\end{definition}

\noindent
（有限）局所自由層の有限直和は（有限）局所自由である。しかし，
局所自由層の無限直和は局所自由とは限らない。

\begin{lemma}
\label{modules-lemma-locally-free-quasi-coherent}
$(X, \mathcal{O}_X)$ を環付き空間とし，$\mathcal{F}$ を
$\mathcal{O}_X$-加群の層とする。$\mathcal{F}$ が局所自由ならば準連接である。
\end{lemma}

\begin{proof}
省略する。
\end{proof}

\begin{lemma}
\label{modules-lemma-pullback-locally-free}
$f : (X, \mathcal{O}_X) \to (Y, \mathcal{O}_Y)$ を環付き空間の射とする。
$\mathcal{G}$ が局所自由 $\mathcal{O}_Y$-加群の層ならば，
$f^*\mathcal{G}$ は局所自由 $\mathcal{O}_X$-加群の層である。
\end{lemma}

\begin{proof}
省略する。
\end{proof}

\begin{lemma}
\label{modules-lemma-rank}
$(X, \mathcal{O}_X)$ を環付き空間とする。$\mathcal{O}_X$ の台が $X$，
すなわち $\mathcal{O}_X$ のすべての茎が零でない環であると仮定する。
$\mathcal{F}$ を局所自由 $\mathcal{O}_X$-加群の層とする。
局所定数関数
$$
\text{rank}_\mathcal{F} :
X \longrightarrow \{0, 1, 2, \ldots\}\cup\{\infty\}
$$
であって，任意の点 $x \in X$ に対し，$x$ のある近傍上で
$\mathcal{F}$ が $\bigoplus_{i\in I} \mathcal{O}_X$ と同型になるような
任意の集合 $I$ の濃度が $\text{rank}_\mathcal{F}(x)$ となるものが存在する。
\end{lemma}

\begin{proof}
補題の仮定のもとでは，$I$ の濃度は，零でない環
$\mathcal{O}_{X, x}$ 上の自由加群 $\mathcal{F}_x$ の階数から読み取れ，
$x$ のある近傍上で一定である。
\end{proof}

\begin{lemma}
\label{modules-lemma-map-finite-locally-free}
$(X, \mathcal{O}_X)$ を環付き空間とし，$r \geq 0$ とする。
$\varphi : \mathcal{F} \to \mathcal{G}$ を，階数 $r$ の有限局所自由
$\mathcal{O}_X$-加群の層の写像とする。このとき，$\varphi$ が同型であることと
$\varphi$ が全射であることは同値である。
\end{lemma}

\begin{proof}
$\varphi$ が全射であると仮定し，$x \in X$ をとる。
$\mathcal{F}|_U$ と $\mathcal{G}|_U$ がともに
$\mathcal{O}_U^{\oplus r}$ と同型になるような $x$ の開近傍 $U$ が存在する。
$\mathcal{G}|_U$ の自由生成元の持ち上げを選ぶと，
$\varphi|_U \circ \psi = \text{id}$ を満たす写像
$\psi : \mathcal{G}|_U \to \mathcal{F}|_U$ を得る。したがって，
$\varphi$ が誘導する写像
$\Gamma(U, \mathcal{F}) \to \Gamma(U, \mathcal{G})$
は全射である。$\Gamma(U, \mathcal{F})$ と $\Gamma(U, \mathcal{G})$ は
ともに $\Gamma(U, \mathcal{O}_U)$-加群として
$\Gamma(U, \mathcal{O}_U)^{\oplus r}$ と同型なので，《代数》補題
\ref{algebra-lemma-fun} を適用すると，
$\Gamma(U, \mathcal{F}) \to \Gamma(U, \mathcal{G})$
が単射であることが分かる。これで証明が完了する。
\end{proof}

\begin{lemma}
\label{modules-lemma-direct-summand-of-locally-free-is-locally-free}
$(X, \mathcal{O}_X)$ を環付き空間とする。すべての茎
$\mathcal{O}_{X, x}$ が局所環ならば，有限局所自由
$\mathcal{O}_X$-加群の層の任意の直和因子は有限局所自由である。
\end{lemma}

\begin{proof}
$\mathcal{F}$ が有限局所自由 $\mathcal{O}_X$-加群の層
$\mathcal{H}$ の直和因子であると仮定する。点 $x \in X$ をとる。
このとき $\mathcal{H}_x$ は有限自由 $\mathcal{O}_{X, x}$-加群である。
$\mathcal{O}_{X, x}$ は局所環なので，ある $r$ に対して
$\mathcal{F}_x \cong \mathcal{O}_{X, x}^{\oplus r}$ である。《代数》補題
\ref{algebra-lemma-finite-projective} を参照せよ。補題
\ref{modules-lemma-finite-presentation-stalk-free} により，$\mathcal{F}$ は
$x$ のある開近傍上で階数 $r$ の自由加群の層である。
（$\mathcal{F}$ は $\mathcal{H}$ の直和因子として有限表示であることに注意する。）
\end{proof}





\section{双線形写像}
\label{modules-section-bilinear}


\noindent
$(X, \mathcal{O}_X)$ を環付き空間とし，$\mathcal{F}$，$\mathcal{G}$，
$\mathcal{H}$ を $\mathcal{O}_X$-加群の層とする。
$\mathcal{O}_X$-加群の層の{\it 双線形写像}
$f : \mathcal{F} \times \mathcal{G} \to \mathcal{H}$ とは，表示された
集合の層の写像であって，任意の開集合 $U \subset X$ に対して誘導写像
$$
\mathcal{F}(U) \times \mathcal{G}(U) \to \mathcal{H}(U)
$$
が加群の $\mathcal{O}_X(U)$-双線形写像となるものをいう。同値な定義として，
加群の双線形写像に対する通常の公理を模倣した，集合の層の写像からなる
いくつかの図式が可換であることを要請してもよい。例えば公理
$f(x + y, z) = f(x, z) + f(y, z)$ は，図式
$$
\xymatrix{
\mathcal{F} \times \mathcal{F} \times \mathcal{G}
\ar[rrr]_{(f \circ \text{pr}_{13}, f \circ \text{pr}_{23})}
\ar[d]_{(+ \circ \text{pr}_{12}, \text{pr}_3)} & & &
\mathcal{H} \times \mathcal{H}
\ar[d]^{+} \\
\mathcal{F} \times \mathcal{G} \ar[rrr]^f & & &
\mathcal{H}
}
$$
の可換性によって表される。別の特徴づけは次のとおりである。
$f : \mathcal{F} \times \mathcal{G} \to \mathcal{H}$
が集合の層の写像であり，$X$ のすべての点で茎上に加群の双線形写像を
誘導するならば，$f$ は加群の層の双線形写像である。これは，局所切断が
等しいかどうかを茎で判定できることから従う。

\medskip\noindent
集合の層の圏における射を $\text{Mor}( - , - )$ と書く。
双線形写像のさらに別の特徴づけは次のとおりである。集合の層の写像
$f : \mathcal{F} \times \mathcal{G} \to \mathcal{H}$
が双線形であるための必要十分条件は，任意の集合の層 $\mathcal{S}$ に対して規則
$$
\text{Mor}(\mathcal{S}, \mathcal{F}) \times
\text{Mor}(\mathcal{S}, \mathcal{G}) \to
\text{Mor}(\mathcal{S}, \mathcal{H}),\quad
(a, b) \mapsto f \circ (a \times b)
$$
が環 $\text{Mor}(\mathcal{S}, \mathcal{O}_X)$ 上の加群の双線形写像と
なることである。局所切断の集合で考える方が容易で，両者が同値であることも
明らかなので，通常はこの観点をとらない。

\medskip\noindent
最後に，定義を述べるさらに別の方法を挙げる。
$\mathcal{O}_X$ は集合の層の圏における環対象であり，
$\mathcal{F}$，$\mathcal{G}$，$\mathcal{H}$ はこの環上の加群対象である。
すると，任意の圏において環対象上の加群対象に対して双線形写像を定義できる。
圏における環対象，環対象上の加群対象，およびそれらの双線形写像を定式化するには，
その圏が有限積をもつと仮定すると便利である（ただし厳密には必要でない）。
集合の層の圏は実際に有限積をもつ。







\section{テンソル積}
\label{modules-section-tensor-product}

\noindent
テンソル積については，《層》\ref{sheaves-section-presheaves-modules} 節および
\ref{sheaves-section-sheafification-presheaves-modules} 節で，環の変更という
設定のもとですでに簡単に論じた。これを加群の層のテンソル積へ一般化しよう。

\medskip\noindent
$(X, \mathcal{O}_X)$ を環付き空間とし，$\mathcal{F}$ と $\mathcal{G}$ を
$\mathcal{O}_X$-加群の層とする。まず，{\it テンソル積前層}
$$
\mathcal{F} \otimes_{p, \mathcal{O}_X} \mathcal{G}
$$
を，開集合 $U \subset X$ に $\mathcal{O}_X(U)$-加群
$\mathcal{F}(U) \otimes_{\mathcal{O}_X(U)} \mathcal{G}(U)$ を対応させる
規則として定義する。次に，{\it テンソル積層}を上の前層の層化
$$
\mathcal{F} \otimes_{\mathcal{O}_X} \mathcal{G}
=
(\mathcal{F} \otimes_{p, \mathcal{O}_X} \mathcal{G})^\#
$$
として定義する。これは，任意の第3の $\mathcal{O}_X$-加群の層
$\mathcal{H}$ に対して
$$
\Hom_{\mathcal{O}_X}
(\mathcal{F} \otimes_{\mathcal{O}_X} \mathcal{G}, \mathcal{H})
=
\text{Bilin}_{\mathcal{O}_X}(\mathcal{F} \times \mathcal{G}, \mathcal{H}).
$$
を満たす $\mathcal{O}_X$-加群の層として特徴づけられる。ここで右辺は，
\ref{modules-section-bilinear} 節で定義した $\mathcal{O}_X$-加群の層の双線形写像の
集合を表す。

\medskip\noindent
環 $R$ 上の加群 $M,N$ のテンソル積は対称性，すなわち
$M \otimes_R N = N \otimes_R M$ を満たすので，加群の層のテンソル積にも
同じことが成り立つ。すなわち
$$
\mathcal{F} \otimes_{\mathcal{O}_X} \mathcal{G}
=
\mathcal{G} \otimes_{\mathcal{O}_X} \mathcal{F}
$$
という $\mathcal{F},\mathcal{G}$ に関手的な同型をもつ。また，加群の
テンソル積は結合律を満たすので，標準的な関手的同型
$$
(\mathcal{F} \otimes_{\mathcal{O}_X} \mathcal{G})
\otimes_{\mathcal{O}_X} \mathcal{H}
=
\mathcal{F} \otimes_{\mathcal{O}_X}
(\mathcal{G} \otimes_{\mathcal{O}_X} \mathcal{H})
$$
も得られ，これは $\mathcal{F},\mathcal{G},\mathcal{H}$ に関手的である。

\begin{lemma}
\label{modules-lemma-stalk-tensor-product}
$(X, \mathcal{O}_X)$ を環付き空間，$\mathcal{F},\mathcal{G}$ を
$\mathcal{O}_X$-加群の層，$x \in X$ とする。このとき
$\mathcal{O}_{X, x}$-加群の標準的同型
$$
(\mathcal{F} \otimes_{\mathcal{O}_X} \mathcal{G})_x
=
\mathcal{F}_x \otimes_{\mathcal{O}_{X, x}} \mathcal{G}_x
$$
があり，$\mathcal{F}$ と $\mathcal{G}$ に関手的である。
\end{lemma}

\begin{proof}
省略する。
\end{proof}

\begin{lemma}
\label{modules-lemma-tensor-product-sheafification}
$(X, \mathcal{O}_X)$ を環付き空間とする。$\mathcal{F}',\mathcal{G}'$ を
$\mathcal{O}_X$-加群の前層とし，その層化をそれぞれ
$\mathcal{F},\mathcal{G}$ とする。このとき
$\mathcal{F} \otimes_{\mathcal{O}_X} \mathcal{G} =
(\mathcal{F}' \otimes_{p, \mathcal{O}_X} \mathcal{G}')^\#$.
\end{lemma}

\begin{proof}
省略する。
\end{proof}


\begin{lemma}
\label{modules-lemma-tensor-product-exact}
$(X, \mathcal{O}_X)$ を環付き空間，$\mathcal{G}$ を
$\mathcal{O}_X$-加群の層とする。
$\mathcal{F}_1
\to \mathcal{F}_2
\to \mathcal{F}_3
\to 0$
が $\mathcal{O}_X$-加群の層の完全列ならば，誘導される列
$$
\mathcal{F}_1 \otimes_{\mathcal{O}_X} \mathcal{G} \to
\mathcal{F}_2 \otimes_{\mathcal{O}_X} \mathcal{G} \to
\mathcal{F}_3 \otimes_{\mathcal{O}_X} \mathcal{G} \to
0
$$
は完全である。
\end{lemma}

\begin{proof}
完全性は茎で確認できること（補題 \ref{modules-lemma-abelian}），茎の記述
（補題 \ref{modules-lemma-stalk-tensor-product}），および加群のテンソル積に関する
対応する結果（《代数》補題 \ref{algebra-lemma-tensor-product-exact}）から従う。
\end{proof}

\begin{lemma}
\label{modules-lemma-tensor-product-pullback}
$f : (X, \mathcal{O}_X) \to (Y, \mathcal{O}_Y)$ を環付き空間の射とし，
$\mathcal{F},\mathcal{G}$ を $\mathcal{O}_Y$-加群の層とする。このとき
$f^*(\mathcal{F} \otimes_{\mathcal{O}_Y} \mathcal{G})
= f^*\mathcal{F} \otimes_{\mathcal{O}_X} f^*\mathcal{G}$
が $\mathcal{F},\mathcal{G}$ に関手的に成り立つ。
\end{lemma}

\begin{proof}
省略する。
\end{proof}

\begin{lemma}
\label{modules-lemma-tensor-commute-colimits}
$(X, \mathcal{O}_X)$ を環付き空間とする。任意の
$\mathcal{O}_X$-加群の層 $\mathcal{F}$ に対し，関手
$$
\textit{Mod}(\mathcal{O}_X) \longrightarrow \textit{Mod}(\mathcal{O}_X)
, \quad
\mathcal{G} \longmapsto \mathcal{F} \otimes_{\mathcal{O}_X} \mathcal{G}
$$
は任意の余極限と可換である。
\end{lemma}

\begin{proof}
$I$ を前順序集合，$\{\mathcal{G}_i\}$ を $I$ 上の系とし，
$\mathcal{G} = \colim_i \mathcal{G}_i$ とおく。$\mathcal{G}$ は前層
$\mathcal{G}' : U \mapsto \colim_i \mathcal{G}_i(U)$ に付随する層であることを
思い出そう（《層》\ref{sheaves-section-limits-sheaves} 節）。補題
\ref{modules-lemma-tensor-product-sheafification} により，テンソル積
$\mathcal{F} \otimes_{\mathcal{O}_X} \mathcal{G}$ は前層
$$
U \longmapsto
\mathcal{F}(U) \otimes_{\mathcal{O}_X(U)} \colim_i \mathcal{G}_i(U) =
\colim_i \mathcal{F}(U) \otimes_{\mathcal{O}_X(U)} \mathcal{G}_i(U)
$$
の層化である。ここで等号は《代数》補題
\ref{algebra-lemma-tensor-products-commute-with-limits} による。したがって，
補題 \ref{modules-lemma-limits-colimits} の $\textit{Mod}(\mathcal{O}_X)$ における
余極限の記述から主張が従う。
\end{proof}

\begin{lemma}
\label{modules-lemma-tensor-product-permanence}
$(X, \mathcal{O}_X)$ を環付き空間とし，$\mathcal{F},\mathcal{G}$ を
$\mathcal{O}_X$-加群の層とする。
\begin{enumerate}
\item $\mathcal{F},\mathcal{G}$ が局所的に切断で生成されるならば，
$\mathcal{F} \otimes_{\mathcal{O}_X} \mathcal{G}$ もそうである。
\item $\mathcal{F},\mathcal{G}$ が有限型ならば，
$\mathcal{F} \otimes_{\mathcal{O}_X} \mathcal{G}$ も有限型である。
\item $\mathcal{F},\mathcal{G}$ が準連接ならば，
$\mathcal{F} \otimes_{\mathcal{O}_X} \mathcal{G}$ も準連接である。
\item $\mathcal{F},\mathcal{G}$ が有限表示ならば，
$\mathcal{F} \otimes_{\mathcal{O}_X} \mathcal{G}$ も有限表示である。
\item $\mathcal{F}$ が有限表示で $\mathcal{G}$ が連接ならば，
$\mathcal{F} \otimes_{\mathcal{O}_X} \mathcal{G}$ は連接である。
\item $\mathcal{F},\mathcal{G}$ が連接ならば，
$\mathcal{F} \otimes_{\mathcal{O}_X} \mathcal{G}$ も連接である。
\item $\mathcal{F},\mathcal{G}$ が局所自由ならば，
$\mathcal{F} \otimes_{\mathcal{O}_X} \mathcal{G}$ も局所自由である。
\end{enumerate}
\end{lemma}

\begin{proof}
まず，局所自由 $\mathcal{O}_X$-加群の層のテンソル積が局所自由であることを
示す。そのためには
$(\bigoplus_{i \in I} \mathcal{O}_X) \otimes_{\mathcal{O}_X}
(\bigoplus_{j \in J} \mathcal{O}_X) \cong
\bigoplus_{(i, j) \in I \times J} \mathcal{O}_X$
を示せばよい。層 $\bigoplus_{i \in I} \mathcal{O}_X$ は前層
$U \mapsto \bigoplus_{i \in I} \mathcal{O}_X(U)$ に付随する層である。
したがって，このテンソル積は前層
$$
U \longmapsto
(\bigoplus\nolimits_{i \in I} \mathcal{O}_X(U))
\otimes_{\mathcal{O}_X(U)}
(\bigoplus\nolimits_{j \in J} \mathcal{O}_X(U)).
$$
に付随する層である。任意の環 $R$ に対して
$(\bigoplus_{i \in I} R) \otimes_R (\bigoplus_{j \in J} R) =
\bigoplus_{(i, j) \in I \times J} R$
が成り立つので，所望の結論を得る。

\medskip\noindent
$\mathcal{F}_2 \to \mathcal{F}_1 \to \mathcal{F} \to 0$ が完全ならば，
補題 \ref{modules-lemma-tensor-product-exact} により複体
$\mathcal{F}_2 \otimes_{\mathcal{O}_X} \mathcal{G} \to
\mathcal{F}_1 \otimes_{\mathcal{O}_X} \mathcal{G} \to
\mathcal{F} \otimes_{\mathcal{O}_X} \mathcal{G} \to 0$
は完全である。これを用いて (5) を示せる。実際，この場合には局所的に，
$\mathcal{F}_i$（$i = 1, 2$）が有限自由であるような上の形の完全列が存在する。
したがって二つの項
$\mathcal{F}_2 \otimes_{\mathcal{O}_X} \mathcal{G}$ と
$\mathcal{F}_1 \otimes_{\mathcal{O}_X} \mathcal{G}$
は $\mathcal{G}$ の有限直和に同型である（例えば補題
\ref{modules-lemma-tensor-commute-colimits} による）。有限直和は連接層なので，
これら二つは連接であり，写像の余核も連接である（補題
\ref{modules-lemma-coherent-abelian}）。

\medskip\noindent
$\mathcal{G}_2 \to \mathcal{G}_1 \to \mathcal{G} \to 0$ も完全ならば，
$$
\mathcal{F}_2 \otimes_{\mathcal{O}_X} \mathcal{G}_1
\oplus
\mathcal{F}_1 \otimes_{\mathcal{O}_X} \mathcal{G}_2
\to
\mathcal{F}_1 \otimes_{\mathcal{O}_X} \mathcal{G}_1
\to
\mathcal{F} \otimes_{\mathcal{O}_X} \mathcal{G}
\to 0
$$
が完全であることが分かる。これを用いて，例えば (3) を示せる。実際，仮定は
局所的に，$\mathcal{F}_i$ と $\mathcal{G}_i$ が自由
$\mathcal{O}_X$-加群の層となる上の形の表示を取れることを意味する。
したがって，表示された列はテンソル積の自由層による表示でもある。

\medskip\noindent
他の主張の証明は省略する。
\end{proof}





\section{平坦加群の層}
\label{modules-section-flat}

\noindent
環上の加群の場合とまったく同様に，平坦加群の層を定義できる。

\begin{definition}
\label{modules-definition-flat}
$(X, \mathcal{O}_X)$ を環付き空間とする。関手
$$
\textit{Mod}(\mathcal{O}_X)
\longrightarrow
\textit{Mod}(\mathcal{O}_X), \quad
\mathcal{G} \mapsto \mathcal{G} \otimes_{\mathcal{O}_X} \mathcal{F}
$$
が完全であるとき，$\mathcal{O}_X$-加群の層 $\mathcal{F}$ は
{\it 平坦}であるという。
\end{definition}

\noindent
平坦性は茎を調べることで特徴づけられる。

\begin{lemma}
\label{modules-lemma-flat-stalks-flat}
$(X, \mathcal{O}_X)$ を環付き空間とする。$\mathcal{O}_X$-加群の層
$\mathcal{F}$ が平坦であるための必要十分条件は，すべての $x \in X$ に対し
茎 $\mathcal{F}_x$ が平坦 $\mathcal{O}_{X, x}$-加群であることである。
\end{lemma}

\begin{proof}
すべての $x \in X$ に対し，$\mathcal{F}_x$ が平坦
$\mathcal{O}_{X, x}$-加群であると仮定する。このとき
$\mathcal{G} \to \mathcal{H} \to \mathcal{K}$ が完全ならば，
$\mathcal{G} \otimes_{\mathcal{O}_X} \mathcal{F} \to
\mathcal{H} \otimes_{\mathcal{O}_X} \mathcal{F} \to
\mathcal{K} \otimes_{\mathcal{O}_X} \mathcal{F}$
も完全である。実際，完全性は茎で確認でき，テンソル積は茎を取ることと
可換である（補題 \ref{modules-lemma-stalk-tensor-product}）。逆に，
$\mathcal{F}$ が平坦であると仮定し，$x \in X$ とする。
$M$ を $\mathcal{O}_{X, x}$-加群として，摩天楼層 $i_{x, *}M$ を考える。
このとき，再び補題 \ref{modules-lemma-stalk-tensor-product} により
$$
M \otimes_{\mathcal{O}_{X, x}} \mathcal{F}_x =
\left(i_{x, *} M \otimes_{\mathcal{O}_X} \mathcal{F}\right)_x
$$
である。$i_{x, *}$ は完全なので，$\mathcal{F}$ の平坦性から
$M \mapsto M \otimes_{\mathcal{O}_{X, x}} \mathcal{F}_x$ が完全であることが
分かる。したがって $\mathcal{F}_x$ は平坦 $\mathcal{O}_{X, x}$-加群である。
\end{proof}

\noindent
したがって，次の定義は意味をもつ。

\begin{definition}
\label{modules-definition-flat-at-point}
$(X, \mathcal{O}_X)$ を環付き空間，$x \in X$ とする。
$\mathcal{F}_x$ が平坦 $\mathcal{O}_{X, x}$-加群であるとき，
$\mathcal{O}_X$-加群の層 $\mathcal{F}$ は {\it $x$ で平坦}であるという。
\end{definition}

\noindent
したがって，$\mathcal{F}$ が平坦 $\mathcal{O}_X$-加群の層であるための
必要十分条件は，すべての点で平坦であることである。

\begin{lemma}
\label{modules-lemma-base-change-flat}
$f : (X, \mathcal{O}_X) \to (Y, \mathcal{O}_Y)$ を環付き空間の射とする。
$\mathcal{G}$ が平坦 $\mathcal{O}_Y$-加群の層ならば，$f^*\mathcal{G}$ は
平坦 $\mathcal{O}_X$-加群の層である。
\end{lemma}

\begin{proof}
補題 \ref{modules-lemma-flat-stalks-flat}，《層》補題
\ref{sheaves-lemma-stalk-pullback-modules}，および《代数》補題
\ref{algebra-lemma-flat-base-change} を組み合わせればよい。
\end{proof}

\begin{lemma}
\label{modules-lemma-colimits-flat}
$(X, \mathcal{O}_X)$ を環付き空間とする。平坦 $\mathcal{O}_X$-加群の層の
フィルター余極限は平坦である。また，平坦 $\mathcal{O}_X$-加群の層の直和は
平坦である。
\end{lemma}

\begin{proof}
補題 \ref{modules-lemma-tensor-commute-colimits}，補題
\ref{modules-lemma-stalk-tensor-product}，《代数》補題
\ref{algebra-lemma-directed-colimit-exact}，および完全性を茎で確認できることから
従う。
\end{proof}

\begin{lemma}
\label{modules-lemma-j-shriek-flat}
$(X, \mathcal{O}_X)$ を環付き空間，$U \subset X$ を開集合とする。
層 $j_{U!}\mathcal{O}_U$ は平坦 $\mathcal{O}_X$-加群の層である。
\end{lemma}

\begin{proof}
$j_{U!}\mathcal{O}_U$ の茎は，零であるか $\mathcal{O}_{X, x}$ に等しい。
補題 \ref{modules-lemma-flat-stalks-flat} を適用すればよい。
\end{proof}

\begin{lemma}
\label{modules-lemma-module-quotient-flat}
$(X, \mathcal{O}_X)$ を環付き空間とする。
\begin{enumerate}
\item 任意の $\mathcal{O}_X$-加群の層は，直和
$\bigoplus j_{U_i!}\mathcal{O}_{U_i}$ の商である。
\item 任意の $\mathcal{O}_X$-加群の層は，平坦
$\mathcal{O}_X$-加群の層の商である。
\end{enumerate}
\end{lemma}

\begin{proof}
$\mathcal{F}$ を $\mathcal{O}_X$-加群の層とする。任意の開集合
$U \subset X$ と任意の $s \in \mathcal{F}(U)$ に対し，射
$\mathcal{O}_U \to \mathcal{F}|_U$，$1 \mapsto s$ の随伴として射
$j_{U!}\mathcal{O}_U \to \mathcal{F}$ を得る。明らかに写像
$$
\bigoplus\nolimits_{(U, s)} j_{U!}\mathcal{O}_U
\longrightarrow
\mathcal{F}
$$
は全射であり，補題 \ref{modules-lemma-colimits-flat} と
\ref{modules-lemma-j-shriek-flat} を組み合わせると，その始域は平坦である。
\end{proof}

\begin{lemma}
\label{modules-lemma-flat-tor-zero}
$(X, \mathcal{O}_X)$ を環付き空間とする。
$$
0 \to \mathcal{F}'' \to \mathcal{F}' \to \mathcal{F} \to 0
$$
を $\mathcal{O}_X$-加群の層の短完全列とする。$\mathcal{F}$ が平坦であると
仮定する。このとき，任意の $\mathcal{O}_X$-加群の層 $\mathcal{G}$ に対し，列
$$
0 \to
\mathcal{F}'' \otimes_\mathcal{O} \mathcal{G} \to
\mathcal{F}' \otimes_\mathcal{O} \mathcal{G} \to
\mathcal{F} \otimes_\mathcal{O} \mathcal{G} \to 0
$$
は完全である。
\end{lemma}

\begin{proof}
すべての $x \in X$ に対し $\mathcal{F}_x$ が平坦
$\mathcal{O}_{X, x}$-加群であること，および完全性を茎で確認できることを
用いれば，《代数》補題 \ref{algebra-lemma-flat-tor-zero} から従う。
\end{proof}

\begin{lemma}
\label{modules-lemma-flat-ses}
\begin{slogan}
環付き空間上の平坦加群の層の全射の核および拡大は，再び平坦である。
\end{slogan}
$(X, \mathcal{O}_X)$ を環付き空間とする。
$$
0 \to
\mathcal{F}_2 \to
\mathcal{F}_1 \to
\mathcal{F}_0 \to 0
$$
を $\mathcal{O}_X$-加群の層の短完全列とする。
\begin{enumerate}
\item $\mathcal{F}_2$ と $\mathcal{F}_0$ が平坦ならば，
$\mathcal{F}_1$ も平坦である。
\item $\mathcal{F}_1$ と $\mathcal{F}_0$ が平坦ならば，
$\mathcal{F}_2$ も平坦である。
\end{enumerate}
\end{lemma}

\begin{proof}
完全性と平坦性は茎のレベルで確認できるので，《代数》補題
\ref{algebra-lemma-flat-ses} から従う。
\end{proof}

\begin{lemma}
\label{modules-lemma-flat-resolution-of-flat}
$(X, \mathcal{O}_X)$ を環付き空間とする。
$$
\ldots \to
\mathcal{F}_2 \to
\mathcal{F}_1 \to
\mathcal{F}_0 \to
\mathcal{Q} \to 0
$$
を $\mathcal{O}_X$-加群の層の完全複体とする。$\mathcal{Q}$ およびすべての
$\mathcal{F}_i$ が平坦 $\mathcal{O}_X$-加群の層ならば，任意の
$\mathcal{O}_X$-加群の層 $\mathcal{G}$ に対し，複体
$$
\ldots \to
\mathcal{F}_2 \otimes_{\mathcal{O}_X} \mathcal{G} \to
\mathcal{F}_1 \otimes_{\mathcal{O}_X} \mathcal{G} \to
\mathcal{F}_0 \otimes_{\mathcal{O}_X} \mathcal{G} \to
\mathcal{Q} \otimes_{\mathcal{O}_X} \mathcal{G} \to 0
$$
も完全である。
\end{lemma}

\begin{proof}
補題 \ref{modules-lemma-flat-tor-zero} から従う。実際，この複体を短完全列に分解し，
補題 \ref{modules-lemma-flat-ses} を用いて
$\Im(\mathcal{F}_{i + 1} \to \mathcal{F}_i)$ が平坦であることを帰納的に示す。
\end{proof}

\noindent
次の補題は平坦性の等式判定法（《代数》補題
\ref{algebra-lemma-flat-eq}）の一方向を与える。

\begin{lemma}
\label{modules-lemma-flat-eq}
$(X, \mathcal{O}_X)$ を環付き空間，$\mathcal{F}$ を平坦
$\mathcal{O}_X$-加群の層，$U \subset X$ を開集合とする。
$$
\mathcal{O}_U \xrightarrow{(f_1, \ldots, f_n)}
\mathcal{O}_U^{\oplus n} \xrightarrow{(s_1, \ldots, s_n)}
\mathcal{F}|_U
$$
を $\mathcal{O}_U$-加群の層の複体とする。すべての $x \in U$ に対し，
$x$ の開近傍 $V \subset U$ と，$(s_1, \ldots, s_n)|_V$ の分解
$$
\mathcal{O}_V^{\oplus n}
\xrightarrow{A}
\mathcal{O}_V^{\oplus m} \xrightarrow{(t_1, \ldots, t_m)}
\mathcal{F}|_V
$$
で $A \circ (f_1, \ldots, f_n)|_V = 0$ を満たすものが存在する。
\end{lemma}

\begin{proof}
$\mathcal{I} \subset \mathcal{O}_U$ を $f_1, \ldots, f_n$ で生成される
イデアルの層とする。このとき $\sum f_i \otimes s_i$ は
$\mathcal{I} \otimes_{\mathcal{O}_U} \mathcal{F}|_U$ の切断であり，
$\mathcal{F}|_U$ において零へ写る。$\mathcal{F}|_U$ は平坦なので，写像
$\mathcal{I} \otimes_{\mathcal{O}_U} \mathcal{F}|_U \to \mathcal{F}|_U$
は単射である。$\mathcal{I} \otimes_{\mathcal{O}_U} \mathcal{F}|_U$ は前層の
テンソル積に付随する層なので，$x$ の開近傍 $V \subset U$ で，
$\sum f_i|_V \otimes s_i|_V$ が
$\mathcal{I}(V) \otimes_{\mathcal{O}(V)} \mathcal{F}(V)$ において零となるものが
存在する。《代数》補題 \ref{algebra-lemma-relations} を用いて定義を展開すると，
$t_1, \ldots, t_m \in \mathcal{F}(V)$ と $a_{ij} \in \mathcal{O}(V)$ で
$\sum a_{ij}f_i|_V = 0$ および $s_i|_V = \sum a_{ij}t_j$ を満たすものを得る。
\end{proof}














\section{双対}
\label{modules-section-duals}

\noindent
$(X, \mathcal{O}_X)$ を環付き空間とする。\ref{modules-section-tensor-product} 節で
構成したテンソル積を備えた $\mathcal{O}_X$-加群の層の圏は，対称モノイド圏で
ある。$\mathcal{O}_X$-加群の層 $\mathcal{F}$ に対し，次は同値である。
\begin{enumerate}
\item $\mathcal{O}_X$-加群の層のモノイド圏において，$\mathcal{F}$ は
左双対をもつ。
\item $\mathcal{F}$ は局所的に有限自由 $\mathcal{O}_X$-加群の層の直和因子で
ある。
\item $\mathcal{F}$ は $\mathcal{O}_X$-加群の層として有限表示かつ平坦である。
\end{enumerate}
これは本節の例 \ref{modules-example-dual}，補題 \ref{modules-lemma-left-dual-module}，および
補題 \ref{modules-lemma-flat-locally-finite-presentation} で証明される。

\begin{example}
\label{modules-example-dual}
$(X, \mathcal{O}_X)$ を環付き空間とする。$\mathcal{F}$ を，局所的に有限自由
$\mathcal{O}_X$-加群の層の直和因子である $\mathcal{O}_X$-加群の層とする。
このとき写像
$$
\mathcal{F} \otimes_{\mathcal{O}_X}
\SheafHom_{\mathcal{O}_X}(\mathcal{F}, \mathcal{O}_X)
\longrightarrow
\SheafHom_{\mathcal{O}_X}(\mathcal{F}, \mathcal{F})
$$
は同型である。実際，これは局所的な問題であり，$\mathcal{F}$ が有限自由ならば
成り立ち，またこの性質をもつ加群の層の任意の直和因子についても成り立つ。
上の同型のもとで $\text{id}_\mathcal{F}$ に対応する切断へ $1$ を送る写像を
$$
\eta :
\mathcal{O}_X
\longrightarrow
\mathcal{F} \otimes_{\mathcal{O}_X}
\SheafHom_{\mathcal{O}_X}(\mathcal{F}, \mathcal{O}_X)
$$
$\eta$ と書く。また，評価写像を
$$
\epsilon : 
\SheafHom_{\mathcal{O}_X}(\mathcal{F}, \mathcal{O}_X)
\otimes_{\mathcal{O}_X} \mathcal{F}
\longrightarrow
\mathcal{O}_X
$$
$\epsilon$ と書く。このとき
$\SheafHom_{\mathcal{O}_X}(\mathcal{F}, \mathcal{O}_X), \eta, \epsilon$ は
《圏》定義 \ref{categories-definition-dual} の意味で $\mathcal{F}$ の左双対である。
$(1 \otimes \epsilon) \circ (\eta \otimes 1) = \text{id}_\mathcal{F}$
および
$(\epsilon \otimes 1) \circ (1 \otimes \eta) =
\text{id}_{\SheafHom_{\mathcal{O}_X}(\mathcal{F}, \mathcal{O}_X)}$
の確認は省略する。
\end{example}

\begin{lemma}
\label{modules-lemma-left-dual-module}
$(X, \mathcal{O}_X)$ を環付き空間，$\mathcal{F}$ を
$\mathcal{O}_X$-加群の層とする。$\mathcal{G},\eta,\epsilon$ を，
$\mathcal{O}_X$-加群の層のモノイド圏における $\mathcal{F}$ の左双対とする
（《圏》定義 \ref{categories-definition-dual}）。このとき
\begin{enumerate}
\item $\mathcal{F}$ は局所的に有限自由 $\mathcal{O}_X$-加群の層の直和因子で
ある。
\item 局所切断 $\lambda$ を $(\lambda \otimes 1)(\eta)$ へ送る写像
$e : \SheafHom_{\mathcal{O}_X}(\mathcal{F}, \mathcal{O}_X) \to \mathcal{G}$
は同型である。
\item $\mathcal{F}$ と $\mathcal{G}$ の局所切断 $f$ と $g$ に対して
$\epsilon(f, g) = e^{-1}(g)(f)$ が成り立つ。
\end{enumerate}
\end{lemma}

\begin{proof}
仮定は，
$$
\mathcal{F} \xrightarrow{\eta \otimes 1}
\mathcal{F} \otimes_{\mathcal{O}_X} \mathcal{G}
\otimes_{\mathcal{O}_X} \mathcal{F}
\xrightarrow{1 \otimes \epsilon} \mathcal{F}
\quad\text{および}\quad
\mathcal{G} \xrightarrow{1 \otimes \eta}
\mathcal{G} \otimes_{\mathcal{O}_X} \mathcal{F}
\otimes_{\mathcal{O}_X} \mathcal{G}
\xrightarrow{\epsilon \otimes 1} \mathcal{G}
$$
が恒等写像であることを意味する。$x \in X$ とする。$x$ の開近傍 $U$ と，
$U$ 上の $\mathcal{F}$ および $\mathcal{G}$ の有限個の切断
$f_1, \ldots, f_n$ および $g_1, \ldots, g_n$ で，
$\eta(1) = \sum f_i g_i$ を満たすものを取れる。$i$ 番目の基底ベクトルを
$f_i$ へ送る写像を
$$
\mathcal{O}_U^{\oplus n} \to \mathcal{F}|_U
$$
と書く。このとき写像 $\eta|_U$ は，写像
$\tilde \eta : \mathcal{O}_U \to
\mathcal{O}_U^{\oplus n} \otimes_{\mathcal{O}_U} \mathcal{G}|_U$.
$\tilde\eta$ を経由して分解する。可換図式
$$
\xymatrix{
\mathcal{F}|_U
\ar[rr]_-{\eta \otimes 1} \ar[rrd]_-{\tilde \eta \otimes 1} & &
\mathcal{F}|_U \otimes \mathcal{G}|_U \otimes \mathcal{F}|_U
\ar[r]_-{1 \otimes \epsilon} &
\mathcal{F}|_U \\
& &
\mathcal{O}_U^{\oplus n} \otimes \mathcal{G}|_U \otimes \mathcal{F}|_U
\ar[u] \ar[r]^-{1 \otimes \epsilon} &
\mathcal{O}_U^{\oplus n} \ar[u]
}
$$
を得る。これは，$\mathcal{F}$ 上の恒等写像が $X$ 上局所的に有限自由加群の層を
経由して分解することを示す。これで (1) が証明された。(2) は《圏》補題
\ref{categories-lemma-left-dual} とその証明から従う。(3) は証明中の最初の等式
から従う。また，左双対の一意性（《圏》注意
\ref{categories-remark-left-dual-adjoint}）および例 \ref{modules-example-dual} の
左双対の構成から (2) と (3) を導くこともできる。
\end{proof}

\begin{lemma}
\label{modules-lemma-flat-locally-finite-presentation}
$(X, \mathcal{O}_X)$ を環付き空間とし，$\mathcal{F}$ を有限表示の平坦
$\mathcal{O}_X$-加群の層とする。このとき $\mathcal{F}$ は局所的に有限自由
$\mathcal{O}_X$-加群の層の直和因子である。
\end{lemma}

\begin{proof}
$X$ を開被覆の各要素で置き換えることにより，表示
$$
\mathcal{O}_X^{\oplus r} \to
\mathcal{O}_X^{\oplus n} \to \mathcal{F} \to 0
$$
が存在すると仮定してよい。$x \in X$ とする。補題 \ref{modules-lemma-flat-eq} により，
$X$ を $x$ の開近傍へ縮小した後，分解
$$
\mathcal{O}_X^{\oplus n} \to
\mathcal{O}_X^{\oplus n_1} \to \mathcal{F}
$$
で，合成
$\mathcal{O}_X^{\oplus r} \to \mathcal{O}_X^{\oplus n} \to
\mathcal{O}_X^{\oplus n_1}$
が $\mathcal{O}_X^{\oplus r}$ の第1直和因子を零化するものが存在すると
仮定できる。この議論をさらに $r - 1$ 回繰り返すと，分解
$$
\mathcal{O}_X^{\oplus n} \to
\mathcal{O}_X^{\oplus n_r} \to \mathcal{F}
$$
で，合成
$\mathcal{O}_X^{\oplus r} \to \mathcal{O}_X^{\oplus n}
\to \mathcal{O}_X^{\oplus n_r}$ が零となるものを得る。これは全射
$\mathcal{O}_X^{\oplus n_r} \to \mathcal{F}$ が切断をもつことを意味し，
結論を得る。
\end{proof}







\section{集合の構成可能層}
\label{modules-section-constructible}

\noindent
$X$ を位相空間とする。集合 $S$ に対し，$\underline{S}$ または
$\underline{S}_X$ は値 $S$ をもつ定数層を表すことを思い出そう
（《層》定義 \ref{sheaves-definition-constant-sheaf}）。$U \subset X$ を
位相空間 $X$ の開集合とする。包含射を $j_U$ と書き，《層》
\ref{sheaves-section-open-immersions} 節で記述した空集合による延長を
$j_{U!} : \Sh(U) \to \Sh(X)$ と書く。

\begin{lemma}
\label{modules-lemma-surjection}
$X$ を位相空間，$\mathcal{B}$ を $X$ の位相の基底，$\mathcal{F}$ を
$X$ 上の集合の層とする。集合 $I$，ならびに各 $i \in I$ に対する
$U_i \in \mathcal{B}$ と有限集合 $S_i$ で，全射
$\coprod_{i \in I} j_{U_i!}\underline{S_i} \to \mathcal{F}$ が存在するものを
取れる。
\end{lemma}

\begin{proof}
$S$ を一元集合とする。$S_i = S$ として主張を証明する。任意の $x \in X$ と
$s \in \mathcal{F}_x$ に対し，$\mathcal{F}_x$ で $s$ へ写る
$U(x, s) \in \mathcal{B}$ と $s(x, s) \in \mathcal{F}(U(x, s))$ を選べる。
《層》補題 \ref{sheaves-lemma-j-shriek} により，切断 $s(x,s)$ は層の写像
$j_{U(x, s)!}\underline{S} \to \mathcal{F}$ に対応する。すると
$$
\coprod\nolimits_{(x, s)} j_{U(x, s)!}\underline{S} \to \mathcal{F}
$$
は茎上で全射であり，したがって全射である。
\end{proof}

\begin{lemma}
\label{modules-lemma-filtered-colimit-constructibles}
$X$ を位相空間，$\mathcal{B}$ を $X$ の位相の基底とし，各
$U \in \mathcal{B}$ は準コンパクトであると仮定する。このとき $X$ 上の
任意の集合の層は，次の形の層のフィルター余極限である。
\begin{equation}
\label{modules-equation-towards-constructible-sets}
\text{余等化子}\left(
\xymatrix{
\coprod\nolimits_{b = 1, \ldots, m} j_{V_b!}\underline{S_b}
\ar@<1ex>[r] \ar@<-1ex>[r] &
\coprod\nolimits_{a = 1, \ldots, n} j_{U_a!}\underline{S_a}
}
\right)
\end{equation}
ここで $U_a,V_b$ は $\mathcal{B}$ の元，$S_a,S_b$ は有限集合である。
\end{lemma}

\begin{proof}
補題 \ref{modules-lemma-surjection} により，任意の集合の層 $\mathcal{F}$ は，始域
$\mathcal{F}_0$ が $U \in \mathcal{B}$ と有限集合 $S$ に対する
$j_{U!}\underline{S}$ の形の層の余積であるような全射の終域となる。これを
$\mathcal{F}_0 \times_\mathcal{F} \mathcal{F}_0$ に適用すると，$\mathcal{F}$ は
一対の写像
$$
\xymatrix{
\coprod\nolimits_{b \in B} j_{V_b!}\underline{S_b}
\ar@<1ex>[r] \ar@<-1ex>[r] &
\coprod\nolimits_{a \in A} j_{U_a!}\underline{S_a}
}
$$
の余等化子であることが分かる。ここで $A,B$ は添字集合，$V_b,U_a$ は
$\mathcal{B}$ の元，$S_a,S_b$ は有限集合である。任意の有限部分集合
$B' \subset B$ に対し，有限部分集合 $A' \subset A$ で，$b \in B'$ 上の余積が
二つの写像のいずれによっても $a \in A'$ 上の余積へ写るものが存在する。
実際，右辺を単射な遷移写像をもつフィルター余極限とみなせる。したがって，
準コンパクト開集合 $V_b$（$b \in B'$）上の切断を取ることはこの余積と可換で
ある（《層》補題 \ref{sheaves-lemma-directed-colimits-sections}）。ゆえに本層は，
有限余積間のこれらの写像の余核の余極限である。
\end{proof}

\begin{lemma}
\label{modules-lemma-constructible-comes-from-finite}
$X$ をスペクトル位相空間，$\mathcal{B}$ を $X$ の準コンパクト開部分集合全体と
する。$\mathcal{F}$ を式 (\ref{modules-equation-towards-constructible-sets}) の形の
集合の層とする。このとき有限ソーバー位相空間 $Y$ への連続スペクトル写像
$f : X \to Y$ と，有限な茎をもつ $Y$ 上の集合の層 $\mathcal{G}$ で，
$f^{-1}\mathcal{G} \cong \mathcal{F}$ を満たすものが存在する。
\end{lemma}

\begin{proof}
$X = \lim X_i$ を有限ソーバー空間の有向系の極限として書ける
（《位相》補題
\ref{topology-lemma-spectral-inverse-limit-finite-sober-spaces}）。もちろん遷移写像
$X_{i'} \to X_i$ はスペクトル写像であり，したがって《位相》補題
\ref{topology-lemma-directed-inverse-limit-spectral-spaces} により，写像
$p_i : X \to X_i$ もスペクトル写像である。ある $i$ に対し，逆像がそれぞれ
$U_a,V_b$ となる $X_i$ の開集合 $U_{a,i},V_{b,i}$ を取れる
（《位相》補題 \ref{topology-lemma-descend-opens}）。余等化子が
$\mathcal{F}$ となる二つの写像
$$
\beta, \gamma :
\coprod\nolimits_{b \in B} j_{V_b!}\underline{S_b}
\longrightarrow
\coprod\nolimits_{a \in A} j_{U_a!}\underline{S_a}
$$
は，随伴により集合の写像の二つの族
$$
\beta_b, \gamma_b :
S_b
\longrightarrow
\Gamma(V_b, \coprod\nolimits_{a \in A} j_{U_a!}\underline{S_a}), \quad
b \in B
$$
に対応する。ここで
$p_i^{-1}(j_{U_{a, i}!}\underline{S_a}) = j_{U_a!}\underline{S_a}$
かつ $(X_{i'} \to X_i)^{-1}(j_{U_{a, i}!}\underline{S_a}) =
j_{U_{a, i'}!}\underline{S_a}$ であることに注意する。《層》補題
\ref{sheaves-lemma-compute-pullback-to-limit}（および $S_b$ と $B$ が有限集合で
あること）から，$i$ を大きくした後，写像
$$
\beta_{b, i}, \gamma_{b, i} :
S_b
\longrightarrow
\Gamma(V_{b, i}, \coprod\nolimits_{a \in A} j_{U_{a, i}!}\underline{S_a})
, \quad b \in B
$$
で，$p_i$ により引き戻すと $\beta_b,\gamma_b$ を与えるものを得る。これらは
さらに $X_i$ 上の層の写像
$$
\beta_i, \gamma_i :
\coprod\nolimits_{b \in B} j_{V_{b, i}!}\underline{S_b}
\longrightarrow
\coprod\nolimits_{a \in A} j_{U_{a, i}!}\underline{S_a}
$$
に対応する。このとき $Y = X_i$ および
$$
\mathcal{G} =
\text{余等化子}\left(
\xymatrix{
\coprod\nolimits_{b = 1, \ldots, m} j_{V_{b, i}!}\underline{S_b}
\ar@<1ex>[r] \ar@<-1ex>[r] &
\coprod\nolimits_{a = 1, \ldots, n} j_{U_{a, i}!}\underline{S_a}
}
\right)
$$
と取ればよい。いくつかの詳細は省略する。
\end{proof}

\begin{lemma}
\label{modules-lemma-constructible-in-constant}
$X$ をスペクトル位相空間，$\mathcal{B}$ を $X$ の準コンパクト開部分集合全体と
する。$\mathcal{F}$ を式 (\ref{modules-equation-towards-constructible-sets}) の形の
集合の層とする。このとき有限個の構成可能閉部分集合
$Z_1, \ldots, Z_n \subset X$ と有限集合 $S_i$ で，$\mathcal{F}$ が
$\prod (Z_i \to X)_*\underline{S_i}$ の部分層に同型となるものが存在する。
\end{lemma}

\begin{proof}
補題 \ref{modules-lemma-constructible-comes-from-finite} により，有限ソーバー位相空間と
有限な茎をもつ層の場合に帰着する。この場合
$\mathcal{F} \subset \prod_{x \in X} i_{x, *}\mathcal{F}_x$
である。ここで $i_x : \{x\} \to X$ は埋め込みである。
$i_{x, *}\mathcal{F}_x$ が $\overline{\{x\}}$ 上の定数層であることの証明は
省略する。
\end{proof}










\section{環付き空間の平坦射}
\label{modules-section-flat-morphisms}

\noindent
点ごとの定義は，上の補題 \ref{modules-lemma-flat-stalks-flat} と定義
\ref{modules-definition-flat-at-point} によって動機づけられる。

\begin{definition}
\label{modules-definition-flat-morphism}
$f : X \to Y$ を環付き空間の射，$x \in X$ とする。環の写像
$\mathcal{O}_{Y, f(x)} \to \mathcal{O}_{X, x}$ が平坦であるとき，$f$ は
{\it $x$ で平坦}であるという。すべての $x \in X$ で $f$ が平坦であるとき，
$f$ は {\it 平坦}であるという。
\end{definition}

\noindent
環の層の写像 $f^\sharp : f^{-1}\mathcal{O}_Y \to \mathcal{O}_X$ を考える。
$x$ におけるその茎は環の写像
$f^\sharp_x : \mathcal{O}_{Y, f(x)} \to \mathcal{O}_{X, x}$ である。したがって，
$f$ が $x$ で平坦であるための必要十分条件は，$\mathcal{O}_X$ が
$f^{-1}\mathcal{O}_Y$-加群の層として $x$ で平坦であることである。また，
$f$ が平坦であるための必要十分条件は，$\mathcal{O}_X$ が
$f^{-1}\mathcal{O}_Y$-加群の層として平坦であることである。平坦射の非常に
特別な場合として，開埋め込みがある。

\begin{lemma}
\label{modules-lemma-pullback-flat}
$f : X \to Y$ を環付き空間の平坦射とする。このとき引き戻し関手
$f^* : \textit{Mod}(\mathcal{O}_Y) \to \textit{Mod}(\mathcal{O}_X)$
は完全である。
\end{lemma}

\begin{proof}
関手 $f^*$ は，完全関手
$f^{-1} : \textit{Mod}(\mathcal{O}_Y) \to \textit{Mod}(f^{-1}\mathcal{O}_Y)$
と環の変更の関手
$$
\textit{Mod}(f^{-1}\mathcal{O}_Y) \to \textit{Mod}(\mathcal{O}_X), \quad
\mathcal{F} \longmapsto
\mathcal{F} \otimes_{f^{-1}\mathcal{O}_Y} \mathcal{O}_X.
$$
の合成である。したがって，定義 \ref{modules-definition-flat-morphism} に続く議論から
結果が従う。
\end{proof}

\begin{definition}
\label{modules-definition-flat-module}
$f : (X, \mathcal{O}_X) \to (Y, \mathcal{O}_Y)$ を環付き空間の射とし，
$\mathcal{F}$ を $\mathcal{O}_X$-加群の層とする。
\begin{enumerate}
\item 茎 $\mathcal{F}_x$ が平坦 $\mathcal{O}_{Y, f(x)}$-加群であるとき，
$\mathcal{F}$ は {\it 点 $x \in X$ で $Y$ 上平坦}であるという。
\item $X$ のすべての点 $x$ で $\mathcal{F}$ が $Y$ 上平坦であるとき，
$\mathcal{F}$ は {\it $Y$ 上平坦}であるという。
\end{enumerate}
\end{definition}

\noindent
この定義により，$\mathcal{F}$ が $x$ で $Y$ 上平坦であるための必要十分条件は，
$\mathcal{F}$ が $f^{-1}\mathcal{O}_Y$-加群の層として $x$ で平坦であることで
ある。実際，《層》補題 \ref{sheaves-lemma-stalk-pullback} により
$(f^{-1}\mathcal{O}_Y)_x = \mathcal{O}_{Y, f(x)}$ である。

\begin{lemma}
\label{modules-lemma-pullback-tensor-flat-module}
$f : X \to Y$ を環付き空間の射とし，$\mathcal{F}$ を $Y$ 上平坦な
$\mathcal{O}_X$-加群の層とする。このとき関手
$$
\textit{Mod}(\mathcal{O}_Y) \to \textit{Mod}(\mathcal{O}_X),\quad
\mathcal{G} \longmapsto f^*\mathcal{G} \otimes_{\mathcal{O}_X} \mathcal{F}
$$
は完全である。
\end{lemma}

\begin{proof}
これは，
$f^*\mathcal{G} \otimes_{\mathcal{O}_X} \mathcal{F} =
f^{-1}\mathcal{G} \otimes_{f^{-1}\mathcal{O}_Y} \mathcal{F}$,
関手 $f^{-1}$ が完全であること，および $\mathcal{F}$ が平坦
$f^{-1}\mathcal{O}_Y$-加群の層であることから従う。
\end{proof}














\section{対称冪と外冪}
\label{modules-section-symmetric-exterior}

\noindent
$(X, \mathcal{O}_X)$ を環付き空間，$\mathcal{F}$ を
$\mathcal{O}_X$-加群の層とする。$\mathcal{F}$ の {\it テンソル代数}を，
非可換 $\mathcal{O}_X$-代数の層
$$
\text{T}(\mathcal{F})
=
\text{T}_{\mathcal{O}_X}(\mathcal{F})
= \bigoplus\nolimits_{n \geq 0} \text{T}^n(\mathcal{F}).
$$
として定義する。ここで $\text{T}^0(\mathcal{F}) = \mathcal{O}_X$，
$\text{T}^1(\mathcal{F}) = \mathcal{F}$ であり，$n \geq 2$ に対して
$$
\text{T}^n(\mathcal{F}) =
\mathcal{F} \otimes_{\mathcal{O}_X} \ldots \otimes_{\mathcal{O}_X} \mathcal{F}
\ \ (n\text{ 個の因子})
$$
である。$\wedge(\mathcal{F})$ を，$\mathcal{F}$ の局所切断 $s$ に対する
$\text{T}^2(\mathcal{F})$ の局所切断 $s \otimes s$ で生成される両側イデアルに
よる $\text{T}(\mathcal{F})$ の商として定義する。これを $\mathcal{F}$ の
{\it 外積代数}という。同様に，$\text{Sym}(\mathcal{F})$ を，
$\text{T}^2(\mathcal{F})$ の $s \otimes t - t \otimes s$ の形の局所切断で
生成される両側イデアルによる $\text{T}(\mathcal{F})$ の商として定義する。

\medskip\noindent
代数 $\text{T}(\mathcal{F})$，$\wedge(\mathcal{F})$，および
$\text{Sym}(\mathcal{F})$ は，《微分次数付き層》定義
\ref{sdga-definition-ga} の意味で次数付き $\mathcal{O}_X$-代数である。さらに
$\text{Sym}(\mathcal{F})$ は可換であり，$\wedge(\mathcal{F})$ は次数付き可換で
ある。

\begin{lemma}
\label{modules-lemma-local-tensor-algebra}
上の状況において，層 $\wedge^n\mathcal{F}$ は前層
$$
U \longmapsto \wedge^n_{\mathcal{O}_X(U)}(\mathcal{F}(U)).
$$
の層化である。《代数》\ref{algebra-section-tensor-algebra} 節を参照せよ。
同様に，層 $\text{Sym}^n\mathcal{F}$ は前層
$$
U \longmapsto \text{Sym}^n_{\mathcal{O}_X(U)}(\mathcal{F}(U)).
$$
の層化である。
\end{lemma}

\begin{proof}
省略する。$\text{Sym}(\mathcal{F})$ と $\wedge(\mathcal{F})$ は，上で与えた
方法よりもこの方法で定義する方が効率的かもしれない。
\end{proof}

\begin{lemma}
\label{modules-lemma-stalk-tensor-algebra}
上の状況において $x \in X$ とする。$\mathcal{O}_{X, x}$-加群の標準的同型
$\text{T}(\mathcal{F})_x = \text{T}(\mathcal{F}_x)$,
$\text{Sym}(\mathcal{F})_x = \text{Sym}(\mathcal{F}_x)$，および
$\wedge(\mathcal{F})_x = \wedge(\mathcal{F}_x)$ が存在する。
\end{lemma}

\begin{proof}
上の補題 \ref{modules-lemma-local-tensor-algebra} および《代数》補題
\ref{algebra-lemma-colimit-tensor-algebra} から明らかである。
\end{proof}

\begin{lemma}
\label{modules-lemma-pullback-tensor-algebra}
$f : (X, \mathcal{O}_X) \to (Y, \mathcal{O}_Y)$ を環付き空間の射とし，
$\mathcal{F}$ を $\mathcal{O}_Y$-加群の層とする。このとき
$f^*\text{T}(\mathcal{F}) = \text{T}(f^*\mathcal{F})$ であり，
$\mathcal{F}$ に付随する外積代数と対称代数についても同様である。
\end{lemma}

\begin{proof}
省略する。
\end{proof}

\begin{lemma}
\label{modules-lemma-presentation-sym-exterior}
$(X, \mathcal{O}_X)$ を環付き空間とする。
$\mathcal{F}_2 \to \mathcal{F}_1 \to \mathcal{F} \to 0$ を
$\mathcal{O}_X$-加群の層の完全列とする。各 $n \geq 1$ に対して完全列
$$
\mathcal{F}_2 \otimes_{\mathcal{O}_X} \text{Sym}^{n - 1}(\mathcal{F}_1)
\to
\text{Sym}^n(\mathcal{F}_1)
\to
\text{Sym}^n(\mathcal{F})
\to
0
$$
および同様の完全列
$$
\mathcal{F}_2 \otimes_{\mathcal{O}_X} \wedge^{n - 1}(\mathcal{F}_1)
\to
\wedge^n(\mathcal{F}_1)
\to
\wedge^n(\mathcal{F})
\to
0
$$
が存在する。
\end{lemma}

\begin{proof}
《代数》補題 \ref{algebra-lemma-presentation-sym-exterior} を参照せよ。
\end{proof}

\begin{lemma}
\label{modules-lemma-tensor-algebra-permanence}
$(X, \mathcal{O}_X)$ を環付き空間とし，$\mathcal{F}$ を
$\mathcal{O}_X$-加群の層とする。
\begin{enumerate}
\item $\mathcal{F}$ が局所的に切断で生成されるならば，各
$\text{T}^n(\mathcal{F})$，$\wedge^n(\mathcal{F})$，および
$\text{Sym}^n(\mathcal{F})$ もそうである。
\item $\mathcal{F}$ が有限型ならば，各 $\text{T}^n(\mathcal{F})$，
$\wedge^n(\mathcal{F})$，および $\text{Sym}^n(\mathcal{F})$ も有限型である。
\item $\mathcal{F}$ が有限表示ならば，各 $\text{T}^n(\mathcal{F})$，
$\wedge^n(\mathcal{F})$，および $\text{Sym}^n(\mathcal{F})$ も有限表示である。
\item $\mathcal{F}$ が連接ならば，$n > 0$ に対する各
$\text{T}^n(\mathcal{F})$，$\wedge^n(\mathcal{F})$，および
$\text{Sym}^n(\mathcal{F})$ は連接である。
\item $\mathcal{F}$ が準連接ならば，各 $\text{T}^n(\mathcal{F})$，
$\wedge^n(\mathcal{F})$，および $\text{Sym}^n(\mathcal{F})$ も準連接である。
\item $\mathcal{F}$ が局所自由ならば，各 $\text{T}^n(\mathcal{F})$，
$\wedge^n(\mathcal{F})$，および $\text{Sym}^n(\mathcal{F})$ も局所自由である。
\end{enumerate}
\end{lemma}

\begin{proof}
$\text{T}^n(\mathcal{F})$ に関するこれらの主張は，補題
\ref{modules-lemma-tensor-product-permanence} から従う。

\medskip\noindent
(1) と (2) は，$\wedge^n(\mathcal{F})$ と $\text{Sym}^n(\mathcal{F})$ が
$\text{T}^n(\mathcal{F})$ の商であることから従う。

\medskip\noindent
(6) は《代数》補題 \ref{algebra-lemma-free-tensor-algebra} から従う。

\medskip\noindent
(3) と (5) には，上の補題 \ref{modules-lemma-presentation-sym-exterior} を用いる。
局所的に，$\mathcal{F}_i$ が自由または有限自由となる表示
$\mathcal{F}_2 \to \mathcal{F}_1 \to \mathcal{F} \to 0$
を選んで補題を適用すると，$\text{Sym}^n(\mathcal{F})$ と
$\wedge^n(\mathcal{F})$ も同様の表示をもつことが分かる。ここでは (6) と
補題 \ref{modules-lemma-tensor-product-permanence} を用いる。

\medskip\noindent
(4) の証明には《代数》補題 \ref{algebra-lemma-present-sym-wedge} を用いる。
$X$ 上で局所化し，$\mathcal{F}$ が大域切断の有限集合
$(s_i)_{i \in I}$ で生成されると仮定してよい。上で挙げた補題と上の補題
\ref{modules-lemma-local-tensor-algebra} を組み合わせると，$n \geq 2$ に対して完全列
$$
\bigoplus\nolimits_{j \in J}
\text{T}^{n - 2}(\mathcal{F})
\to
\text{T}^n(\mathcal{F})
\to
\text{Sym}^n(\mathcal{F})
\to
0
$$
が存在し，添字集合 $J$ は有限である。ここで
$\text{T}^{n - 2}(\mathcal{F})$ は有限生成であることが分かっているので，
最初の矢印の像は $\text{T}^n(\mathcal{F})$ の連接部分層である
（補題 \ref{modules-lemma-coherent-abelian}）。同じ補題により，
$\text{Sym}^n(\mathcal{F})$ は連接であると結論する。
\end{proof}

\begin{lemma}
\label{modules-lemma-whole-tensor-algebra-permanence}
$(X, \mathcal{O}_X)$ を環付き空間とし，$\mathcal{F}$ を
$\mathcal{O}_X$-加群の層とする。
\begin{enumerate}
\item $\mathcal{F}$ が準連接ならば，$\text{T}(\mathcal{F})$，
$\wedge(\mathcal{F})$，および $\text{Sym}(\mathcal{F})$ も準連接である。
\item $\mathcal{F}$ が局所自由ならば，$\text{T}(\mathcal{F})$，
$\wedge(\mathcal{F})$，および $\text{Sym}(\mathcal{F})$ も局所自由である。
\end{enumerate}
\end{lemma}

\begin{proof}
局所自由加群の層の無限直和 $\bigoplus \mathcal{G}_i$ が局所自由であるとは
限らず，準連接加群の層の無限直和が準連接であるとも限らない。問題は，
点 $x \in X$ を与えたとき，$\mathcal{G}_i$ が自由となる（それぞれ，適切な
表示をもつ）$x$ の開近傍 $U_i$ の共通部分が $x$ の開近傍であるとは限らない
ことである。しかし，補題 \ref{modules-lemma-tensor-algebra-permanence} の証明では，
$\mathcal{F}$ に対する適切な開近傍を一度選べば，この開近傍が層
$\text{T}^n(\mathcal{F})$，$\wedge^n(\mathcal{F})$，および
$\text{Sym}^n(\mathcal{F})$ のそれぞれに対して機能することを見た。
これで補題が従う。
\end{proof}







\section{内部 Hom}
\label{modules-section-internal-hom}

\noindent
$(X, \mathcal{O}_X)$ を環付き空間とし，$\mathcal{F},\mathcal{G}$ を
$\mathcal{O}_X$-加群の層とする。規則
$$
U \longmapsto \Hom_{\mathcal{O}_X|_U}(\mathcal{F}|_U, \mathcal{G}|_U).
$$
を考える。《層》\ref{sheaves-section-glueing-sheaves} 節の議論から，これは
アーベル群の層である。さらに，
$\varphi \in \Hom_{\mathcal{O}_X|_U}(\mathcal{F}|_U, \mathcal{G}|_U)$ と
$f \in \mathcal{O}_X(U)$ が与えられたとき，$\mathcal{F}|_U$ 上の $f$ 倍を
前合成するか，$\mathcal{G}|_U$ 上の $f$ 倍を後合成することにより
$f\varphi \in \Hom_{\mathcal{O}_X|_U}(\mathcal{F}|_U, \mathcal{G}|_U)$ を
定義できる（どちらも同じ結果を与える）。したがって，実際には
$\mathcal{O}_X$-加群の層を得る。この層を
$\SheafHom_{\mathcal{O}_X}(\mathcal{F}, \mathcal{G})$ と書く。
標準的な「評価」射
$$
\mathcal{F}
\otimes_{\mathcal{O}_X}
\SheafHom_{\mathcal{O}_X}(\mathcal{F}, \mathcal{G})
\longrightarrow
\mathcal{G}.
$$
がある。また，すべての $x \in X$ に対して標準的射
$$
\SheafHom_{\mathcal{O}_X}(\mathcal{F}, \mathcal{G})_x
\to
\Hom_{\mathcal{O}_{X, x}}(\mathcal{F}_x, \mathcal{G}_x)
$$
があるが，これが同型であることは稀である。

\begin{lemma}
\label{modules-lemma-internal-hom}
$(X, \mathcal{O}_X)$ を環付き空間とし，$\mathcal{F},\mathcal{G},\mathcal{H}$
を $\mathcal{O}_X$-加群の層とする。標準的同型
$$
\SheafHom_{\mathcal{O}_X}
(\mathcal{F} \otimes_{\mathcal{O}_X} \mathcal{G}, \mathcal{H})
\longrightarrow
\SheafHom_{\mathcal{O}_X}
(\mathcal{F}, \SheafHom_{\mathcal{O}_X}(\mathcal{G}, \mathcal{H}))
$$
があり，三つの変数すべてに関手的である（三つの位置すべてで層 Hom を取る）。
特に，射 $\mathcal{F} \otimes_{\mathcal{O}_X} \mathcal{G} \to \mathcal{H}$ を
与えることは，射
$\mathcal{F} \to \SheafHom_{\mathcal{O}_X}(\mathcal{G}, \mathcal{H})$ を
与えることと同じである。
\end{lemma}

\begin{proof}
これは《代数》補題 \ref{algebra-lemma-hom-from-tensor-product} の類似である。
証明も同じなので省略する。
\end{proof}

\begin{lemma}
\label{modules-lemma-internal-hom-exact}
$(X, \mathcal{O}_X)$ を環付き空間とし，$\mathcal{F},\mathcal{G}$ を
$\mathcal{O}_X$-加群の層とする。
\begin{enumerate}
\item $\mathcal{F}_2 \to \mathcal{F}_1 \to \mathcal{F} \to 0$ が
$\mathcal{O}_X$-加群の層の完全列ならば，
$$
0 \to
\SheafHom_{\mathcal{O}_X}(\mathcal{F}, \mathcal{G}) \to
\SheafHom_{\mathcal{O}_X}(\mathcal{F}_1, \mathcal{G}) \to
\SheafHom_{\mathcal{O}_X}(\mathcal{F}_2, \mathcal{G})
$$
は完全である。
\item $0 \to \mathcal{G} \to \mathcal{G}_1 \to \mathcal{G}_2$ が
$\mathcal{O}_X$-加群の層の完全列ならば，
$$
0 \to
\SheafHom_{\mathcal{O}_X}(\mathcal{F}, \mathcal{G}) \to
\SheafHom_{\mathcal{O}_X}(\mathcal{F}, \mathcal{G}_1) \to
\SheafHom_{\mathcal{O}_X}(\mathcal{F}, \mathcal{G}_2)
$$
は完全である。
\end{enumerate}
\end{lemma}

\begin{proof}
(1) のような $\mathcal{F}_2 \to \mathcal{F}_1 \to \mathcal{F} \to 0$ を
取る。任意の開集合 $U \subset X$ に対し，列
$$
0 \to \Hom_{\mathcal{O}_U}(\mathcal{F}|_U, \mathcal{G}|_U) \to
\Hom_{\mathcal{O}_U}(\mathcal{F}_1|_U, \mathcal{G}|_U) \to
\Hom_{\mathcal{O}_U}(\mathcal{F}_2|_U, \mathcal{G}|_U)
$$
は《ホモロジー》補題 \ref{homology-lemma-check-exactness} により完全である。
これは，(1) の層の列の $U$ 上の切断を取るとアーベル群の完全列が得られることを
意味する。したがって定義により (1) の列は完全である。(2) の証明もまったく
同じである。
\end{proof}

\begin{lemma}
\label{modules-lemma-adjoint-tensor-restrict}
$X$ を位相空間とし，$\mathcal{O}_1 \to \mathcal{O}_2$ を環の層の準同型とする。
このとき
$$
\Hom_{\mathcal{O}_1}(\mathcal{F}_{\mathcal{O}_1}, \mathcal{G}) =
\Hom_{\mathcal{O}_2}(\mathcal{F},
\SheafHom_{\mathcal{O}_1}(\mathcal{O}_2, \mathcal{G}))
$$
が $\mathcal{F} \in \textit{Mod}(\mathcal{O}_2)$ と
$\mathcal{G} \in \textit{Mod}(\mathcal{O}_1)$ に双関手的に成り立つ。
\end{lemma}

\begin{proof}
省略する。これは《代数》補題 \ref{algebra-lemma-adjoint-hom-restrict} の
類似であり，まったく同じ方法で証明される。
\end{proof}

\begin{lemma}
\label{modules-lemma-stalk-internal-hom}
$(X, \mathcal{O}_X)$ を環付き空間とし，$\mathcal{F},\mathcal{G}$ を
$\mathcal{O}_X$-加群の層とする。$\mathcal{F}$ が有限型ならば，標準写像
$$
\SheafHom_{\mathcal{O}_X}(\mathcal{F}, \mathcal{G})_x
\to
\Hom_{\mathcal{O}_{X, x}}(\mathcal{F}_x, \mathcal{G}_x)
$$
は単射である。$\mathcal{F}$ が有限表示ならば，この標準射は同型である。
\end{lemma}

\begin{proof}
この写像は，$\SheafHom_{\mathcal{O}_X}(\mathcal{F}, \mathcal{G})_x$ における
$(U,\varphi)$ の同値類を，$x$ における茎上の誘導写像
$\varphi_x : \mathcal{F}_x \to \mathcal{G}_x$ へ送る。ここで
$x \in U \subset X$ は開であり，
$\varphi \in \Hom_{\mathcal{O}_U}(\mathcal{F}|_U, \mathcal{G}|_U)$ である。

\medskip\noindent
$\mathcal{F}$ が有限型であると仮定する。写像の核の元 $\sigma$，すなわち
$\varphi_x = 0$ となるものの代表 $(U,\varphi)$ を取る。必要なら $U$ を
縮小し，$\mathcal{F}|_U$ を生成する切断
$s^1, \ldots, s^n \in \mathcal{F}(U)$ を選ぶ。$\varphi_x(s^i_x)=0$ であり，
切断は有限個なので，$x$ の開近傍 $V \subset U$ で，すべての
$i = 1, \ldots, n$ に対し $\varphi_V(s^i|_V)=0$ となるものを取れる。
$s^i|_V$（$i = 1, \ldots, n$）は $\mathcal{F}|_V$ を生成するので，これは
$\varphi|_V=0$ を意味する。$(U,\varphi)$ は $(V,\varphi|_V)$ と同値なので，
$\sigma=0$ と結論でき，写像の単射性が従う。

\medskip\noindent
次に，$\mathcal{F}$ が有限表示であると仮定する。$X$ 上で局所化することにより，
$\mathcal{F}$ が表示
$$
\bigoplus\nolimits_{j = 1, \ldots, m}
\mathcal{O}_X
\longrightarrow
\bigoplus\nolimits_{i = 1, \ldots, n}
\mathcal{O}_X
\to
\mathcal{F}
\to
0.
$$
をもつと仮定してよい。補題 \ref{modules-lemma-internal-hom-exact} により，これは完全列
$
0 \to
\SheafHom_{\mathcal{O}_X}(\mathcal{F}, \mathcal{G}) \to
\bigoplus\nolimits_{i = 1, \ldots, n} \mathcal{G}
\longrightarrow
\bigoplus\nolimits_{j = 1, \ldots, m} \mathcal{G}.
$
を与える。茎を取ると完全列
$
0 \to
\SheafHom_{\mathcal{O}_X}(\mathcal{F}, \mathcal{G})_x \to
\bigoplus\nolimits_{i = 1, \ldots, n} \mathcal{G}_x
\longrightarrow
\bigoplus\nolimits_{j = 1, \ldots, m} \mathcal{G}_x
$
を得る。一方，$\mathcal{F}_x$ は完全列
$
\bigoplus\nolimits_{j = 1, \ldots, m}
\mathcal{O}_{X, x}
\longrightarrow
\bigoplus\nolimits_{i = 1, \ldots, n}
\mathcal{O}_{X, x}
\to
\mathcal{F}_x
\to
0
$
に入り，これは完全列
$
0 \to
\Hom_{\mathcal{O}_{X, x}}(\mathcal{F}_x, \mathcal{G}_x) \to
\bigoplus\nolimits_{i = 1, \ldots, n} \mathcal{G}_x
\longrightarrow
\bigoplus\nolimits_{j = 1, \ldots, m} \mathcal{G}_x
$
を誘導する。これは上の列と同じなので，結果が従う。
\end{proof}

\begin{lemma}
\label{modules-lemma-pullback-internal-hom}
$f : (X, \mathcal{O}_X) \to (Y, \mathcal{O}_Y)$ を環付き空間の射とし，
$\mathcal{F},\mathcal{G}$ を $\mathcal{O}_Y$-加群の層とする。
$\mathcal{F}$ が有限表示で $f$ が平坦ならば，標準写像
$$
f^*\SheafHom_{\mathcal{O}_Y}(\mathcal{F}, \mathcal{G})
\longrightarrow
\SheafHom_{\mathcal{O}_X}(f^*\mathcal{F}, f^*\mathcal{G})
$$
は同型である。
\end{lemma}

\begin{proof}
$f^*\mathcal{F}$ も有限表示であることに注意する
（補題 \ref{modules-lemma-pullback-finite-presentation}）。$x \in X$ が
$y \in Y$ へ写るとする。$x$ における茎を見ると，補題
\ref{modules-lemma-stalk-internal-hom} および《代数詳論》補題
\ref{more-algebra-lemma-pseudo-coherence-and-base-change-ext} により同型を得る。
実際，この場合 $\Hom$ は
$\mathcal{O}_{Y, y} \to \mathcal{O}_{X, x}$ による基底変換と可換である。
第2の証明：補題 \ref{modules-lemma-stalk-internal-hom} の証明とまったく同じ議論を
用いる。
\end{proof}

\begin{lemma}
\label{modules-lemma-internal-hom-locally-kernel-direct-sum}
$(X, \mathcal{O}_X)$ を環付き空間とし，$\mathcal{F},\mathcal{G}$ を
$\mathcal{O}_X$-加群の層とする。$\mathcal{F}$ が有限表示ならば，層
$\SheafHom_{\mathcal{O}_X}(\mathcal{F}, \mathcal{G})$ は局所的に，
$\mathcal{G}$ の有限直和の間の写像の核である。特に $\mathcal{G}$ が
連接ならば，$\SheafHom_{\mathcal{O}_X}(\mathcal{F}, \mathcal{G})$ も連接で
ある。
\end{lemma}

\begin{proof}
最初の主張は補題 \ref{modules-lemma-stalk-internal-hom} の証明で見た。
連接層に関する結果は補題 \ref{modules-lemma-coherent-abelian} から従う。
\end{proof}

\begin{lemma}
\label{modules-lemma-hom-finite-presentation-colimit}
$(X,\mathcal{O}_X)$ を環付き空間とし，$\mathcal{F}$ を有限表示の
$\mathcal{O}_X$-加群の層とする。
$\mathcal{G} = \colim_{\lambda \in \Lambda} \mathcal{G}_\lambda$
を $\mathcal{O}_X$-加群の層のフィルター余極限とする。このとき標準写像
$$
\colim_\lambda \SheafHom_{\mathcal{O}_X}(\mathcal{F}, \mathcal{G}_\lambda)
\longrightarrow
\SheafHom_{\mathcal{O}_X}(\mathcal{F}, \mathcal{G})
$$
は同型である。
\end{lemma}

\begin{proof}
加群の層の余極限を取ることは開集合への制限と可換である
（《層》\ref{sheaves-section-limits-sheaves} 節）。したがって，
$\mathcal{F}$ が大域的表示
$$
\bigoplus\nolimits_{j = 1, \ldots, m}
\mathcal{O}_X
\longrightarrow
\bigoplus\nolimits_{i = 1, \ldots, n}
\mathcal{O}_X
\to
\mathcal{F}
\to
0
$$
をもつと仮定してよい。関手 $\SheafHom_{\mathcal{O}_X}(-,-)$ はどちらの
変数についても有限直和と可換であり，
$\SheafHom_{\mathcal{O}_X}(\mathcal{O}_X,-)$ は恒等関手である。これと補題
\ref{modules-lemma-internal-hom-exact} により，完全列
$$
0 \to
\SheafHom_{\mathcal{O}_X}(\mathcal{F}, \mathcal{G}) \to
\bigoplus\nolimits_{i = 1, \ldots, n} \mathcal{G} \to
\bigoplus\nolimits_{j = 1, \ldots, m} \mathcal{G}
$$
を得る。フィルター余極限は $\textit{Mod}(\mathcal{O}_X)$ において完全なので，
次の可換図式の上段も完全である。
$$
\xymatrix{
0 \ar[r] &
\colim_\lambda \SheafHom_{\mathcal{O}_X}(\mathcal{F}, \mathcal{G}_\lambda)
\ar[r] \ar[d] &
\colim_\lambda \bigoplus\nolimits_{i = 1, \ldots, n} \mathcal{G}_\lambda
\ar[r] \ar[d] &
\colim_\lambda \bigoplus\nolimits_{j = 1, \ldots, m} \mathcal{G}_\lambda
\ar[d] \\
0 \ar[r] &
\SheafHom_{\mathcal{O}_X}(\mathcal{F}, \mathcal{G}) \ar[r] &
\bigoplus\nolimits_{i = 1, \ldots, n} \mathcal{G} \ar[r] &
\bigoplus\nolimits_{j = 1, \ldots, m} \mathcal{G}
}
$$
右側の二つの垂直矢印は同型なので，結論を得る。
\end{proof}

\begin{lemma}
\label{modules-lemma-finite-presentation-quasi-compact-colimit}
$(X,\mathcal{O}_X)$ を環付き空間，$I$ を前順序集合とし，
$(\mathcal{F}_i,\varphi_{ii'})$ を $\mathcal{O}_X$-加群の層からなる $I$ 上の
系とする（《圏》\ref{categories-section-posets-limits} 節）。次を仮定する。
\begin{enumerate}
\item $I$ は有向である。
\item $\mathcal{G}$ は有限表示の $\mathcal{O}_X$-加群の層である。
\item $X$ は，$J$ が有限で，すべての $j,j' \in J$ に対し
$U_j \cap U_{j'}$ が準コンパクトであるような開被覆
$\mathcal{U} : X = \bigcup_{j\in J} U_j$ の共終系をもつ。
\end{enumerate}
このとき
$$
\colim_i \Hom_X(\mathcal{G}, \mathcal{F}_i)
=
\Hom_X(\mathcal{G}, \colim_i \mathcal{F}_i).
$$
が成り立つ。
\end{lemma}

\begin{proof}
$\mathcal{H} = \SheafHom_{\mathcal{O}_X}(\mathcal{G}, \colim \mathcal{F}_i)$
および
$\mathcal{H}_i = \SheafHom_{\mathcal{O}_X}(\mathcal{G}, \mathcal{F}_i)$ とおく。
構成により
$$
\Hom_X(\mathcal{G}, \mathcal{F}) = \Gamma(X, \mathcal{H})
\quad\text{および}\quad
\Hom_X(\mathcal{G}, \mathcal{F}_i) = \Gamma(X, \mathcal{H}_i)
$$
である。補題 \ref{modules-lemma-hom-finite-presentation-colimit} により
$\mathcal{H} = \colim \mathcal{H}_i$ である。したがって補題は《層》補題
\ref{sheaves-lemma-directed-colimits-sections} から従う。
\end{proof}

\begin{remark}
\label{modules-remark-condition-necessary}
上の補題では，$X$ が準コンパクトであるという条件に加えて何らかの条件が
必要である。《層》例 \ref{sheaves-example-conditions-needed-colimit} を
参照せよ。
\end{remark}




\section{加群の層の零化イデアル}
\label{modules-section-annihilator}

\noindent
$(X,\mathcal{O}_X)$ を環付き空間とし，$\mathcal{F}$ を
$\mathcal{O}_X$-加群の層とする。$\mathcal{O}_X$-加群の層の標準写像
$$
\mathcal{O}_X
\longrightarrow
\SheafHom_{\mathcal{O}_X}(\mathcal{F}, \mathcal{F})
$$
があり，局所切断 $f \in \mathcal{O}_X(U)$ を，$f$ 倍で与えられる写像
$f : \mathcal{F}|_U \to \mathcal{F}|_U$ へ送る。

\begin{definition}
\label{modules-definition-annihilator-sheaf}
$(X,\mathcal{O}_X)$ を環付き空間とし，$\mathcal{F}$ を
$\mathcal{O}_X$-加群の層とする。$\mathcal{F}$ の {\it 零化イデアル}
$\text{Ann}_{\mathcal{O}_X}(\mathcal{F})$ を，上で論じた写像
$\mathcal{O}_X \to \SheafHom_{\mathcal{O}_X}(\mathcal{F}, \mathcal{F})$
の核として定義する。
\end{definition}

\noindent
各 $x\in X$ に対し，$\mathcal{O}_{X,x}$ のイデアルの包含
\begin{equation}
\label{modules-equation-inclusion-of-annihilator-ideals}
(\text{Ann}_{\mathcal{O}_X}(\mathcal{F}))_x
\subset
\text{Ann}_{\mathcal{O}_{X, x}}(\mathcal{F}_x)
\end{equation}
がある。実際，$\text{Ann}_{\mathcal{O}_X}(\mathcal{F})$ の任意の切断は，
それが定義されるすべての点で $\mathcal{F}$ の茎を零化する。次に，
(\ref{modules-equation-inclusion-of-annihilator-ideals}) が等号となる簡単な状況を
示す。

\begin{lemma}
\label{modules-lemma-finite-type-annihilator-stalk}
$(X,\mathcal{O}_X)$ を環付き空間とし，$\mathcal{F}$ を
$\mathcal{O}_X$-加群の層とする。$\mathcal{F}$ が有限型ならば，
$(\text{Ann}_{\mathcal{O}_X}(\mathcal{F}))_x =
\text{Ann}_{\mathcal{O}_{X, x}}(\mathcal{F}_x)$.
\end{lemma}

\begin{proof}
補題 \ref{modules-lemma-stalk-internal-hom} により写像
$$
\SheafHom_{\mathcal{O}_X}(\mathcal{F}, \mathcal{F})_x \longrightarrow
\Hom_{\mathcal{O}_{X,x}}(\mathcal{F}_x, \mathcal{F}_x)
$$
は単射である。したがって，$x$ の開近傍 $U$ 上の
$\mathcal{O}_X$ の切断 $f$ が $\mathcal{F}_x$ 上で零として作用するならば，
ある開集合 $U \supset V \ni x$ に対し $\mathcal{F}|_V$ 上でも零として
作用する。ゆえに包含 (\ref{modules-equation-inclusion-of-annihilator-ideals}) は
等号である。
\end{proof}

\begin{lemma}
\label{modules-lemma-induced-quotient-ring-sheaf-module-structure}
$(X,\mathcal{O}_X)$ を環付き空間，$\mathcal{F}$ を
$\mathcal{O}_X$-加群の層，$\mathcal{I} \subset \mathcal{O}_X$ をイデアル層と
する。$\mathcal{I} \subset \text{Ann}_{\mathcal{O}_X}(\mathcal{F})$ ならば，
$\mathcal{F}$ は自然な $\mathcal{O}_X/\mathcal{I}$-加群の層の構造をもち，
茎上では通常の可換代数の構成と一致する。
\end{lemma}

\begin{proof}
包含 $\mathcal{I} \to \mathcal{O}_X$ の余核の普遍性を適用すると，可換図式
$$
\xymatrix{
\mathcal{O}_X \ar[r] \ar[d] &
\SheafHom_{\mathcal{O}_X}(\mathcal{F}, \mathcal{F}) \\
\mathcal{O}_X/\mathcal{I} \ar@{-->}[ur]
}
$$
を $\mathcal{O}_X$-加群の層において得る。補題 \ref{modules-lemma-internal-hom} により，
得られた写像 $\mathcal{O}_X/\mathcal{I} \to
\SheafHom_{\mathcal{O}_X}(\mathcal{F}, \mathcal{F})$
は $\mathcal{O}_X$-加群の層の写像
$$
\mathcal{O}_X/\mathcal{I} \otimes_{\mathcal{O}_X} \mathcal{F}
\longrightarrow
\mathcal{F}
$$
に対応する。これは，与えられた $\mathcal{O}_X$-加群の層の構造と両立する
$\mathcal{F}$ 上の $\mathcal{O}_X/\mathcal{I}$-加群の層の構造を得たことを
意味する。茎に関する主張の確認は省略する。
\end{proof}

\begin{lemma}
\label{modules-lemma-coherent-annihilator}
$(X,\mathcal{O}_X)$ を環付き空間とする。$\mathcal{O}_X$ と $\mathcal{F}$ が
連接ならば，$\text{Ann}_{\mathcal{O}_X}(\mathcal{F})$ も連接である。
\end{lemma}

\begin{proof}
$\text{Ann}_{\mathcal{O}_X}(\mathcal{F})$ は
$\mathcal{O}_X \to \SheafHom_{\mathcal{O}_X}(\mathcal{F}, \mathcal{F})$ の
核なので，補題 \ref{modules-lemma-coherent-abelian} により，
$\SheafHom_{\mathcal{O}_X}(\mathcal{F}, \mathcal{F})$ が連接であることを示せば
十分である。これは補題 \ref{modules-lemma-internal-hom-locally-kernel-direct-sum}，
および $\mathcal{F}$ が連接，したがって特に有限表示であること
（補題 \ref{modules-lemma-coherent-finite-presentation}）から従う。
\end{proof}





\section{Koszul 複体}
\label{modules-section-koszul-complex}

\noindent
まず《代数詳論》\ref{more-algebra-section-koszul} 節の Koszul 複体に関する節を
読むことを勧める。$\mathcal{O}_X$-加群の層の圏における Koszul 複体を
次のように定義する。

\begin{definition}
\label{modules-definition-koszul}
$(X,\mathcal{O}_X)$ を環付き空間とし，
$\varphi : \mathcal{E} \to \mathcal{O}_X$ を $\mathcal{O}_X$-加群の層の写像と
する。$\varphi$ に付随する {\it Koszul 複体} $K_\bullet(\varphi)$ を，
次のように定義される可換微分次数付き代数の層とする。
\begin{enumerate}
\item 台となる次数付き代数は外積代数
$K_\bullet(\varphi) = \wedge(\mathcal{E})$ である。
\item 微分 $d : K_\bullet(\varphi) \to K_\bullet(\varphi)$ は，
$\mathcal{E}=K_1(\varphi)$ のすべての局所切断 $e$ に対して
$d(e)=\varphi(e)$ を満たす一意な導分である。
\end{enumerate}
\end{definition}

\noindent
具体的には，$e_1 \wedge \ldots \wedge e_n$ が $\mathcal{E}$ の局所切断の
外積ならば，
$$
d(e_1 \wedge \ldots \wedge e_n) =
\sum\nolimits_{i = 1, \ldots, n} (-1)^{i + 1}
\varphi(e_i)e_1 \wedge \ldots \wedge \widehat{e_i} \wedge \ldots \wedge e_n.
$$
である。これがテンソル代数上の良定義な導分を与え，$e \wedge e$ を零化し，
したがって外積代数を経由することは直ちに分かる。

\begin{definition}
\label{modules-definition-koszul-complex}
$(X,\mathcal{O}_X)$ を環付き空間とし，
$f_1, \ldots, f_n \in \Gamma(X, \mathcal{O}_X)$ とする。
{\it $f_1,\ldots,f_n$ 上の Koszul 複体}とは，写像
$(f_1, \ldots, f_n) : \mathcal{O}_X^{\oplus n} \to \mathcal{O}_X$
に付随する Koszul 複体である。これを
$K_\bullet(\mathcal{O}_X, f_1, \ldots, f_n)$ または
$K_\bullet(\mathcal{O}_X, f_\bullet)$ と書く。
\end{definition}

\noindent
もちろん，$\mathcal{O}_X$-加群の層の写像
$\varphi : \mathcal{E} \to \mathcal{O}_X$ が与えられ，
$\mathcal{E}$ が有限局所自由ならば，$K_\bullet(\varphi)$ は $X$ 上局所的に
Koszul 複体 $K_\bullet(\mathcal{O}_X, f_1, \ldots, f_n)$ に同型である。



\section{可逆加群の層}
\label{modules-section-invertible}

\noindent
環上の加群の場合（《代数詳論》\ref{more-algebra-section-picard} 節）と同様に，
次の定義を置く。

\begin{definition}
\label{modules-definition-invertible}
$(X,\mathcal{O}_X)$ を環付き空間とする。{\it 可逆
$\mathcal{O}_X$-加群の層}とは，関手
$$
\textit{Mod}(\mathcal{O}_X) \longrightarrow \textit{Mod}(\mathcal{O}_X),\quad
\mathcal{F} \longmapsto \mathcal{L} \otimes_{\mathcal{O}_X} \mathcal{F}
$$
が圏同値となるような $\mathcal{O}_X$-加群の層 $\mathcal{L}$ のことである。
$\mathcal{L}$ が $\mathcal{O}_X$-加群の層として $\mathcal{O}_X$ に同型で
あるとき，$\mathcal{L}$ は {\it 自明}であるという。
\end{definition}

\noindent
下の補題 \ref{modules-lemma-invertible-is-locally-free-rank-1} は，階数 $1$ の
局所自由加群の層との関係を説明する。

\begin{lemma}
\label{modules-lemma-invertible}
$(X,\mathcal{O}_X)$ を環付き空間とし，$\mathcal{L}$ を
$\mathcal{O}_X$-加群の層とする。次は同値である。
\begin{enumerate}
\item $\mathcal{L}$ は可逆である。
\item $\mathcal{O}_X$-加群の層 $\mathcal{N}$ で
$\mathcal{L} \otimes_{\mathcal{O}_X} \mathcal{N} \cong \mathcal{O}_X$.
を満たすものが存在する。
\end{enumerate}
この場合，$\mathcal{L}$ は局所的に有限自由 $\mathcal{O}_X$-加群の層の
直和因子であり，(2) の加群の層 $\mathcal{N}$ は
$\SheafHom_{\mathcal{O}_X}(\mathcal{L}, \mathcal{O}_X)$ に同型である。
\end{lemma}

\begin{proof}
(1) を仮定する。このとき関手 $- \otimes_{\mathcal{O}_X} \mathcal{L}$ は
本質的全射なので，(2) のような $\mathcal{O}_X$-加群の層 $\mathcal{N}$ が
存在する。(2) が成り立つならば，関手
$- \otimes_{\mathcal{O}_X} \mathcal{N}$ は
$- \otimes_{\mathcal{O}_X} \mathcal{L}$ の擬逆であるから，(1) が成り立つ。

\medskip\noindent
(1) と (2) が成り立つと仮定する。与えられた同型を
$\psi : \mathcal{L} \otimes_{\mathcal{O}_X} \mathcal{N} \to \mathcal{O}_X$
と書く。$x \in X$ とする。開近傍 $U$，整数 $n \geq 1$，および切断
$s_i \in \mathcal{L}(U)$，$t_i \in \mathcal{N}(U)$ で
$\psi(\sum s_i \otimes t_i) = 1$ を満たすものを選ぶ。同型
$$
\mathcal{L}|_U \to
\mathcal{L}|_U \otimes_{\mathcal{O}_U}
\mathcal{L}|_U \otimes_{\mathcal{O}_U} \mathcal{N}|_U \to \mathcal{L}|_U
$$
を考える。ここで第1の矢印は $s$ を $\sum s_i \otimes s \otimes t_i$ へ送り，
第2の矢印は $s \otimes s' \otimes t$ を $\psi(s' \otimes t)s$ へ送る。
$s \mapsto \sum \psi(s \otimes t_i)s_i$ は $\mathcal{L}|_U$ の自己同型であると
結論できる。この自己同型は
$$
\mathcal{L}|_U \to \mathcal{O}_U^{\oplus n} \to \mathcal{L}|_U
$$
と分解する。ここで第1の矢印は
$s \mapsto (\psi(s \otimes t_1), \ldots, \psi(s \otimes t_n))$ で，
第2の矢印は $(a_1, \ldots, a_n) \mapsto \sum a_i s_i$ で与えられる。
このようにして，$\mathcal{L}|_U$ は有限自由 $\mathcal{O}_U$-加群の層の
直和因子であると結論する。

\medskip\noindent
(1) と (2) が成り立つと仮定する。評価写像
$$
\mathcal{L} \otimes_{\mathcal{O}_X}
\SheafHom_{\mathcal{O}_X}(\mathcal{L}, \mathcal{O}_X)
\longrightarrow \mathcal{O}_X
$$
を考える。補題の証明を終えるため，これが茎上に同型を誘導することを確認して
同型であると示す。$x \in X$ とする。前段落により $\mathcal{L}$ は有限表示の
$\mathcal{O}_X$-加群の層であることが分かっているので，補題
\ref{modules-lemma-stalk-internal-hom} により，
$$
\mathcal{L}_x \otimes_{\mathcal{O}_{X, x}}
\Hom_{\mathcal{O}_{X, x}}(\mathcal{L}_x, \mathcal{O}_{X, x})
\longrightarrow \mathcal{O}_{X, x}
$$
が同型であることを示せば十分である。
$\mathcal{L}_x \otimes_{\mathcal{O}_{X, x}} \mathcal{N}_x =
(\mathcal{L} \otimes_{\mathcal{O}_X} \mathcal{N})_x =
\mathcal{O}_{X, x}$ なので（補題 \ref{modules-lemma-stalk-tensor-product}），所望の結果は
《代数詳論》補題 \ref{more-algebra-lemma-invertible} から従う。
\end{proof}

\begin{lemma}
\label{modules-lemma-pullback-invertible}
$f : (X,\mathcal{O}_X) \to (Y,\mathcal{O}_Y)$ を環付き空間の射とする。可逆
$\mathcal{O}_Y$-加群の層 $\mathcal{L}$ の引き戻し $f^*\mathcal{L}$ は可逆で
ある。
\end{lemma}

\begin{proof}
補題 \ref{modules-lemma-invertible} により，
$\mathcal{L} \otimes_{\mathcal{O}_Y} \mathcal{N} \cong \mathcal{O}_Y$ を
満たす $\mathcal{O}_Y$-加群の層 $\mathcal{N}$ が存在する。引き戻すと，補題
\ref{modules-lemma-tensor-product-pullback} により
$f^*\mathcal{L} \otimes_{\mathcal{O}_X} f^*\mathcal{N} \cong \mathcal{O}_X$
を得る。したがって補題 \ref{modules-lemma-invertible} により
$f^*\mathcal{L}$ は可逆である。
\end{proof}

\begin{lemma}
\label{modules-lemma-invertible-is-locally-free-rank-1}
$(X,\mathcal{O}_X)$ を環付き空間とする。階数 $1$ の任意の局所自由
$\mathcal{O}_X$-加群の層は可逆である。すべての茎 $\mathcal{O}_{X,x}$ が
局所環ならば逆も成り立つ（ただし一般には成り立たない）。
\end{lemma}

\begin{proof}
括弧内の主張は，環 $R$ が与える環の層 $\mathcal{O}_X$ を備えた一点空間
$X$ を考えれば従う。このとき可逆 $\mathcal{O}_X$-加群の層は可逆
$R$-加群に対応するので，$\Pic(R)$ が自明群でなければ例を得る。

\medskip\noindent
$\mathcal{L}$ が階数 $1$ の局所自由加群の層であると仮定し，評価写像
$$
\mathcal{L} \otimes_{\mathcal{O}_X}
\SheafHom_{\mathcal{O}_X}(\mathcal{L}, \mathcal{O}_X)
\longrightarrow \mathcal{O}_X
$$
を考える。$\mathcal{L}$ を自明化する開被覆上で見れば，この写像は同型である。
したがって補題 \ref{modules-lemma-invertible} により $\mathcal{L}$ は可逆である。

\medskip\noindent
すべての茎 $\mathcal{O}_{X,x}$ が局所環で，$\mathcal{L}$ が可逆であると
仮定する。補題 \ref{modules-lemma-invertible} の証明で，すべての $x \in X$ に対し
$\mathcal{L}_x$ は可逆 $\mathcal{O}_{X,x}$-加群であることを見た。
$\mathcal{O}_{X,x}$ は局所環なので，
$\mathcal{L}_x \cong \mathcal{O}_{X, x}$
である（《代数詳論》\ref{more-algebra-section-picard} 節）。補題
\ref{modules-lemma-invertible} により $\mathcal{L}$ は有限表示なので，補題
\ref{modules-lemma-finite-presentation-stalk-free} により $\mathcal{L}$ は階数 $1$ の
局所自由加群の層であると結論する。
\end{proof}

\begin{lemma}
\label{modules-lemma-constructions-invertible}
$(X,\mathcal{O}_X)$ を環付き空間とする。
\begin{enumerate}
\item $\mathcal{L},\mathcal{N}$ が可逆 $\mathcal{O}_X$-加群の層ならば，
$\mathcal{L} \otimes_{\mathcal{O}_X} \mathcal{N}$ も可逆である。
\item $\mathcal{L}$ が可逆 $\mathcal{O}_X$-加群の層ならば，
$\SheafHom_{\mathcal{O}_X}(\mathcal{L}, \mathcal{O}_X)$ も可逆であり，評価写像
$\mathcal{L} \otimes_{\mathcal{O}_X}
\SheafHom_{\mathcal{O}_X}(\mathcal{L}, \mathcal{O}_X) \to \mathcal{O}_X$
は同型である。
\end{enumerate}
\end{lemma}

\begin{proof}
(1) は定義から明らかであり，(2) は補題 \ref{modules-lemma-invertible} とその証明から
従う。
\end{proof}

\begin{definition}
\label{modules-definition-powers}
$(X,\mathcal{O}_X)$ を環付き空間とする。$X$ 上の可逆層
$\mathcal{L}$ と $n \in \mathbf{Z}$ に対し，$\mathcal{L}$ の第 $n$
{\it テンソル冪} $\mathcal{L}^{\otimes n}$ を，同値
$\mathcal{F} \mapsto\mathcal{F} \otimes_{\mathcal{O}_X} \mathcal{L}$
をちょうど $n$ 回適用したときの $\mathcal{O}_X$ の像として定義する。
\end{definition}

\noindent
可逆 $\mathcal{O}_X$-加群の層を，それとのテンソル積が同値となるものとして
定義したので，これは負の $n$ に対しても意味をもつ。より明示的には，
$$
\mathcal{L}^{\otimes n} =
\left\{
\begin{matrix}
\mathcal{O}_X & \text{（} & n = 0\text{ のとき）} \\
\SheafHom_{\mathcal{O}_X}(\mathcal{L}, \mathcal{O}_X) & \text{（} & n = -1\text{ のとき）}\\
\mathcal{L} \otimes_{\mathcal{O}_X} \ldots \otimes_{\mathcal{O}_X} \mathcal{L}
& \text{（} & n > 0\text{ のとき）} \\
\mathcal{L}^{\otimes -1} \otimes_{\mathcal{O}_X} \ldots
\otimes_{\mathcal{O}_X} \mathcal{L}^{\otimes -1}
& \text{（} & n < -1\text{ のとき）}
\end{matrix}
\right.
$$
である（補題 \ref{modules-lemma-constructions-invertible} 参照）。
この定義のもとで標準同型
$\mathcal{L}^{\otimes n} \otimes_{\mathcal{O}_X}
\mathcal{L}^{\otimes m} \to
\mathcal{L}^{\otimes n + m}$ があり，これらの同型は可換性制約と
結合性制約を満たす（定式化は省略する）。

\medskip\noindent
$(X,\mathcal{O}_X)$ を環付き空間とする。
$s \in \Gamma(X, \mathcal{L}^{\otimes n})$ と
$t \in \Gamma(X, \mathcal{L}^{\otimes m})$ を，
$\Gamma(X, \mathcal{L}^{\otimes n + m})$ において $s \otimes t$ に
対応する切断へ写すことにより，
$\bigoplus \Gamma(X, \mathcal{L}^{\otimes n})$ 上に
$\mathbf{Z}$-次数付き環の構造を定義できる。これが単位元 $1$ をもつ
可換結合的環を定めることの確認は省略する。しかし，《代数》
\ref{algebra-section-graded} 節の規約では，次数付き環は負の次数に
非零元をもたない。そこで次の定義を置く。

\begin{definition}
\label{modules-definition-gamma-star}
$(X,\mathcal{O}_X)$ を環付き空間とする。$X$ 上の可逆層
$\mathcal{L}$ に対し，{\it 付随次数付き環}を
$$
\Gamma_*(X, \mathcal{L})
=
\bigoplus\nolimits_{n \geq 0} \Gamma(X, \mathcal{L}^{\otimes n})
$$
と定義する。$\mathcal{O}_X$-加群の層 $\mathcal{F}$ が与えられたとき，
$$
\Gamma_*(X, \mathcal{L}, \mathcal{F})
=
\bigoplus\nolimits_{n \in \mathbf{Z}} \Gamma(X,
\mathcal{F} \otimes_{\mathcal{O}_X} \mathcal{L}^{\otimes n})
$$
と置き，これを次数付き $\Gamma_*(X,\mathcal{L})$-加群とみなす。
\end{definition}

\noindent
しばしば単に $\Gamma_*(\mathcal{L})$ および $\Gamma_*(\mathcal{F})$
と書く（ただし $\mathcal{F}$ が可逆ならば曖昧である）。
$\Gamma_*(\mathcal{F})$ への $\Gamma_*(\mathcal{L})$ の乗法は，
上の同型を用いて定義する。
$\gamma : \mathcal{F} \to \mathcal{G}$ が $\mathcal{O}_X$-加群の層の射ならば，
$\Gamma_*(\mathcal{L})$-加群準同型
$\gamma : \Gamma_*(\mathcal{F}) \to \Gamma_*(\mathcal{G})$
を得る。$\alpha : \mathcal{L} \to \mathcal{N}$ が可逆
$\mathcal{O}_X$-加群の層の間の射ならば，次数付き環準同型
$\Gamma_*(\mathcal{L}) \to \Gamma_*(\mathcal{N})$ を得る。
$f : (Y,\mathcal{O}_Y) \to (X,\mathcal{O}_X)$ が環付き空間の射で，
$\mathcal{L}$ が $X$ 上可逆ならば，$Y$ 上の可逆層
$f^*\mathcal{L}$（補題 \ref{modules-lemma-pullback-invertible}）と，
誘導された次数付き環準同型
$$
f^* :
\Gamma_*(X, \mathcal{L})
\longrightarrow
\Gamma_*(Y, f^*\mathcal{L})
$$
を得る。さらに，上の構成の間にはいくつかの整合性があるが，その記述は省略する。

\begin{lemma}
\label{modules-lemma-pic-set}
$(X,\mathcal{O}_X)$ を環付き空間とする。可逆加群の層の集合
$\{\mathcal{L}_i\}_{i \in I}$ で，$X$ 上の各可逆加群の層がちょうど一つの
$\mathcal{L}_i$ と同型になるものが存在する。
\end{lemma}

\begin{proof}
補題 \ref{modules-lemma-invertible} により，任意の可逆 $\mathcal{O}_X$-加群の層は，
局所的に有限自由 $\mathcal{O}_X$-加群の層の直和因子であることを思い出そう。
各開被覆 $\mathcal{U} : X = \bigcup_{j \in J} U_j$ と写像
$r : J \to \mathbf{N}$ に対し，
$\mathcal{F}_j = \mathcal{F}|_{U_j}$ が
$\mathcal{O}_{U_j}^{\oplus r(j)}$ の直和因子となる
$\mathcal{O}_X$-加群の層 $\mathcal{F}$ を考える。
$\Hom_{\mathcal{O}_U}(\mathcal{O}_U^{\oplus r}, \mathcal{O}_U^{\oplus r})$
は集合なので，$\mathcal{F}_j$ の同型類全体は集合をなす。層
$\mathcal{F}$ は $\mathcal{F}_j$ を貼り合わせて得られる
（《層》\ref{sheaves-section-glueing-sheaves} 節参照）。すべての貼り合わせ
データ全体が集合をなすことに注意する。写像 $J \to \mathcal{P}(X)$，
$j \mapsto U_j$ が単射となる被覆
$\mathcal{U} : X = \bigcup_{j \in J} U_i$ 全体も集合をなす。
各被覆に対し，写像 $r : J \to \mathbf{N}$ 全体は集合である。
したがって，すべての $\mathcal{F}$ 全体は集合をなす。
\end{proof}

\noindent
大まかにいえば，この補題は可逆層の同型類全体が集合をなすことを述べている。
補題 \ref{modules-lemma-constructions-invertible} により，テンソル積はこの集合上に
アーベル群の構造を定める。

\begin{definition}
\label{modules-definition-pic}
$(X,\mathcal{O}_X)$ を環付き空間とする。$X$ の {\it Picard 群}
$\Pic(X)$ とは，可逆 $\mathcal{O}_X$-加群の層の同型類を元とし，
テンソル積に対応する演算を加法とするアーベル群である。
\end{definition}

\begin{lemma}
\label{modules-lemma-s-open}
\begin{slogan}
直線束の（局所）自明化は（局所）非消滅切断と同じものである。
\end{slogan}
$X$ を環付き空間とする。各茎 $\mathcal{O}_{X,x}$ は極大イデアル
$\mathfrak m_x$ をもつ局所環であると仮定する。
$\mathcal{L}$ を可逆 $\mathcal{O}_X$-加群の層とする。
任意の切断 $s \in \Gamma(X,\mathcal{L})$ に対し，集合
$$
X_s = \{x \in X \mid \text{像 }s \not\in \mathfrak m_x\mathcal{L}_x\}
$$
は $X$ の開集合である。写像
$s : \mathcal{O}_{X_s} \to \mathcal{L}|_{X_s}$ は同型であり，
$s'(s|_{X_s}) = 1$ を満たす $\mathcal{L}^{\otimes -1}$ の
$X_s$ 上の切断 $s'$ が存在する。
\end{lemma}

\begin{proof}
$x \in X_s$ とする。補題 \ref{modules-lemma-constructions-invertible} により同型
$$
\mathcal{L}_x \otimes_{\mathcal{O}_{X, x}} (\mathcal{L}^{\otimes -1})_x
\longrightarrow
\mathcal{O}_{X, x}
$$
がある。$\mathcal{L}_x$ と $(\mathcal{L}^{\otimes -1})_x$ はともに階数
$1$ の自由 $\mathcal{O}_{X,x}$-加群である。《代数》の Nakayama の補題
\ref{algebra-lemma-NAK} により，$s_x$ は $\mathcal{L}_x$ の基底であると
結論する。したがって，$s_x \otimes t_x$ が $1$ へ写るような基底元
$t_x \in (\mathcal{L}^{\otimes -1})_x$ が存在する。
$t_x$ が $\mathcal{L}^{\otimes -1}$ の $U$ 上の切断 $t$ から来て，
$s \otimes t$ が $1 \in \mathcal{O}_X(U)$ へ写るような $x$ の開近傍
$U$ を選ぶ。明らかに，すべての $x' \in U$ に対して $s$ は
$\mathcal{L}_{x'}$ を生成する。したがって $U \subset X_s$ であり，
$X_s$ は開である。さらに，上で $U$ 上に構成した切断 $t$ は一意なので，
これらを貼り合わせて補題の切断 $s'$ を得る。
\end{proof}

\begin{remark}
\label{modules-remark-pullback-s-open}
局所環付き空間の射 $f : Y \to X$ が与えられたとする
（《スキーム》定義 \ref{schemes-definition-locally-ringed-space} 参照）。
このとき逆像 $f^{-1}(X_s)$ は $Y_{f^*s}$ に等しい。ここで
$f^*s \in \Gamma(Y,f^*\mathcal{L})$ は $s$ の引き戻しである。
\end{remark}


\section{階数と行列式}
\label{modules-section-rank-and-det}

\noindent
$(X,\mathcal{O}_X)$ を環付き空間とする。有限局所自由
$\mathcal{O}_X$-加群の層の圏 $\textit{Vect}(X)$ を考える。
これは，許容全射が全射で，許容単射が全射の核である完全圏である
（《入射対象》注 \ref{injectives-remark-embed-exact-category} 参照）。
さらに，$\textit{Vect}(X)$ の対象の同型類全体は集合をなす
（証明は省略する）。したがって，第 $0$ Grothendieck $K$-群
$K_0(\textit{Vect}(X))$ を構成できる。明示的には，この場合
$K_0(\textit{Vect}(X))$ は，有限局所自由 $\mathcal{O}_X$-加群の層
$\mathcal{E}$ に対する $[\mathcal{E}]$ で生成され，有限局所自由
$\mathcal{O}_X$-加群の層の短完全列
$0 \to \mathcal{E}' \to \mathcal{E} \to \mathcal{E}'' \to 0$
があるごとに関係式
$$
[\mathcal{E}] = [\mathcal{E}'] + [\mathcal{E}'']
$$
を課したアーベル群である。

\medskip\noindent
{\bf 階数。} すべての茎 $\mathcal{O}_{X,x}$ が零環でないと仮定する。
有限局所自由 $\mathcal{O}_X$-加群の層 $\mathcal{E}$ に対し，
{\it 階数}は局所定数関数
$$
\text{rank}_\mathcal{E} : X \longrightarrow \mathbf{Z}_{\geq 0},\quad
x \longmapsto \text{rank}_{\mathcal{O}_{X, x}} \mathcal{E}_x
$$
である（補題 \ref{modules-lemma-rank} 参照）。局所自由加群の層の定義により，
関数 $\text{rank}_\mathcal{E}$ は局所定数である。有限局所自由
$\mathcal{O}_X$-加群の層の短完全列
$0 \to \mathcal{E}' \to \mathcal{E} \to \mathcal{E}'' \to 0$
があれば，
$\text{rank}_\mathcal{E} = \text{rank}_{\mathcal{E}'} +
\text{rank}_{\mathcal{E}''}$ である。したがって階数は準同型
$$
K_0(\textit{Vect}(X)) \longrightarrow \text{Map}_{cont}(X, \mathbf{Z}),\quad
[\mathcal{E}] \longmapsto \text{rank}_\mathcal{E}
$$
を定める。

\medskip\noindent
{\bf 行列式。} 有限局所自由 $\mathcal{O}_X$-加群の層 $\mathcal{E}$ に
対し，非交和分解
$$
X = X_0 \amalg X_1 \amalg X_2 \amalg \ldots
$$
を得る。ここで各 $X_i$ は開かつ閉であり，$\mathcal{E}$ は $X_i$ 上で
階数 $i$ の有限局所自由加群の層である（これは
$\text{rank}_\mathcal{E}$ が局所定数であることとまったく同じである）。
この場合，$\det(\mathcal{E})$ を，各 $i \geq 0$ に対して $X_i$ 上で
$\wedge^i(\mathcal{E}|_{X_i})$ に等しい $X$ 上の可逆層として定義する。
上の分解は非交和なので，確認すべき貼り合わせ条件はない。下の補題
\ref{modules-lemma-det-ses} により，これはアーベル群の準同型
$$
\det : K_0(\textit{Vect}(X)) \longrightarrow \Pic(X),\quad
[\mathcal{E}] \longmapsto \det(\mathcal{E})
$$
を定める。このようにして得られる $\Pic(X)$ の元は階数 $1$ の局所自由層で
表される（一般化については補題の後を参照）。

\begin{lemma}
\label{modules-lemma-det-ses}
$X$ を環付き空間とする。有限局所自由 $\mathcal{O}_X$-加群の層の短完全列
$0 \to \mathcal{E}' \to \mathcal{E} \to \mathcal{E}'' \to 0$
があるとする。このとき $\mathcal{O}_X$-加群の層の標準同型
$$
\det(\mathcal{E}') \otimes_{\mathcal{O}_X}\det(\mathcal{E}'')
\longrightarrow
\det(\mathcal{E})
$$
がある。
\end{lemma}

\begin{proof}
$X$ を互いに素な開閉集合に分解し，その各々の上で
$\mathcal{E}'$ と $\mathcal{E}''$ の階数を定数にできる。
したがって，$\mathcal{E}'$ と $\mathcal{E}''$ の階数がそれぞれ定数
$r'$ と $r''$ である場合に帰着する。この場合，写像
$$
\wedge^{r'}(\mathcal{E}') \otimes_{\mathcal{O}_X} \wedge^{r''}(\mathcal{E}'')
\longrightarrow
\wedge^{r' + r''}(\mathcal{E})
$$
を次のように定義する。$\mathcal{E}'$ の局所切断
$s'_1,\ldots,s'_{r'}$ と $\mathcal{E}''$ の局所切断
$s''_1,\ldots,s''_{r''}$ に対し，
$$
s'_1 \wedge \ldots \wedge s'_{r'} \otimes
s''_1 \wedge \ldots \wedge s''_{r''}
\quad\text{を}\quad
s'_1 \wedge \ldots \wedge s'_{r'} \wedge
\tilde s''_1 \wedge \ldots \wedge \tilde s''_{r''}
$$
へ写す。ここで $\tilde s''_i$ は切断 $s''_i$ の $\mathcal{E}$ の切断への
局所持ち上げである。詳細は省略する。
\end{proof}

\noindent
$(X,\mathcal{O}_X)$ を環付き空間とする。有限局所自由
$\mathcal{O}_X$-加群の層の代わりに，$X$ 上局所的に有限自由
$\mathcal{O}_X$-加群の層の直和因子となる $\mathcal{O}_X$-加群の層
$\mathcal{F}$ を考えることができる。これは，$\mathcal{F}$ が有限表示の
平坦 $\mathcal{O}_X$-加群の層であることと同じである
（補題 \ref{modules-lemma-flat-locally-finite-presentation} 参照）。
すべての茎 $\mathcal{O}_{X,x}$ が局所環ならば，そのような
$\mathcal{F}$ は有限局所自由である
（補題 \ref{modules-lemma-direct-summand-of-locally-free-is-locally-free} 参照）。
しかし一般にはそうならない。例えば，$X$ が一点で，
$\Gamma(X,\mathcal{O}_X)$ が二つの非零環の積 $A \times B$ であり，
$\mathcal{F}$ が $A \times 0$ に対応する場合がある。
したがって，そのような加群の層では階数関数は定義されない。
それでも $\det(\mathcal{F})$ を定義でき，これは定義
\ref{modules-definition-invertible} の意味で可逆 $\mathcal{O}_X$-加群の層となる
（階数 $1$ の局所自由層とは限らない）。$\mathcal{F}$ が有限局所自由ならば，
この構成は上の構成と一致する。読者には本節の残りを読み飛ばすことを勧める。

\begin{lemma}
\label{modules-lemma-determinant-as-socle}
$(X,\mathcal{O}_X)$ を環付き空間とし，$\mathcal{F}$ を有限表示の平坦
$\mathcal{O}_X$-加群の層とする。
$$
\det(\mathcal{F}) \subset
\wedge^*_{\mathcal{O}_X}(\mathcal{F})
$$
で，$\mathcal{F} \subset \wedge^*_{\mathcal{O}_X}(\mathcal{F})$ の
零化イデアルを表す。このとき $\det(\mathcal{F})$ は可逆
$\mathcal{O}_X$-加群の層である。
\end{lemma}

\begin{proof}
これは $X$ 上局所的に証明してよい。したがって補題
\ref{modules-lemma-flat-locally-finite-presentation} により，$\mathcal{F}$ は
有限自由加群の層の直和因子であると仮定してよい。
$\mathcal{F} \oplus \mathcal{G} = \mathcal{O}_X^{\oplus n}$
と書き，$R = \mathcal{O}_X(X)$ と置く。このとき
$\mathcal{F}(X) \oplus \mathcal{G}(X) = R^{\oplus n}$
であり，これに応じてすべての開集合 $U \subset X$ に対し
$\mathcal{F}(U) \oplus \mathcal{G}(U) = \mathcal{O}_X(U)^{\oplus n}$
である。補題 \ref{modules-lemma-construct-quasi-coherent-sheaves} の記法で
$\mathcal{F} = \mathcal{F}_M$ であり，
$M = \mathcal{F}(X)$ は有限射影 $R$-加群であると結論する。
すなわち，$\mathcal{F}(U) = M \otimes_R \mathcal{O}_X(U)$ である。
したがって，すべての開集合 $U \subset X$ に対し
$\det(M) \otimes_R \mathcal{O}_X(U) = \det(\mathcal{F}(U))$
である。ここで $\det$ は《代数詳論》
\ref{more-algebra-section-determinants} 節のものである。《代数詳論》注
\ref{more-algebra-remark-determinant-as-socle} により，
$$
\det(M) \otimes_R \mathcal{O}_X(U) =
\det(\mathcal{F}(U)) \subset
\wedge^*_{\mathcal{O}_X(U)}(\mathcal{F}(U))
$$
は $\mathcal{F}(U)$ の零化イデアルである。したがって，補題の主張で
定義した $\det(\mathcal{F})$ は $\mathcal{F}_{\det(M)}$ に等しい。
いくつかの詳細は省略する。零化イデアルは，開集合上の切断について
零化イデアルをとってから層化することでは定義できないので，注意が必要である。
ゆえに $\det(\mathcal{F})$ は可逆加群の引き戻しであり，結論を得る。
\end{proof}
















\section{環の層の局所化}
\label{modules-section-localizing-sheaves-rings}

\noindent
$X$ を位相空間，$\mathcal{O}_X$ を環の前層とする。
$\mathcal{S} \subset \mathcal{O}_X$ を $\mathcal{O}_X$ に含まれる
集合の前層とする。すべての開集合 $U \subset X$ に対し，
$\mathcal{S}(U) \subset \mathcal{O}_X(U)$ は乗法的部分集合であると仮定する
（《代数》定義 \ref{algebra-definition-multiplicative-subset} 参照）。
この場合，環の前層
$$
\mathcal{S}^{-1}\mathcal{O}_X :
U \longmapsto \mathcal{S}(U)^{-1}\mathcal{O}_X(U).
$$
を考えることができる。$V \subset U$ が $X$ の開集合ならば，制限写像は
切断 $f/s$（$f \in \mathcal{O}_X(U)$，$s \in \mathcal{S}(U)$）を
$(f|_V)/(s|_V)$ へ写す。

\begin{lemma}
\label{modules-lemma-simple-invert}
$X$ を位相空間，$\mathcal{O}_X$ を環の前層とする。
$\mathcal{S} \subset \mathcal{O}_X$ を $\mathcal{O}_X$ に含まれる
集合の前層とする。すべての開集合 $U \subset X$ に対し，
$\mathcal{S}(U) \subset \mathcal{O}_X(U)$ は乗法的部分集合であると仮定する。
\begin{enumerate}
\item 環の前層の射
$\mathcal{O}_X \to \mathcal{S}^{-1}\mathcal{O}_X$
で，$\mathcal{S}$ の各局所切断が $\mathcal{O}_X$ の可逆な切断へ
写るものが存在する。
\item $\mathcal{S}$ の各局所切断を $\mathcal{A}$ の可逆な切断へ写す
環の前層の任意の準同型 $\mathcal{O}_X \to \mathcal{A}$ に対し，一意な分解
$\mathcal{S}^{-1}\mathcal{O}_X \to \mathcal{A}$
が存在する。
\item 任意の $x \in X$ に対し，
$$
(\mathcal{S}^{-1}\mathcal{O}_X)_x = \mathcal{S}_x^{-1} \mathcal{O}_{X, x}.
$$
\item 層化 $(\mathcal{S}^{-1}\mathcal{O}_X)^\#$ は環の層であり，
環の層の射
$(\mathcal{O}_X)^\# \to (\mathcal{S}^{-1}\mathcal{O}_X)^\#$
を備える。この射は，$\mathcal{S}$ の各局所切断を可逆な切断へ写す
$(\mathcal{O}_X)^\#$ から環の層への射に関して普遍である。
\item 任意の $x \in X$ に対し，
$$
(\mathcal{S}^{-1}\mathcal{O}_X)^\#_x = \mathcal{S}_x^{-1} \mathcal{O}_{X, x}.
$$
\end{enumerate}
\end{lemma}

\begin{proof}
省略する。
\end{proof}

\noindent
$X$ を位相空間，$\mathcal{O}_X$ を環の前層とする。
$\mathcal{S} \subset \mathcal{O}_X$ を $\mathcal{O}_X$ に含まれる
集合の前層とする。すべての開集合 $U \subset X$ に対し，
$\mathcal{S}(U) \subset \mathcal{O}_X(U)$ は乗法的部分集合であると仮定する。
$\mathcal{F}$ を $\mathcal{O}_X$-加群の前層とする。この場合，
$\mathcal{S}^{-1}\mathcal{O}_X$-加群の前層
$$
\mathcal{S}^{-1}\mathcal{F} :
U \longmapsto \mathcal{S}(U)^{-1}\mathcal{F}(U).
$$
を考えることができる。$V \subset U$ が $X$ の開集合ならば，制限写像は
切断 $t/s$（$t \in \mathcal{F}(U)$，$s \in \mathcal{S}(U)$）を
$(t|_V)/(s|_V)$ へ写す。

\begin{lemma}
\label{modules-lemma-simple-invert-module}
$X$ を位相空間，$\mathcal{O}_X$ を環の前層とする。
$\mathcal{S} \subset \mathcal{O}_X$ を $\mathcal{O}_X$ に含まれる
集合の前層とする。すべての開集合 $U \subset X$ に対し，
$\mathcal{S}(U) \subset \mathcal{O}_X(U)$ は乗法的部分集合であると仮定する。
任意の $\mathcal{O}_X$-加群の前層 $\mathcal{F}$ に対し，
$$
\mathcal{S}^{-1}\mathcal{F}
=
\mathcal{S}^{-1}\mathcal{O}_X \otimes_{p, \mathcal{O}_X} \mathcal{F}
$$
である（記法については《層》
\ref{sheaves-section-presheaves-modules} 節参照）。また，
$\mathcal{F}$ と $\mathcal{O}_X$ が層ならば，
$$
(\mathcal{S}^{-1}\mathcal{F})^\#
=
(\mathcal{S}^{-1}\mathcal{O}_X)^\# \otimes_{\mathcal{O}_X} \mathcal{F}
$$
である（記法については《層》
\ref{sheaves-section-sheafification-presheaves-modules} 節参照）。
\end{lemma}

\begin{proof}
省略する。
\end{proof}








\section{微分加群の層}
\label{modules-section-differentials}

\noindent
本節では，環付き空間の射に対する相対微分加群の層の定義を簡潔に説明する。
読者には，可換代数の章の対応する節
（《代数》\ref{algebra-section-differentials} 節）も参照することを勧める。

\begin{definition}
\label{modules-definition-derivation}
$X$ を位相空間とし，
$\varphi : \mathcal{O}_1 \to \mathcal{O}_2$ を環の層の準同型とする。
$\mathcal{F}$ を $\mathcal{O}_2$-加群の層とする。$\mathcal{F}$ への
{\it $\mathcal{O}_1$-導分}，より正確には {\it $\varphi$-導分}とは，
加法的で，$\mathcal{O}_1 \to \mathcal{O}_2$ の像を零化し，
{\it Leibniz 則}
$$
D(ab) = aD(b) + D(a)b
$$
を満たす写像 $D : \mathcal{O}_2 \to \mathcal{F}$ のことである。
ここで $a,b$ は $\mathcal{O}_2$ の，同時に定義される任意の局所切断である。
$\mathcal{F}$ への $\varphi$-導分全体の集合を
$\text{Der}_{\mathcal{O}_1}(\mathcal{O}_2,\mathcal{F})$ と書く。
\end{definition}

\noindent
これは《代数》定義 \ref{algebra-definition-derivation} の層論的類似である。
定義のような導分 $D : \mathcal{O}_2 \to \mathcal{F}$ が与えられると，
大域切断上の写像
$$
D : \Gamma(X, \mathcal{O}_2) \longrightarrow \Gamma(X, \mathcal{F})
$$
は代数における定義の意味で $\Gamma(X,\mathcal{O}_1)$-導分である。
$\alpha : \mathcal{F} \to \mathcal{G}$ が $\mathcal{O}_2$-加群の層の射ならば，
誘導された写像
$$
\text{Der}_{\mathcal{O}_1}(\mathcal{O}_2, \mathcal{F})
\longrightarrow
\text{Der}_{\mathcal{O}_1}(\mathcal{O}_2, \mathcal{G})
$$
があり，$D \mapsto \alpha \circ D$ で与えられることに注意する。
言い換えれば，関手を得る。

\begin{lemma}
\label{modules-lemma-universal-module}
$X$ を位相空間とし，
$\varphi : \mathcal{O}_1 \to \mathcal{O}_2$ を環の層の準同型とする。関手
$$
\textit{Mod}(\mathcal{O}_2) \longrightarrow \textit{Ab}, \quad
\mathcal{F} \longmapsto \text{Der}_{\mathcal{O}_1}(\mathcal{O}_2, \mathcal{F})
$$
は表現可能である。
\end{lemma}

\begin{proof}
代数における類似の主張とまったく同じ方法で証明する。
この証明では，$X$ 上の任意の集合の層 $\mathcal{F}$ に対し，
前層 $U \mapsto \mathcal{O}_2(U)[\mathcal{F}(U)]$ の層化を
$\mathcal{O}_2[\mathcal{F}]$ と書く。ここで
$\mathcal{O}_2(U)[\mathcal{F}(U)]$ は集合 $\mathcal{F}(U)$ を基底とする
自由 $\mathcal{O}_2(U)$-加群である。$s \in \mathcal{F}(U)$ に対し，
対応する $\mathcal{O}_2[\mathcal{F}]$ の $U$ 上の切断を $[s]$ と書く。
$\mathcal{F}$ が $\mathcal{O}_2$-加群の層ならば，標準写像
$$
c : \mathcal{O}_2[\mathcal{F}] \longrightarrow \mathcal{F}
$$
があり，前層のレベルでは
$\sum f_s[s] \mapsto \sum f_s s$ で与えられる。この写像，および以下の
他の写像を記述する際には，略記 $[s] \mapsto s$ を用いる。
$\mathcal{O}_2$-加群の層の写像
\begin{equation}
\label{modules-equation-define-module-differentials}
\begin{matrix}
\mathcal{O}_2[\mathcal{O}_2 \times \mathcal{O}_2] \oplus
\mathcal{O}_2[\mathcal{O}_2 \times \mathcal{O}_2] \oplus
\mathcal{O}_2[\mathcal{O}_1] &
\longrightarrow &
\mathcal{O}_2[\mathcal{O}_2] \\
[(a, b)] \oplus [(f, g)] \oplus [h] & \longmapsto & [a + b] - [a] - [b] + \\
& & [fg] - g[f] - f[g] + \\
& & [\varphi(h)]
\end{matrix}
\end{equation}
を考える。ここでは上の略記を用いた。この写像の余核を
$\Omega_{\mathcal{O}_2/\mathcal{O}_1}$ と置く。このとき集合の層の写像
$$
\text{d} : \mathcal{O}_2 \longrightarrow \Omega_{\mathcal{O}_2/\mathcal{O}_1}
$$
で，局所切断 $f$ を $\Omega_{\mathcal{O}_2/\mathcal{O}_1}$ における
$[f]$ の像へ写すものが明らかに存在する。構成により $\text{d}$ は
$\mathcal{O}_1$-導分である。次に，$\mathcal{F}$ を
$\mathcal{O}_2$-加群の層，
$D : \mathcal{O}_2 \to \mathcal{F}$ を $\mathcal{O}_1$-導分とする。
$[g]$ を $D(g)$ へ写す $\mathcal{O}_2$-線形写像
$\mathcal{O}_2[\mathcal{O}_2] \to \mathcal{F}$ を考えることができる。
導分の定義により，この写像は写像
(\ref{modules-equation-define-module-differentials}) の像に属する切断を零化するので，
写像
$$
\alpha_D : \Omega_{\mathcal{O}_2/\mathcal{O}_1} \longrightarrow \mathcal{F}
$$
を定める。$D = \alpha_D \circ \text{d}$ は明らかなので，補題が従う。
\end{proof}

\begin{definition}
\label{modules-definition-module-differentials}
$X$ を位相空間とし，
$\varphi : \mathcal{O}_1 \to \mathcal{O}_2$ を $X$ 上の環の層の準同型とする。
$\varphi$ の {\it 微分加群の層}とは，関手
$\mathcal{F} \mapsto \text{Der}_{\mathcal{O}_1}(\mathcal{O}_2, \mathcal{F})$
を表現する対象であり，補題 \ref{modules-lemma-universal-module} により存在する。
これを $\Omega_{\mathcal{O}_2/\mathcal{O}_1}$ と書き，
{\it 普遍 $\varphi$-導分}を
$\text{d} : \mathcal{O}_2 \to \Omega_{\mathcal{O}_2/\mathcal{O}_1}$
と書く。
\end{definition}

\noindent
$\Omega_{\mathcal{O}_2/\mathcal{O}_1}$ は $\mathcal{O}_2$-加群の層の写像
(\ref{modules-equation-define-module-differentials}) の余核であることに注意する。
また，写像 $\text{d}$ は，$\text{d}f$ を局所切断 $[f]$ の像とする規則で
記述される。

\begin{lemma}
\label{modules-lemma-differentials-sheafify}
$X$ を位相空間とし，
$\varphi : \mathcal{O}_1 \to \mathcal{O}_2$ を $X$ 上の環の層の準同型とする。
このとき $\Omega_{\mathcal{O}_2/\mathcal{O}_1}$ は前層
$U \mapsto \Omega_{\mathcal{O}_2(U)/\mathcal{O}_1(U)}$ に付随する層である。
\end{lemma}

\begin{proof}
写像 (\ref{modules-equation-define-module-differentials}) を考える。同様の前層の
写像があり，開集合 $U$ における値は
$$
\mathcal{O}_2(U)[\mathcal{O}_2(U) \times \mathcal{O}_2(U)] \oplus
\mathcal{O}_2(U)[\mathcal{O}_2(U) \times \mathcal{O}_2(U)] \oplus
\mathcal{O}_2(U)[\mathcal{O}_1(U)]
\longrightarrow
\mathcal{O}_2(U)[\mathcal{O}_2(U)]
$$
である。《代数》定義 \ref{algebra-definition-differentials} における
微分加群の構成により，この写像の余核の $U$ における値は
$\Omega_{\mathcal{O}_2(U)/\mathcal{O}_1(U)}$ である。一方，
(\ref{modules-equation-define-module-differentials}) に現れる層は，上の前層の
層化である。層化は完全なので，結果が従う。
\end{proof}

\begin{lemma}
\label{modules-lemma-localize-differentials}
$X$ を位相空間とし，
$\varphi : \mathcal{O}_1 \to \mathcal{O}_2$ を環の層の準同型とする。
開集合 $U \subset X$ に対し，普遍導分と整合する標準同型
$$
\Omega_{\mathcal{O}_2/\mathcal{O}_1}|_U =
\Omega_{(\mathcal{O}_2|_U)/(\mathcal{O}_1|_U)}
$$
がある。
\end{lemma}

\begin{proof}
$\Omega_{\mathcal{O}_2/\mathcal{O}_1}$ は写像
(\ref{modules-equation-define-module-differentials}) の余核なので成り立つ。
\end{proof}

\begin{lemma}
\label{modules-lemma-pullback-differentials}
$f : Y \to X$ を位相空間の連続写像とし，
$\varphi : \mathcal{O}_1 \to \mathcal{O}_2$ を $X$ 上の環の層の
準同型とする。このとき，普遍導分と整合する標準同一視
$f^{-1}\Omega_{\mathcal{O}_2/\mathcal{O}_1} =
\Omega_{f^{-1}\mathcal{O}_2/f^{-1}\mathcal{O}_1}$
がある。
\end{lemma}

\begin{proof}
$\Omega_{\mathcal{O}_2/\mathcal{O}_1}$ は写像
(\ref{modules-equation-define-module-differentials}) の余核であり，
$\Omega_{f^{-1}\mathcal{O}_2/f^{-1}\mathcal{O}_1}$
についても同様だからである。これは，関手 $f^{-1}$ が完全であり，かつ
$f^{-1}(\mathcal{O}_2[\mathcal{O}_2]) =
f^{-1}\mathcal{O}_2[f^{-1}\mathcal{O}_2]$，
$f^{-1}(\mathcal{O}_2[\mathcal{O}_2 \times \mathcal{O}_2]) =
f^{-1}\mathcal{O}_2[f^{-1}\mathcal{O}_2 \times f^{-1}\mathcal{O}_2]$，および
$f^{-1}(\mathcal{O}_2[\mathcal{O}_1]) =
f^{-1}\mathcal{O}_2[f^{-1}\mathcal{O}_1]$
であることから従う。
\end{proof}

\begin{lemma}
\label{modules-lemma-stalk-module-differentials}
$X$ を位相空間とし，
$\mathcal{O}_1 \to \mathcal{O}_2$ を $X$ 上の環の層の準同型とする。
$x \in X$ に対し，
$\Omega_{\mathcal{O}_2/\mathcal{O}_1, x} =
\Omega_{\mathcal{O}_{2, x}/\mathcal{O}_{1, x}}$
が成り立つ。
\end{lemma}

\begin{proof}
これは包含写像 $\{x\} \to X$ に対する補題
\ref{modules-lemma-pullback-differentials} の特別な場合である。別の証明として，
補題 \ref{modules-lemma-differentials-sheafify}，《層》補題
\ref{sheaves-lemma-stalk-sheafification}，および《代数》補題
\ref{algebra-lemma-colimit-differentials} を用いることもできる。
\end{proof}

\begin{lemma}
\label{modules-lemma-functoriality-differentials}
$X$ を位相空間とする。
$$
\xymatrix{
\mathcal{O}_2 \ar[r]_\varphi & \mathcal{O}_2' \\
\mathcal{O}_1 \ar[r] \ar[u] & \mathcal{O}'_1 \ar[u]
}
$$
を $X$ 上の環の層の可換図式とする。写像
$\mathcal{O}_2 \to \mathcal{O}'_2$ と写像
$\text{d} : \mathcal{O}'_2 \to \Omega_{\mathcal{O}'_2/\mathcal{O}'_1}$
の合成は $\mathcal{O}_1$-導分である。したがって，
$\mathcal{O}_2$-加群の層の標準写像
$\Omega_{\mathcal{O}_2/\mathcal{O}_1} \to
\Omega_{\mathcal{O}'_2/\mathcal{O}'_1}$
を得る。これは，$\mathcal{O}_2$ の任意の局所切断 $f$ に対して
$\text{d}(f) \mapsto \text{d}(\varphi(f))$
となるという性質によって一意に特徴づけられる。このようにして
$\Omega_{-/-}$ は環の層の射の圏上の関手となる。
\end{lemma}

\begin{proof}
定義から直ちに従う。
\end{proof}

\begin{lemma}
\label{modules-lemma-differential-seq}
補題 \ref{modules-lemma-functoriality-differentials} において，
$\mathcal{O}_2 \to \mathcal{O}'_2$ は全射で，核が
$\mathcal{I} \subset \mathcal{O}_2$ であり，
$\mathcal{O}_1 = \mathcal{O}'_1$ であると仮定する。このとき，
$\mathcal{O}'_2$-加群の層の標準完全列
$$
\mathcal{I}/\mathcal{I}^2
\longrightarrow
\Omega_{\mathcal{O}_2/\mathcal{O}_1} \otimes_{\mathcal{O}_2} \mathcal{O}'_2
\longrightarrow
\Omega_{\mathcal{O}'_2/\mathcal{O}_1}
\longrightarrow
0
$$
がある。最左の写像は，$\mathcal{I}$ の局所切断 $f$ を
$\text{d}f \otimes 1$ へ写すという規則によって特徴づけられる。
\end{lemma}

\begin{proof}
$\mathcal{I}$ の局所切断 $f$ に対し，$\mathcal{I}/\mathcal{I}^2$ における
$f$ の像を $\overline f$ と書く。写像
$\overline f \mapsto \text{d}f \otimes 1$ が良定義であることを示すには，
$f_1,f_2$ が $\mathcal{I}$ の局所切断ならば
$\text{d}f_1f_2 \otimes 1 = 0$ であることを確認すればよい。これは
Leibniz 則
$\text{d} f_1f_2 \otimes 1 =
(f_1 \text{d}f_2 + f_2 \text{d} f_1 )\otimes 1 =
\text{d}f_2 \otimes f_1 + \text{d}f_1 \otimes f_2 = 0$
から明らかである。同様の計算により，この写像は
$\mathcal{O}'_2 = \mathcal{O}_2/\mathcal{I}$-線形である。右側の写像は
補題 \ref{modules-lemma-functoriality-differentials} のものである。列の完全性は，
茎上で確認できる（補題 \ref{modules-lemma-abelian}）。補題
\ref{modules-lemma-stalk-module-differentials} により，これは《代数》補題
\ref{algebra-lemma-differential-seq} から従う。
\end{proof}

\begin{definition}
\label{modules-definition-differentials}
$(f,f^\sharp) : (X,\mathcal{O}_X) \to (S,\mathcal{O}_S)$ を
環付き空間の射とする。
\begin{enumerate}
\item $\mathcal{F}$ を $\mathcal{O}_X$-加群の層とする。$\mathcal{F}$ への
{\it $S$-導分}とは $f^{-1}\mathcal{O}_S$-導分，より正確には定義
\ref{modules-definition-derivation} の意味での $f^\sharp$-導分である。
$\mathcal{F}$ への $S$-導分全体の集合を
$\text{Der}_S(\mathcal{O}_X,\mathcal{F})$ と書く。
\item {\it $S$ 上の $X$ の微分の層 $\Omega_{X/S}$}とは，普遍
$S$-導分 $\text{d}_{X/S} : \mathcal{O}_X \to \Omega_{X/S}$ を備えた
微分加群の層 $\Omega_{\mathcal{O}_X/f^{-1}\mathcal{O}_S}$ である。
\end{enumerate}
\end{definition}

\noindent
導分が自然に現れる特別な状況を次に述べる。

\begin{lemma}
\label{modules-lemma-double-structure-gives-derivation}
$(f,f^\sharp) : (X,\mathcal{O}_X) \to (S,\mathcal{O}_S)$ を
環付き空間の射とする。短完全列
$$
0 \to \mathcal{I} \to \mathcal{A} \to \mathcal{O}_X \to 0
$$
を考える。ここで $\mathcal{A}$ は $f^{-1}\mathcal{O}_S$-代数の層，
$\pi : \mathcal{A} \to \mathcal{O}_X$ は
$f^{-1}\mathcal{O}_S$-代数の層の全射であり，
$\mathcal{I} = \Ker(\pi)$ はその核である。$\mathcal{I}$ は
$\mathcal{A}$ において平方零のイデアルの層であると仮定する。したがって
$\mathcal{I}$ は自然な $\mathcal{O}_X$-加群の層の構造をもつ。
$\pi$ の切断 $s : \mathcal{O}_X \to \mathcal{A}$ とは，
$\pi \circ s = \text{id}$ を満たす $f^{-1}\mathcal{O}_S$-代数の層の射である。
$\pi$ の任意の切断 $s : \mathcal{O}_X \to \mathcal{A}$ と任意の
$S$-導分 $D : \mathcal{O}_X \to \mathcal{I}$ に対し，写像
$$
s + D : \mathcal{O}_X \to \mathcal{A}
$$
は $\pi$ の切断である。また，すべての切断 $s'$ は，一意な $S$-導分
$D$ により $s + D$ の形に書ける。
\end{lemma}

\begin{proof}
$\mathcal{I}$ 上の $\mathcal{O}_X$-加群の層の構造は
$h\tau=\widetilde h\tau$（$\mathcal{A}$ における積）で与えられることを
思い出そう。ここで $h$ は $\mathcal{O}_X$ の局所切断，
$\widetilde h$ は $h$ の $\mathcal{A}$ の局所切断への局所持ち上げ，
$\tau$ は $\mathcal{I}$ の局所切断である。特に $s$ が与えられれば，
$\widetilde h=s(h)$ としてよい。$s+D$ が環の層の準同型であることを
確かめるため，次を計算する。
\begin{eqnarray*}
(s + D)(ab) & = & s(ab) + D(ab) \\
& = & s(a)s(b) + aD(b) + D(a)b \\
& = & s(a) s(b) + s(a)D(b) + D(a)s(b) \\
& = & (s(a) + D(a))(s(b) + D(b))
\end{eqnarray*}
最後の等式では Leibniz 則を用いた。同様に，$D$ は $S$-導分なので
$s+D$ は $f^{-1}\mathcal{O}_S$-代数の層の射であることが分かる。逆に，
$s'$ が与えられたとき $D=s'-s$ と置く。詳細は省略する。
\end{proof}

\begin{lemma}
\label{modules-lemma-functoriality-differentials-ringed-spaces}
$$
\xymatrix{
X' \ar[d]_{h'} \ar[r]_f & X \ar[d]^h \\
S' \ar[r]^g & S
}
$$
を環付き空間の可換図式とする。
\begin{enumerate}
\item 標準写像 $\mathcal{O}_X \to f_*\mathcal{O}_{X'}$ と
$f_*\text{d}_{X'/S'} : f_*\mathcal{O}_{X'} \to f_*\Omega_{X'/S'}$
の合成は $S$-導分であり，$\mathcal{O}_X$-加群の層の標準写像
$\Omega_{X/S} \to f_*\Omega_{X'/S'}$
を得る。
\item 可換図式
$$
\xymatrix{
f^{-1}\mathcal{O}_X \ar[r] & \mathcal{O}_{X'} \\
f^{-1}h^{-1}\mathcal{O}_S \ar[u] \ar[r] & (h')^{-1}\mathcal{O}_{S'} \ar[u]
}
$$
は補題 \ref{modules-lemma-pullback-differentials} および
\ref{modules-lemma-functoriality-differentials} により，標準写像
$f^{-1}\Omega_{X/S} \to \Omega_{X'/S'}$ を誘導する。
\end{enumerate}
これら二つの写像は，$f^*$ と $f_*$ の随伴，等式
$f^*\Omega_{X/S} =
f^{-1}\Omega_{X/S} \otimes_{f^{-1}\mathcal{O}_X} \mathcal{O}_{X'}$，
および《層》補題 \ref{sheaves-lemma-adjointness-tensor-restrict} を通じて，
同じ $\mathcal{O}_{X'}$-加群の層の準同型
$$
c_f : f^*\Omega_{X/S} \longrightarrow \Omega_{X'/S'}
$$
に対応する。これは，$\mathcal{O}_X$ の任意の局所切断 $a$ に対し，
$f^*\text{d}_{X/S}(a)$ を $\text{d}_{X'/S'}(f^*a)$ へ写すという性質で
一意に特徴づけられる。
\end{lemma}

\begin{proof}
省略する。
\end{proof}

\begin{lemma}
\label{modules-lemma-check-functoriality-differentials}
$$
\xymatrix{
X'' \ar[d] \ar[r]_g & X' \ar[d] \ar[r]_f & X \ar[d] \\
S'' \ar[r] & S' \ar[r] & S
}
$$
を環付き空間の可換図式とする。補題
\ref{modules-lemma-functoriality-differentials-ringed-spaces} の記法のもとで，
$$
c_{f \circ g} = c_g \circ g^* c_f
$$
が写像 $(f \circ g)^*\Omega_{X/S} \to \Omega_{X''/S''}$ として成り立つ。
\end{lemma}

\begin{proof}
省略する。
\end{proof}

\begin{lemma}
\label{modules-lemma-triangle-differentials}
$f : X \to Y$，$g : Y \to S$ を環付き空間の射とする。
このとき標準完全列
$$
f^*\Omega_{Y/S} \to \Omega_{X/S} \to \Omega_{X/Y} \to 0
$$
がある。ここで各写像は補題
\ref{modules-lemma-functoriality-differentials-ringed-spaces} を適用して得られる。
\end{lemma}

\begin{proof}
$x \in X$ における茎上の誘導写像をとり，補題
\ref{modules-lemma-stalk-module-differentials} を用いると，列
$$
\mathcal{O}_{X, x} \otimes_{\mathcal{O}_{Y, f(x)}}
\Omega_{\mathcal{O}_{Y, f(x)}/\mathcal{O}_{S, g(f(x))}} \to
\Omega_{\mathcal{O}_{X, x}/\mathcal{O}_{S, g(f(x))}} \to
\Omega_{\mathcal{O}_{X, x}/\mathcal{O}_{Y, f(x)}} \to 0
$$
を得る。この列の写像が《代数》補題
\ref{algebra-lemma-exact-sequence-differentials} の写像と同じであることを
確認すれば十分である。これは補題
\ref{modules-lemma-functoriality-differentials-ringed-spaces} における写像の特徴づけと，
《層》補題 \ref{sheaves-lemma-stalk-pullback-modules} の同型を通じて
$\mathcal{O}_Y$-加群の層の任意の局所切断 $s$ に対し
$(f^*s)_x = s_x \otimes 1$ となることから従う。
\end{proof}


















\section{有限階微分作用素}
\label{modules-section-differential-operators}

\noindent
本節では有限階微分作用素を導入する。読者には，可換代数の章の対応する節
（《代数》\ref{algebra-section-differential-operators} 節）も
参照することを勧める。

\begin{definition}
\label{modules-definition-differential-operators}
$X$ を位相空間とし，
$\varphi : \mathcal{O}_1 \to \mathcal{O}_2$ を $X$ 上の環の層の
準同型とする。$k \geq 0$ を整数とし，$\mathcal{F},\mathcal{G}$ を
$\mathcal{O}_2$-加群の層とする。{\it $k$ 階微分作用素
$D : \mathcal{F} \to \mathcal{G}$}とは，$\mathcal{O}_1$-線形写像で，
$\mathcal{O}_2$ のすべての局所切断 $g$ に対して写像
$s \mapsto D(gs)-gD(s)$ が $k-1$ 階微分作用素となるものである。
基底の場合 $k=0$ については，$0$ 階微分作用素を
$\mathcal{O}_2$-線形写像と定義する。
\end{definition}

\noindent
$D : \mathcal{F} \to \mathcal{G}$ が $k$ 階微分作用素ならば，
$\mathcal{O}_2$ のすべての局所切断 $g$ に対し $gD$ も
$k$ 階微分作用素である。二つの $k$ 階微分作用素の和も同じ階の
微分作用素である。したがって，これら全体の集合
$$
\text{Diff}^k(\mathcal{F}, \mathcal{G}) =
\text{Diff}^k_{\mathcal{O}_2/\mathcal{O}_1}(\mathcal{F}, \mathcal{G})
$$
は $\Gamma(X,\mathcal{O}_2)$-加群である。また，
$$
\text{Diff}^0(\mathcal{F}, \mathcal{G}) \subset
\text{Diff}^1(\mathcal{F}, \mathcal{G}) \subset
\text{Diff}^2(\mathcal{F}, \mathcal{G}) \subset \ldots
$$
である。開集合 $U \subset X$ を，$k$ 階微分作用素
$D : \mathcal{F}|_U \to \mathcal{G}|_U$ の加群へ写す規則は，
$X$ 上の $\mathcal{O}_2$-加群の層である。したがって微分作用素の層を得る
（これが必要になれば，ここに定義を追加する）。

\begin{lemma}
\label{modules-lemma-composition-differential-operators}
$X$ を位相空間とし，$\mathcal{O}_1 \to \mathcal{O}_2$ を
$X$ 上の環の層の写像とする。$\mathcal{E},\mathcal{F},\mathcal{G}$ を
$\mathcal{O}_2$-加群の層とする。
$D : \mathcal{E} \to \mathcal{F}$ と
$D' : \mathcal{F} \to \mathcal{G}$ がそれぞれ $k$ 階および $k'$ 階の
微分作用素ならば，$D' \circ D$ は $k+k'$ 階微分作用素である。
\end{lemma}

\begin{proof}
$g$ を $\mathcal{O}_2$ の局所切断とする。このとき
$\mathcal{E}$ の局所切断 $x$ を
$$
D'(D(gx)) - gD'(D(x)) = D'(D(gx)) - D'(gD(x)) + D'(gD(x)) - gD'(D(x))
$$
へ写す写像は，より低い階の微分作用素の二つの合成の和である。
したがって，$k+k'$ に関する帰納法により補題が従う。
\end{proof}

\begin{lemma}
\label{modules-lemma-module-principal-parts}
$X$ を位相空間とし，$\mathcal{O}_1 \to \mathcal{O}_2$ を
$X$ 上の環の層の写像とする。$\mathcal{F}$ を
$\mathcal{O}_2$-加群の層とし，$k \geq 0$ とする。
$\mathcal{O}_2$-加群の層
$\mathcal{P}^k_{\mathcal{O}_2/\mathcal{O}_1}(\mathcal{F})$
と，$\mathcal{O}_2$-加群の層 $\mathcal{G}$ に関して関手的な標準同型
$$
\text{Diff}^k_{\mathcal{O}_2/\mathcal{O}_1}(\mathcal{F}, \mathcal{G}) =
\Hom_{\mathcal{O}_2}(
\mathcal{P}^k_{\mathcal{O}_2/\mathcal{O}_1}(\mathcal{F}), \mathcal{G})
$$
が存在する。
\end{lemma}

\begin{proof}
存在は一般的な圏論的議論から従う（将来ここに参照を挿入する）が，
この構成は後の証明で有用なので，直接的な構成も与える。
補題 \ref{modules-lemma-universal-module} の証明で導入した記法を自由に用いる。
任意の微分作用素 $D : \mathcal{F} \to \mathcal{G}$ に対し，
$[m]$ を $D(m)$ へ写す $\mathcal{O}_2$-線形写像
$L_D : \mathcal{O}_2[\mathcal{F}] \to \mathcal{G}$
を得る。$D$ が $0$ 階ならば，$L_D$ は局所切断
$$
[m + m'] - [m] - [m'],\quad
g_0[m] - [g_0m]
$$
を零化する。ここで $g_0$ は $\mathcal{O}_2$ の局所切断，
$m,m'$ は $\mathcal{F}$ の局所切断である。$D$ が $1$ 階ならば，
$L_D$ は局所切断
$$
[m + m' - [m] - [m'],\quad
f[m] - [fm], \quad
g_0g_1[m] - g_0[g_1m] - g_1[g_0m] + [g_1g_0m]
$$
を零化する。ここで $f$ は $\mathcal{O}_1$ の局所切断，
$g_0,g_1$ は $\mathcal{O}_2$ の局所切断，$m,m'$ は
$\mathcal{F}$ の局所切断である。$D$ が $k$ 階ならば，$L_D$ は
局所切断 $[m + m']-[m]-[m']$，$f[m]-[fm]$，および局所切断
$$
g_0g_1\ldots g_k[m] - \sum g_0 \ldots \hat g_i \ldots g_k[g_im] + \ldots
+(-1)^{k + 1}[g_0\ldots g_km]
$$
を零化する。逆に，$L : \mathcal{O}_2[\mathcal{F}] \to \mathcal{G}$ が
前文に列挙したすべての局所切断を零化する
$\mathcal{O}_2$-線形写像ならば，$m \mapsto L([m])$ は $k$ 階微分作用素で
ある。したがって，
$\mathcal{P}^k_{\mathcal{O}_2/\mathcal{O}_1}(\mathcal{F})$
は，$\mathcal{O}_2[\mathcal{F}]$ をこれらの局所切断が生成する
$\mathcal{O}_2$-部分加群の層で割った商である。
\end{proof}

\begin{definition}
\label{modules-definition-module-principal-parts}
$X$ を位相空間とし，$\mathcal{O}_1 \to \mathcal{O}_2$ を
$X$ 上の環の層の写像とする。$\mathcal{F}$ を
$\mathcal{O}_2$-加群の層とする。補題
\ref{modules-lemma-module-principal-parts} で構成した加群の層
$\mathcal{P}^k_{\mathcal{O}_2/\mathcal{O}_1}(\mathcal{F})$ を，
$\mathcal{F}$ の {\it $k$ 階主部加群の層}という。
\end{definition}

\noindent
包含
$$
\text{Diff}^0(\mathcal{F}, \mathcal{G}) \subset
\text{Diff}^1(\mathcal{F}, \mathcal{G}) \subset
\text{Diff}^2(\mathcal{F}, \mathcal{G}) \subset \ldots
$$
は，Yoneda の補題（《圏》補題 \ref{categories-lemma-yoneda}）を通じて全射
$$
\ldots \to \mathcal{P}^2_{\mathcal{O}_2/\mathcal{O}_1}(\mathcal{F})
\to \mathcal{P}^1_{\mathcal{O}_2/\mathcal{O}_1}(\mathcal{F})
\to \mathcal{P}^0_{\mathcal{O}_2/\mathcal{O}_1}(\mathcal{F}) = \mathcal{F}
$$
に対応することに注意する。

\begin{lemma}
\label{modules-lemma-differential-operators-sheafify}
$X$ を位相空間とし，$\mathcal{O}_1 \to \mathcal{O}_2$ を
$X$ 上の環の前層の準同型とする。$\mathcal{F}$ を
$\mathcal{O}_2$-加群の前層とする。このとき
$\mathcal{P}^k_{\mathcal{O}_2^\#/\mathcal{O}_1^\#}(\mathcal{F}^\#)$
は前層
$U \mapsto P^k_{\mathcal{O}_2(U)/\mathcal{O}_1(U)}(\mathcal{F}(U))$
に付随する層である。
\end{lemma}

\begin{proof}
補題 \ref{modules-lemma-differentials-sheafify} で微分の層について行ったのと
まったく同じ方法で証明できる。あるいは，補題
\ref{modules-lemma-module-principal-parts} の普遍性を直接用いて等式を示す方が
自然かもしれない。詳細は省略する。
\end{proof}

\begin{lemma}
\label{modules-lemma-sequence-of-principal-parts}
$X$ を位相空間とし，$\mathcal{O}_1 \to \mathcal{O}_2$ を
$X$ 上の環の層の準同型とする。$\mathcal{F}$ を
$\mathcal{O}_2$-加群の層とする。$\mathcal{F}$ に関して関手的な標準短完全列
$$
0 \to
\Omega_{\mathcal{O}_2/\mathcal{O}_1} \otimes_{\mathcal{O}_2} \mathcal{F} \to
\mathcal{P}^1_{\mathcal{O}_2/\mathcal{O}_1}(\mathcal{F}) \to
\mathcal{F} \to 0
$$
があり，{\it 主部列}と呼ぶ。
\end{lemma}

\begin{proof}
可換代数における版（《代数》補題
\ref{algebra-lemma-sequence-of-principal-parts}）と，補題
\ref{modules-lemma-differentials-sheafify} および
\ref{modules-lemma-differential-operators-sheafify} から従う。
\end{proof}

\begin{remark}
\label{modules-remark-functoriality-principal-parts}
$X$ を位相空間とする。$X$ 上の環の層の可換図式
$$
\xymatrix{
\mathcal{B} \ar[r] & \mathcal{B}' \\
\mathcal{A} \ar[u] \ar[r] & \mathcal{A}' \ar[u]
}
$$
と，$\mathcal{B}$-加群の層 $\mathcal{F}$，
$\mathcal{B}'$-加群の層 $\mathcal{F}'$，および
$\mathcal{B}$-線形写像 $\mathcal{F} \to \mathcal{F}'$ が
与えられたとする。このとき，整合的な加群の層の写像の系
$$
\xymatrix{
\ldots \ar[r] &
\mathcal{P}^2_{\mathcal{B}'/\mathcal{A}'}(\mathcal{F}') \ar[r] &
\mathcal{P}^1_{\mathcal{B}'/\mathcal{A}'}(\mathcal{F}') \ar[r] &
\mathcal{P}^0_{\mathcal{B}'/\mathcal{A}'}(\mathcal{F}') \\
\ldots \ar[r] &
\mathcal{P}^2_{\mathcal{B}/\mathcal{A}}(\mathcal{F}) \ar[r] \ar[u] &
\mathcal{P}^1_{\mathcal{B}/\mathcal{A}}(\mathcal{F}) \ar[r] \ar[u] &
\mathcal{P}^0_{\mathcal{B}/\mathcal{A}}(\mathcal{F}) \ar[u]
}
$$
を得る。これらの写像は，この種の写像のさらなる合成と整合する。
これを見る最も簡単な方法は，補題 \ref{modules-lemma-module-principal-parts} の
証明における $\mathcal{P}^k_{\mathcal{B}/\mathcal{A}}(\mathcal{M})$ の
（局所）生成元と関係式による記述を用いることであるが，これらの加群の層の
普遍性から直接見ることもできる。さらに，これらの写像は補題
\ref{modules-lemma-sequence-of-principal-parts} の短完全列と整合する。
\end{remark}

\noindent
次に，定義を環付き空間の射へ拡張する。

\begin{definition}
\label{modules-definition-relative-differential-operators}
$(f,f^\sharp) : (X,\mathcal{O}_X) \to (S,\mathcal{O}_S)$ を
環付き空間の射とする。$\mathcal{F}$ と $\mathcal{G}$ を
$\mathcal{O}_X$-加群の層とし，$k \geq 0$ を整数とする。
{\it $X/S$ 上の $k$ 階微分作用素}とは，
$f^\sharp : f^{-1}\mathcal{O}_S \to \mathcal{O}_X$ に関する微分作用素
$D : \mathcal{F} \to \mathcal{G}$ のことである。これらの微分作用素全体の
集合を $\text{Diff}^k_{X/S}(\mathcal{F},\mathcal{G})$ と書く。
\end{definition}












\section{de Rham 複体}
\label{modules-section-de-rham-complex}

\noindent
本節は，《代数》\ref{algebra-section-de-rham-complex} 節の環付き空間の射に
対する類似である。読者には，まずその節を読むことを強く勧める。

\medskip\noindent
$X$ を位相空間とし，$\mathcal{A} \to \mathcal{B}$ を
環の層の準同型とする。
$\text{d} : \mathcal{B} \to \Omega_{\mathcal{B}/\mathcal{A}}$
を，\ref{modules-section-differentials} 節で構成した微分加群の層とその普遍
$\mathcal{A}$-導分とする。$i \geq 0$ に対し，
$$
\Omega_{\mathcal{B}/\mathcal{A}}^i =
\wedge^i_\mathcal{B}(\Omega_{\mathcal{B}/\mathcal{A}})
$$
を \ref{modules-section-symmetric-exterior} 節における第 $i$ 外冪とする。

\begin{definition}
\label{modules-definition-de-rham-complex}
上の状況で，{\it $\mathcal{A}$ 上の $\mathcal{B}$ の de Rham 複体}とは，
$0$ 次の微分が
$\text{d} : \mathcal{B} \to \Omega_{\mathcal{B}/\mathcal{A}}$ で与えられ，
高次の微分が次の性質をもつ，$\mathcal{A}$-加群の層の一意な複体
$$
\Omega_{\mathcal{B}/\mathcal{A}}^0 \to
\Omega_{\mathcal{B}/\mathcal{A}}^1 \to
\Omega_{\mathcal{B}/\mathcal{A}}^2 \to \ldots
$$
のことである。
\begin{equation}
\label{modules-equation-rule}
\text{d}\left(b_0\text{d}b_1 \wedge \ldots \wedge \text{d}b_p\right) =
\text{d}b_0 \wedge \text{d}b_1 \wedge \ldots \wedge \text{d}b_p
\end{equation}
ここで $b_0,\ldots,b_p \in \mathcal{B}(U)$ は共通の開集合
$U \subset X$ 上の切断である。
\end{definition}

\noindent
《代数》\ref{algebra-section-de-rham-complex} 節で与えた煩雑な議論を
繰り返してこの複体を構成することもできる。その代わり，
$\Omega_{\mathcal{B}/\mathcal{A}}$ は前層
$U \mapsto \Omega_{\mathcal{B}(U)/\mathcal{A}(U)}$ の層化であることを
思い出そう（補題 \ref{modules-lemma-differentials-sheafify} 参照）。
したがって，$\Omega_{\mathcal{B}/\mathcal{A}}^i$ は前層
$U \mapsto \Omega^i_{\mathcal{B}(U)/\mathcal{A}(U)}$ の層化である
（補題 \ref{modules-lemma-local-tensor-algebra} 参照）。ゆえに de Rham 複体を，
規則
$$
U \longmapsto
\Omega^\bullet_{\mathcal{B}(U)/\mathcal{A}(U)}
$$
の層化として定義できる。

\begin{lemma}
\label{modules-lemma-pullback-de-rham-complex}
$f : Y \to X$ を位相空間の連続写像とし，
$\mathcal{A} \to \mathcal{B}$ を $X$ 上の環の層の準同型とする。
このとき de Rham 複体の標準同一視
$f^{-1}\Omega^\bullet_{\mathcal{B}/\mathcal{A}} =
\Omega^\bullet_{f^{-1}\mathcal{B}/f^{-1}\mathcal{A}}$
がある。
\end{lemma}

\begin{proof}
省略する。ヒント：補題 \ref{modules-lemma-pullback-differentials} と比較せよ。
\end{proof}

\begin{lemma}
\label{modules-lemma-differentials-de-rham-complex-order-1}
$X$ を位相空間とし，$\mathcal{A} \to \mathcal{B}$ を
$X$ 上の環の層の準同型とする。微分
$\text{d} : \Omega^i_{\mathcal{B}/\mathcal{A}}  \to
\Omega^{i + 1}_{\mathcal{B}/\mathcal{A}}$
は $1$ 階微分作用素である。
\end{lemma}

\begin{proof}
上で de Rham 複体を規則
$U \mapsto \Omega^\bullet_{\mathcal{B}(U)/\mathcal{A}(U)}$ の層化として
構成したことにより，これは《代数》補題
\ref{algebra-lemma-differentials-de-rham-complex-order-1} から従う。
\end{proof}

\noindent
$X$ を位相空間とし，
$$
\xymatrix{
\mathcal{B} \ar[r] & \mathcal{B}' \\
\mathcal{A} \ar[r] \ar[u] & \mathcal{A}' \ar[u]
}
$$
を $X$ 上の環の層の可換図式とする。de Rham 複体の自然な写像
$$
\Omega^\bullet_{\mathcal{B}/\mathcal{A}} \longrightarrow
\Omega^\bullet_{\mathcal{B}'/\mathcal{A}'}
$$
がある。すなわち，$0$ 次では写像 $\mathcal{B} \to \mathcal{B}'$，
$1$ 次では \ref{modules-section-differentials} 節で構成した写像
$\Omega_{\mathcal{B}/\mathcal{A}} \to \Omega_{\mathcal{B}'/\mathcal{A}'}$
であり，$p \geq 2$ に対しては誘導された写像
$\Omega^p_{\mathcal{B}/\mathcal{A}} =
\wedge^p_\mathcal{B}(\Omega_{\mathcal{B}/\mathcal{A}}) \to
\wedge^p_{\mathcal{B}'}(\Omega_{\mathcal{B}'/\mathcal{A}'}) =
\Omega^p_{\mathcal{B}'/\mathcal{A}'}$
である。微分との整合性は，式 (\ref{modules-equation-rule}) による微分の
特徴づけから従う。

\begin{definition}
\label{modules-definition-de-rham-complex-morphism-ringed-spaces}
$f : (X,\mathcal{O}_X) \to (Y,\mathcal{O}_Y)$ を環付き空間の射とする。
$f$ の，または $Y$ 上の $X$ の {\it de Rham 複体}とは，複体
$$
\Omega^\bullet_{X/Y} = \Omega^\bullet_{\mathcal{O}_X/f^{-1}\mathcal{O}_Y}
$$
のことである。
\end{definition}

\noindent
環付き空間の可換図式
$$
\xymatrix{
X' \ar[d]_{h'} \ar[r]_f & X \ar[d]^h \\
S' \ar[r]^g & S
}
$$
を考える。このとき de Rham 複体の標準写像
$$
\Omega^\bullet_{X/S} \to f_*\Omega^\bullet_{X'/S'}
$$
を得る。実際，$X'$ 上の環の層の可換図式
$$
\xymatrix{
f^{-1}\mathcal{O}_X \ar[r] & \mathcal{O}_{X'} \\
f^{-1}h^{-1}\mathcal{O}_S \ar[u] \ar[r] & (h')^{-1}\mathcal{O}_{S'} \ar[u]
}
$$
は，複体の写像（上を参照）
$$
f^{-1}\Omega^\bullet_{X/S} =
\Omega^\bullet_{f^{-1}\mathcal{O}_X/f^{-1}h^{-1}\mathcal{O}_S}
\longrightarrow
\Omega^\bullet_{\mathcal{O}_{X'}/(h')^{-1}\mathcal{O}_{S'}} =
\Omega^\bullet_{X'/S'}
$$
を与える（最初の等式には補題
\ref{modules-lemma-pullback-de-rham-complex} を用いる）。その後，随伴を用いればよい。

\begin{lemma}
\label{modules-lemma-differentials-relative-de-rham-complex-order-1}
$f : X \to Y$ を環付き空間の射とする。微分
$\text{d} : \Omega^i_{X/Y}  \to \Omega^{i + 1}_{X/Y}$
は $X/Y$ 上の $1$ 階微分作用素である。
\end{lemma}

\begin{proof}
補題 \ref{modules-lemma-differentials-de-rham-complex-order-1} と定義から直ちに従う。
\end{proof}





















\section{素朴な余接複体}
\label{modules-section-netherlander}

\noindent
本節は，《代数》\ref{algebra-section-netherlander} 節の環付き空間の射に
対する類似である。読者には，まずその節を読むことを強く勧める。

\medskip\noindent
$X$ を位相空間とし，$\mathcal{A} \to \mathcal{B}$ を
環の層の準同型とする。本節では，$X$ 上の任意の集合の層
$\mathcal{E}$ に対し，前層
$U \mapsto \mathcal{A}(U)[\mathcal{E}(U)]$ の層化を
$\mathcal{A}[\mathcal{E}]$ と書く。ここで
$\mathcal{A}(U)[\mathcal{E}(U)]$
は，変数が $\mathcal{E}(U)$ の元に対応する
$\mathcal{A}(U)$ 上の多項式代数を表す。
$e \in \mathcal{E}(U)$ に対応する変数を
$[e] \in \mathcal{A}(U)[\mathcal{E}(U)]$ と書く。
$\mathcal{A}$-代数の層の標準全射
\begin{equation}
\label{modules-equation-canonical-presentation}
\mathcal{A}[\mathcal{B}] \longrightarrow \mathcal{B},\quad [b] \longmapsto b
\end{equation}
があり，その核を $\mathcal{I} \subset \mathcal{A}[\mathcal{B}]$ と書く。
$\mathcal{I}$ が局所切断 $[b][b']-[bb']$ および $[a]-a$ で生成される
ことは容易に分かる。補題 \ref{modules-lemma-differential-seq} により標準写像
\begin{equation}
\label{modules-equation-naive-cotangent-complex}
\mathcal{I}/\mathcal{I}^2
\longrightarrow
\Omega_{\mathcal{A}[\mathcal{B}]/\mathcal{A}}
\otimes_{\mathcal{A}[\mathcal{B}]} \mathcal{B}
\end{equation}
があり，その余核は $\Omega_{\mathcal{B}/\mathcal{A}}$ に標準的に同型である。

\begin{definition}
\label{modules-definition-naive-cotangent-complex}
$X$ を位相空間とし，$\mathcal{A} \to \mathcal{B}$ を
環の層の準同型とする。{\it 素朴な余接複体}
$\NL_{\mathcal{B}/\mathcal{A}}$ とは，鎖複体
(\ref{modules-equation-naive-cotangent-complex})
$$
\NL_{\mathcal{B}/\mathcal{A}} =
\left(\mathcal{I}/\mathcal{I}^2
\longrightarrow
\Omega_{\mathcal{A}[\mathcal{B}]/\mathcal{A}}
\otimes_{\mathcal{A}[\mathcal{B}]} \mathcal{B}\right)
$$
であり，$\mathcal{I}/\mathcal{I}^2$ を次数 $-1$ に，
$\Omega_{\mathcal{A}[\mathcal{B}]/\mathcal{A}}
\otimes_{\mathcal{A}[\mathcal{B}]} \mathcal{B}$
を次数 $0$ に置いたものである。
\end{definition}

\noindent
この構成は，微分加群の層について補題
\ref{modules-lemma-functoriality-differentials} で論じたものと同様の関手性を満たす。
すなわち，$X$ 上の環の層の可換図式
\begin{equation}
\label{modules-equation-commutative-square-sheaves}
\vcenter{
\xymatrix{
\mathcal{B} \ar[r] & \mathcal{B}' \\
\mathcal{A} \ar[u] \ar[r] & \mathcal{A}' \ar[u]
}
}
\end{equation}
が与えられると，複体の標準 $\mathcal{B}$-線形写像
$$
\NL_{\mathcal{B}/\mathcal{A}} \longrightarrow \NL_{\mathcal{B}'/\mathcal{A}'}
$$
がある。実際，可換図式の写像は標準写像
$\mathcal{A}[\mathcal{B}] \to \mathcal{A}'[\mathcal{B}']$
を与え，これは $\mathcal{I}$ を
$\mathcal{I}'=\Ker(\mathcal{A}'[\mathcal{B}']\to\mathcal{B}')$ へ写す。
したがって，写像
$\mathcal{I}/\mathcal{I}^2 \to \mathcal{I}'/(\mathcal{I}')^2$ と
微分加群の層の間の写像を得て，両者を合わせると素朴な余接複体の間の
所望の写像が得られる。この写像は次の意味で合成と整合する。
環の層の可換図式
$$
\xymatrix{
\mathcal{B} \ar[r] &
\mathcal{B}' \ar[r] &
\mathcal{B}'' \\
\mathcal{A} \ar[u] \ar[r] &
\mathcal{A}' \ar[u] \ar[r] &
\mathcal{A}'' \ar[u]
}
$$
が与えられたとき，合成
$$
\NL_{\mathcal{B}/\mathcal{A}} \longrightarrow
\NL_{\mathcal{B}'/\mathcal{A}'} \longrightarrow
\NL_{\mathcal{B}''/\mathcal{A}''}
$$
は外側の長方形に対する写像である。

\medskip\noindent
$\mathcal{B}$ を $\mathcal{A}$ 上の多項式代数の商として別の表示に
取り替えても，$D(\mathcal{B})$ の同じ対象を得る。このことを説明しよう。
$\mathcal{E}$ を $X$ 上の集合の層，
$\alpha : \mathcal{E} \to \mathcal{B}$ を集合の層の写像とする。
すると $\mathcal{A}$-代数の層の準同型
$\mathcal{A}[\mathcal{E}] \to \mathcal{B}$ を得る。この写像が全射ならば，
すなわち $\alpha(\mathcal{E})$ が $\mathcal{B}$ を
$\mathcal{A}$-代数の層として生成するならば，
$$
\NL(\alpha) = \left(
\mathcal{J}/\mathcal{J}^2
\longrightarrow
\Omega_{\mathcal{A}[\mathcal{E}]/\mathcal{A}}
\otimes_{\mathcal{A}[\mathcal{E}]} \mathcal{B}\right)
$$
と置く。ここで $\mathcal{J} \subset \mathcal{A}[\mathcal{E}]$ は全射
$\mathcal{A}[\mathcal{E}] \to \mathcal{B}$ の核である。次が成り立つ。

\begin{lemma}
\label{modules-lemma-NL-up-to-qis}
上の状況で，$D(\mathcal{B})$ における標準同型
$\NL(\alpha) = \NL_{\mathcal{B}/\mathcal{A}}$ がある。
\end{lemma}

\begin{proof}
$\NL_{\mathcal{B}/\mathcal{A}}=\NL(\text{id}_\mathcal{B})$ であることに
注意する。したがって，上のような二つの写像
$\alpha_i : \mathcal{E}_i \to \mathcal{B}$ が与えられたとき，
$D(\mathcal{B})$ における標準擬同型
$\NL(\alpha_1)=\NL(\alpha_2)$ があることを示せば十分である。
$\mathcal{E}=\mathcal{E}_1\amalg\mathcal{E}_2$，
$\alpha=\alpha_1\amalg\alpha_2 : \mathcal{E}\to\mathcal{B}$ と置く。また，
$\mathcal{J}_i = \Ker(\mathcal{A}[\mathcal{E}_i] \to \mathcal{B})$
および
$\mathcal{J} = \Ker(\mathcal{A}[\mathcal{E}] \to \mathcal{B})$.
と置く。$\mathcal{J}_i$ を $\mathcal{J}$ へ写す写像
$\mathcal{A}[\mathcal{E}_i] \to \mathcal{A}[\mathcal{E}]$ を得る。
したがって，複体の標準写像
$$
\NL(\alpha_i) \longrightarrow \NL(\alpha)
$$
を得るので，これらの写像が擬同型であることを示せば十分である。
これは茎上で確認すればよい（補題 \ref{modules-lemma-abelian}）。
$x \in X$ ならば，$\NL(\alpha)$ の茎は，写像
$\alpha_x : \mathcal{E}_x \to \mathcal{B}_x$ から来る表示
$\mathcal{A}_x[\mathcal{E}_x] \to \mathcal{B}_x$ に付随する，
《代数》\ref{algebra-section-netherlander} 節の複体 $\NL(\alpha_x)$ である。
（いくつかの詳細は省略する。微分を作る操作と茎をとる操作の整合性には補題
\ref{modules-lemma-stalk-module-differentials} を用いよ。）したがって《代数》補題
\ref{algebra-lemma-NL-homotopy} により結果が従う。
\end{proof}

\begin{lemma}
\label{modules-lemma-pullback-NL}
$f : X \to Y$ を位相空間の連続写像とし，
$\mathcal{A} \to \mathcal{B}$ を $Y$ 上の環の層の準同型とする。
このとき
$f^{-1}\NL_{\mathcal{B}/\mathcal{A}} =
\NL_{f^{-1}\mathcal{B}/f^{-1}\mathcal{A}}$ である。
\end{lemma}

\begin{proof}
省略する。ヒント：補題 \ref{modules-lemma-pullback-differentials} を用いよ。
\end{proof}

\begin{lemma}
\label{modules-lemma-stalk-NL}
$X$ を位相空間とし，$\mathcal{A} \to \mathcal{B}$ を
$X$ 上の環の層の準同型とする。$x \in X$ に対し，
$\NL_{\mathcal{B}/\mathcal{A},x} =
\NL_{\mathcal{B}_x/\mathcal{A}_x}$ である。
\end{lemma}

\begin{proof}
これは包含写像 $\{x\} \to X$ に対する補題
\ref{modules-lemma-pullback-NL} の特別な場合である。
\end{proof}

\begin{lemma}
\label{modules-lemma-exact-sequence-NL}
$X$ を位相空間とし，
$\mathcal{A} \to \mathcal{B} \to \mathcal{C}$
を環の層の写像とする。$C$ を複体の写像
$\NL_{\mathcal{C}/\mathcal{A}} \to \NL_{\mathcal{C}/\mathcal{B}}$ の錐
（《導来圏》定義 \ref{derived-definition-cone}）とする。
$\mathcal{C}$-加群の層の複体の標準写像
$$
c :
\NL_{\mathcal{B}/\mathcal{A}} \otimes_\mathcal{B} \mathcal{C}
\longrightarrow
C[-1]
$$
があり，これはコホモロジー層の標準六項完全列
$$
\xymatrix{
H^0(\NL_{\mathcal{B}/\mathcal{A}} \otimes_\mathcal{B} \mathcal{C}) \ar[r] &
H^0(\NL_{\mathcal{C}/\mathcal{A}}) \ar[r] &
H^0(\NL_{\mathcal{C}/\mathcal{B}}) \ar[r] &
0 \\
H^{-1}(\NL_{\mathcal{B}/\mathcal{A}} \otimes_\mathcal{B} \mathcal{C}) \ar[r] &
H^{-1}(\NL_{\mathcal{C}/\mathcal{A}}) \ar[r] &
H^{-1}(\NL_{\mathcal{C}/\mathcal{B}}) \ar[llu]
}
$$
を与える。
\end{lemma}

\begin{proof}
$c$ を与えるには，写像
$c_1 : \NL_{\mathcal{B}/\mathcal{A}} \otimes_\mathcal{B} \mathcal{C}
\to \NL_{\mathcal{C}/\mathcal{A}}$ と，合成
$$
\NL_{\mathcal{B}/\mathcal{A}} \otimes_\mathcal{B} \mathcal{C} \to
\NL_{\mathcal{C}/\mathcal{A}} \to \NL_{\mathcal{C}/\mathcal{B}}
$$
と零写像の間の明示的ホモトピーを与えなければならない
（《導来圏》補題 \ref{derived-lemma-map-from-cone} 参照）。
$c_1$ には，明らかな図式に対する上記の関手性を用いる。
ホモトピーには写像
$$
\NL_{\mathcal{B}/\mathcal{A}}^0 \otimes_\mathcal{B} \mathcal{C}
\longrightarrow
\NL_{\mathcal{C}/\mathcal{B}}^{-1},\quad
\text{d}[b] \otimes 1 \longmapsto [\varphi(b)] - b[1]
$$
を用いる。ここで $\varphi : \mathcal{B} \to \mathcal{C}$ は与えられた写像で
ある。《代数》注
\ref{algebra-remark-composition-homotopy-equivalent-to-zero} と比較されたい。
コホモロジー層に関する帰結を得るには，$H^0(c)$ が同型で，
$H^{-1}(c)$ が全射であることを示せば十分である。これは茎上で調べることが
でき（補題 \ref{modules-lemma-stalk-NL} 参照），可換代数における対応する結果
（《代数》補題 \ref{algebra-lemma-exact-sequence-NL}）を用いればよい。
いくつかの詳細は省略する。
\end{proof}

\noindent
環付き空間の射の余接複体を，上で定義した余接複体を用いて定義する。

\begin{definition}
\label{modules-definition-cotangent-complex-morphism-ringed-topoi}
環付き空間の射
$f : (X,\mathcal{O}_X) \to (Y,\mathcal{O}_Y)$ の
{\it 素朴な余接複体} $\NL_f=\NL_{X/Y}$ とは，
$\NL_{\mathcal{O}_X/f^{-1}\mathcal{O}_Y}$ のことである。
\end{definition}

\noindent
環付き空間の可換図式
$$
\xymatrix{
X' \ar[r]_g \ar[d]_{f'} & X \ar[d]^f \\
Y' \ar[r]^h & Y
}
$$
が与えられたとき，標準写像
$c : g^*\NL_{X/Y} \to \NL_{X'/Y'}$ がある。すなわち，これは写像
$$
g^*\NL_{X/Y} =
\mathcal{O}_{X'} \otimes_{g^{-1}\mathcal{O}_X}
\NL_{g^{-1}\mathcal{O}_X/g^{-1}f^{-1}\mathcal{O}_Y}
\longrightarrow
\NL_{\mathcal{O}_{X'}/(f')^{-1}\mathcal{O}_{Y'}} = \NL_{X'/Y'}
$$
であり，矢印は環の層の可換図式
$$
\xymatrix{
g^{-1}\mathcal{O}_X \ar[r]_{g^\sharp} &
\mathcal{O}_{X'} \\
g^{-1}f^{-1}\mathcal{O}_Y
\ar[r]^{g^{-1}h^\sharp}
\ar[u]^{g^{-1}f^\sharp} &
(f')^{-1}\mathcal{O}_{Y'} \ar[u]_{(f')^\sharp}
}
$$
から，上の (\ref{modules-equation-commutative-square-sheaves}) と同様に得られる。
さらにこのような図式
$$
\xymatrix{
X'' \ar[r]_{g'} \ar[d] & X' \ar[d] \\
Y'' \ar[r] & Y'
}
$$
が与えられたとき，$(g')^*c$ と写像
$c' : (g')^*\NL_{X'/Y'} \to \NL_{X''/Y''}$
の合成は写像
$(g \circ g')^*\NL_{X''/Y''} \to \NL_{X/Y}$ である。

\begin{lemma}
\label{modules-lemma-exact-sequence-NL-ringed-topoi}
$f : X \to Y$ と $g : Y \to Z$ を環付き空間の射とする。
$C$ を $\mathcal{O}_X$-加群の層の複体の写像
$\NL_{X/Z} \to \NL_{X/Y}$ の錐とする。標準写像
$$
f^*\NL_{Y/Z} \to C[-1]
$$
があり，これはコホモロジー層の標準六項完全列
$$
\xymatrix{
H^0(f^*\NL_{Y/Z}) \ar[r] &
H^0(\NL_{X/Z}) \ar[r] &
H^0(\NL_{X/Y}) \ar[r] &
0 \\
H^{-1}(f^*\NL_{Y/Z}) \ar[r] &
H^{-1}(\NL_{X/Z}) \ar[r] &
H^{-1}(\NL_{X/Y}) \ar[llu]
}
$$
を与える。
\end{lemma}

\begin{proof}
環の層の写像
$$
(g \circ f)^{-1}\mathcal{O}_Z \to f^{-1}\mathcal{O}_Y \to \mathcal{O}_X
$$
を考え，補題 \ref{modules-lemma-exact-sequence-NL} を適用する。
\end{proof}
